% Preambel mit Einstellungen importieren
% Document type and used packages
\documentclass[open=right, % Kapitel darf nur auf rechten Seite beginnen
    paper=A4,               % DIN-A4-Papier
    a4paper,                % DIN-A4-Papier
    12pt,                   % Schriftgöße
    headings=small,         % Kleine Überschriften
    headsepline=true,       % Trennlinie am Kopf der Seite
    footsepline=false,      % Trennlinie am Fuß der Seite
    bibliography=totoc,     % Literaturverzeichnis in das Inhaltsverzeichnis aufnehmen
    twoside=off,            % Doppelseitiger Druck - auf off stellen für einseitig
    DIV=7,
    chapterprefix=true,     % Kapitel x vor dem Kapitelnamen
    cleardoublepage=plain
    ]{scrbook} 

% Pakete einbinden, die benötigt werden
%\usepackage{scrpage2}				% oboselet!!!!!
\usepackage{scrlayer-scrpage}		% new scrpage2 package
\usepackage[utf8]{inputenc}       % Dateien in UTF-8 benutzen
\usepackage[T1]{fontenc}          % Zeichenkodierung
\usepackage{graphicx}             % Bilder einbinden
\usepackage[ngerman,english]{babel}       % Deutsch und Englisch unterstützen
\usepackage{xcolor}               % Color support
\usepackage{amsmath}              % Matheamtische Formeln
\usepackage{amsfonts}             % Mathematische Zeichensätze
\usepackage{amssymb}              % Mathematische Symbole
\usepackage{float}                % Fließende Objekte (Tabellen, Grafiken etc.)
\usepackage{booktabs}             % Korrekter Tabellensatz
\usepackage[printonlyused]{acronym}   % Abkürzungsverzeichnis [nur verwendete Abkürzugen]
\usepackage{makeidx}              % Sachregister
\usepackage{listings}             % Source Code listings
\usepackage{listingsutf8}         % Listings in UTF8
\usepackage[hang,font={sf,footnotesize},labelfont={footnotesize,bf}]{caption} % Beschriftungen
\usepackage[scaled]{helvet}       % Schrift Helvetia laden
\usepackage[absolute]{textpos}	  % Absolute Textpositionen (für Deckblatt)
\usepackage{calc}                 % Berechnung von Positionen
\usepackage{blindtext}            % Blindtexte
\usepackage[bottom=40mm,left=35mm,right=35mm,top=30mm]{geometry} % Ränder ändern
%\usepackage[square]{natbib}       % Literaturverzeichnis nach DIN mit eckigen Klammern bei \citep
\usepackage{setspace}             % Abstände korrigieren
\usepackage{ifthen}               % Logische Bedingungen mit ifthenelse
\usepackage{scrhack}              % Get rid of tocbasic warnings
\usepackage[pagebackref=false]{hyperref}  % Hyperlinks
\usepackage[all]{hypcap}          % Korrekte Verlinkung von Floats

%von Max eingebunden:
%\usepackage{times}
\usepackage{url}
\usepackage{hyperref}
%\usepackage{underscore}

% Farben definieren
\definecolor{linkblue}{RGB}{0, 0, 100}
\definecolor{linkblack}{RGB}{0, 0, 0}
\definecolor{comment}{RGB}{63, 127, 95}
\definecolor{darkgreen}{RGB}{14, 144, 102}
\definecolor{darkblue}{RGB}{0,0,168}
\definecolor{darkred}{RGB}{128,0,0}
\definecolor{javadoccomment}{RGB}{0,0,240}

% Einstellungen für das Hyperlink-Paket
\hypersetup{
    colorlinks=true,      % Farbige links verwenden       
%    allcolors=linkblue,
    linktoc=all,          % Links im Inhaltsverzeichnis
    linkcolor=linkblack,  % Querverweise
    citecolor=linkblack,  % Literaturangaben
	filecolor=linkblack,  % Dateilinks
	urlcolor	=linkblack    % URLs
}

% Einstellungen für Quelltexte
\lstset{     
      xleftmargin=0.2cm,     
      basicstyle=\footnotesize\ttfamily,
      keywordstyle=\color{darkgreen},
      identifierstyle=\color{darkblue},
      commentstyle=\color{comment}, 
      stringstyle=\color{darkred}, 
      tabsize=2,
      lineskip={2pt},
      columns=flexible,
      inputencoding=utf8,
      captionpos=b,
      breakautoindent=true,
	  breakindent=2em,
	  breaklines=true,
	  prebreak=,
	  postbreak=,
      numbers=none,
      numberstyle=\tiny,
      showspaces=false,      % Keine Leerzeichensymbole
      showtabs=false,        % Keine Tabsymbole
      showstringspaces=false,% Leerzeichen in Strings
      morecomment=[s][\color{javadoccomment}]{/**}{*/},
      literate={Ö}{{\"O}}1 {Ä}{{\"A}}1 {Ü}{{\"U}}1 {ß}{{\ss}}2 {ü}{{\"u}}1 {ä}{{\"a}}1 {ö}{{\"o}}1
}

\urlstyle{same}

% Einstellungen für Überschriften
\renewcommand*{\chapterformat}{%
  \Large\chapapp~\thechapter   % Große Schrift
  \vspace{0.3cm}               % Abstand zum Titel des Kapitels
}

% Abstände für die Überschriften setzen
\renewcommand{\chapterheadstartvskip}{\vspace*{0cm}}
\renewcommand{\chapterheadendvskip}{\vspace*{1.5cm}}%1.5

\RedeclareSectionCommand[
  beforeskip=-1.8\baselineskip,
  afterskip=0.25\baselineskip]{section}

\RedeclareSectionCommand[
  beforeskip=-1.8\baselineskip,
  afterskip=0.15\baselineskip]{subsection}

\RedeclareSectionCommand[
  beforeskip=-1.8\baselineskip,
  afterskip=0.15\baselineskip]{subsubsection}


% In der Kopfzeile nur die kurze Kapitelbezeichnung (ohne Kapitel davor)
\renewcommand*\chaptermarkformat{\thechapter\autodot\enskip}
\automark[chapter]{chapter}

% Einstellungen für Schriftarten
\setkomafont{pagehead}{\normalfont\sffamily}
\setkomafont{pagenumber}{\normalfont\sffamily}
\setkomafont{paragraph}{\sffamily\bfseries\small}
\setkomafont{subsubsection}{\sffamily\itshape\bfseries\small}
\addtokomafont{footnote}{\footnotesize}
\setkomafont{chapter}{\LARGE\selectfont\bfseries}

% Wichtige Abstände
\setlength{\parskip}{0.2cm}  % 2mm Abstand zwischen zwei Absätzen
\setlength{\parindent}{0mm}  % Absätze nicht einziehen
\clubpenalty = 10000         % Keine "Schusterjungen"
\widowpenalty = 10000        % Keine "Hurenkinder"
\displaywidowpenalty = 10000 % Keine "Hurenkinder"
\renewcommand{\footnotesize}{\fontsize{9}{10}\selectfont} % Größe der Fußnoten
\setlength{\footnotesep}{8pt} % Abstand zwischen den Fußnoten

% Index erzeugen
\makeindex

% Einfacher Font-Wechsel über dieses Makro
\newcommand{\changefont}[3]{
\fontfamily{#1} \fontseries{#2} \fontshape{#3} \selectfont}

% Eigenes Makro für Bilder
\newcommand{\bild}[3]{
\begin{figure}[hbt!]
\centering
  \includegraphics[width=#2]{#1}
  \caption{#3}
  \label{#1}
\end{figure}}

% Wo liegt Sourcecode?
\newcommand{\srcloc}{src/}

% Wo sind die Bilder?
\graphicspath{{bilder/}}

% Makros für typographisch korrekte Abkürzungen
\newcommand{\zb}[0]{z.\,B.\ }
\newcommand{\dahe}[0]{d.\,h.\ }
\newcommand{\ua}[0]{u.\,a.\ }

% Flags für Veröffentlichung und Sperrvermerk
\newboolean{hsmapublizieren}

\newboolean{hsmasperrvermerk}


% Dokumenteninfos importieren
% -------------------------------------------------------
% Daten für die Arbeit
% Wenn hier alles korrekt eingetragen wurde, wird das Titelblatt
% automatisch generiert. D.h. die Datei titelblatt.tex muss nicht mehr
% angepasst werden.

\newcommand{\hsmasprache}{de} % de oder en für Deutsch oder Englisch

% Titel der Arbeit auf Deutsch
\newcommand{\hsmatitelde}{Medizinische Simulatoren in virtueller Realität ohne haptisches Feedback, mit Position Based Dynamics und Game Engines}

% Titel der Arbeit auf Englisch
\newcommand{\hsmatitelen}{Medical simulators in Virtual Reality without haptic feedback, using Position Based Dynamics and Game Engines}
%implementation of a buoyancy simulation and an interactive fluid surface for video games

% Weitere Informationen zur Arbeit
\newcommand{\hsmaort}{Heidelberg}    % Ort
\newcommand{\hsmaautorvname}{Maximilian} % Vorname(n)
\newcommand{\hsmaautornname}{Legnar} % Nachname(n)
\newcommand{\hsmaautornmatr}{3544557} % Nachname(n)
\newcommand{\hsmadatum}{30.04.2021} % <----------Datum der Abgabe 01.11.2020 - 30.04.2021
\newcommand{\hsmajahr}{2021} % Jahr der Abgabe
\newcommand{\hsmafirma}{Ruprecht-Karls-Universität Heidelberg} % Firma bei der die Arbeit durchgeführt wurde
\newcommand{\hsmabetreuer}{Prof. Dr. Jürgen Hesser} % Betreuer an der Hochschule
\newcommand{\hsmazweitkorrektor}{Prof. Dr. Susanne Krömker} % Betreuer im Unternehmen oder Zweitkorrektor
\newcommand{\hsmafakultaet}{I} % I für Informatik
\newcommand{\hsmastudiengang}{MEB} % IB IMB UIB IM MTB

% Zustimmung zur Veröffentlichung
\setboolean{hsmapublizieren}{true}   % Einer Veröffentlichung wird zugestimmt
\setboolean{hsmasperrvermerk}{false} % Die Arbeit hat keinen Sperrvermerk

% -------------------------------------------------------
% Abstract

% Kurze (maximal halbseitige) Beschreibung, worum es in der Arbeit geht auf Deutsch
\newcommand{\hsmaabstractde}
{
In dieser Arbeit wird untersucht, welches Potenzial Position Based Dynamics (PBD)-basierte Echtzeitsimulationen \cite{PBD} in Verbindung mit Game Engines für Serious Games im medizinischem Bereich haben. 

Hierfür wird die Unreal Engine 4 zusammen mit NVIDIA FleX \cite{UPP} genutzt, um Software zu untersuchen, die auf einen aktuellen technischen Stand ist. Zuerst wird die Performance und das Simulationsverhalten von NVIDIA FleX  genauer untersucht, um in Erfahrung zu bringen, wie PBD-Simulationen von unterschiedlichen Faktoren beeinflusst werden. Außerdem werden fehlende Funktionalitäten identifiziert, die für die Entwicklung von Serious Games wichtig sind. Daraufhin werden Lösungen präsentiert, mit denen die fehlenden Funktionalitäten implementiert werden können. 
%Dabei wird vor allem auf die umfangreichen Funktionen der Unreal Engine 4 gesetzt. 
So erhalten wir ein umfangreiches Softwareframework, mit dem sich schnell und einfach medizinische Echtzeitsimulationen erstellen lassen.

Außerdem wird eine Lösung präsentiert, mit der sich einfache chirurgische Simulatoren erstellen lassen, ohne Verwendung von haptischer Hardware. Hierfür wird nur eine handelsübliche VR-Anlage (z. B. Oculus Quest 2) mit Motion Controllern benötigt. Dabei kann der Nutzer mithilfe seiner Motion Controller nach virtuellen Werkzeugen greifen und diese intuitiv benutzen. Mithilfe von simulierten Kräften übt ein virtuelles Werkzeug Interaktionskräfte auf die virtuelle Umgebung aus.

Es wird auch eine optisch realistische Nadel-Gewebe-Interaktion umgesetzt. Gewebe und Nadel werden dabei mithilfe von Partikeln und \textit{Shape Matching Constraints} \cite{Shape} modelliert. Die Partikelstruktur des Gewebemodells wird dabei nicht verändert, das heißt, das virtuelle Gewebe wird nicht von der Nadel verletzt beziehungsweise geschnitten.

%Basierend auf Reynolds Boid-System \cite{Reynolds87flocks} wird ein Modell für die Simulation von Tierschwärmen entwickelt. Wesentlicher Fokus liegt dabei auf lebendig wirkende Schwarmbewegungen, interaktives Verhalten und Performance. 
%
%Statt den simulierten Schwarm wie ein Partikelsystem zu behandeln, wird Reynolds Boid-Modell in ein Zellensystem, das mit den zellulären Automaten verwandt ist, transferiert. Diese Anpassung wird vorgenommen, damit die Simulation von herkömmlichen Shadern ausgeführt werden kann.
%
%Das Simulationsmodell unterteilt den Schwarm nicht in einzelne separate Individuen. Stattdessen ist der Schwarm eine im Raum ungleichmäßig verteilte Menge.
%Der Schwarm wird daher als leuchtende Punkte-Wolke, von einer Textur dargestellt.
%
%Die konzipierte Schwarmsimulation wird mittels Content-Driven Multipass Rendering prototypisch implementiert und in Echtzeit ausgeführt.




%Im erstem Teil der Arbeit wird eine Auftriebssimulation für Videospiele entwickelt, um 3D-Modelle auf Wasserflächen schwimmen lassen zu können. Das zugrundeliegende mathematische Modell diese Simulation basiert auf dem archimedischem Prinzip. Die Auftriebssimulation wird mit der Unreal Engine 4 implementiert und ihre integrierte Physik Engine, nvidia PhysX3.3, wird genutzt um Auftriebskräfte auf virtuelle Schwimmkörper wirken zu lassen, welche einem 3D-Modell zugeordnet werden.
%
%Im zweitem Teil der Arbeit wird eine interaktive Wasseroberfläche für Videospiele entwickelt. Lösungsansatz ist hierbei das sogenannte \glqq IWave-Verfahren\grqq{} \cite{tessendorf}. Dieses Verfahren wird mit der Unreal Engine 4 und einer Technik namens \glqq Content Driven Multipass Rendering\grqq{} umgesetzt.

%\cite{tessendorf}
}

% Kurze (maximal halbseitige) Beschreibung, worum es in der Arbeit geht auf Englisch

\newcommand{\hsmaabstracten}
{
This thesis examines the potential of Position Based Dynamics (PBD)-based real-time simulations \cite{PBD} in connection with game engines for serious games with medical background.

For this purpose, the Unreal Engine 4 is used together with NVIDIA FleX \newline \cite{UPP} to examine software that is state-of-the-art. First, the performance and simulation behavior of FleX is examined in more detail in order to find out how FleX simulations are influenced by different factors. Then missing functionalities are identified, which could be useful for serious game development. After that, we present solutions which can be used to implement these missing functionalities.
In that way we optain a comprehensive and easy-to-use framework which serves as a development platform for medical real-time simulations.

In addition, a solution is presented how to implement immersive surgical simulators without haptic feedback. Such surgical simulations can be used with standard VR systems, like an Oculus Quest 2, HTC vive or valve index. The user can use his motion controllers to grab virtual tools and use them. With the help of simulated forces, a virtual tool can apply accurate interaction forces to virtual environment.

An optically realistic needle-tissue interaction is also implemented. Tissue and needle are modeled with the help of particles and \textit{Shape Matching Constraints} \newline \cite{Shape}. The particle structure of the tissue model is not changed.
}





\begin{document}
\frontmatter

% Römische Ziffern ür die "Front-Matter"
\setcounter{page}{0}
\changefont{ptm}{m}{n}  % Times New Roman für den Fließtext
\renewcommand{\rmdefault}{ptm}

% Titelblatt
% -------------------------------------------------------
% In dieser Datei sollten eigentlich keine Veränderungen mehr
% notwendig sein.
% -------------------------------------------------------

\thispagestyle{empty}

% Fakultäten der HS-Mannheim
% -------------------------------------------------------
\ifthenelse{\equal{\hsmafakultaet}{I}}%
  {\newcommand{\hsmafakultaetlangde}{Fakultät für Mathematik und Informatik}%
   \newcommand{\hsmafakultaetlangen}{Interdisciplinary Center for Scientific Computing}}{}

\ifthenelse{\equal{\hsmastudiengang}{MEB}}%%%%--------------------------------------------------------------------
  {\newcommand{\hsmastudienganglangde}{Fakultät für Mathematik und Informatik}%
   \newcommand{\hsmastudienganglangen}{Applied Computer Science}%Mechatronic
   \newcommand{\hsmatypde}{Master-Thesis}%
   \newcommand{\hsmatypen}{Master-Thesis}%
   \newcommand{\hsmagrad}{\hsmabsc}}{}

\newcommand{\hsmabsc}{Bachelor of Science (B.Sc.)}
\newcommand{\hsmaba}{Bachelor of Arts (B.A.)}
\newcommand{\hsmamaster}{Master of Science (M.Sc.)}
\newcommand{\hsmamastera}{Master of Arts (M.A.)}
\newcommand{\hsmamasterba}{Master of Business Administration (MBA)}

\newcommand{\hsmakoerperschaftde}{Ruprecht-Karls-Universität Heidelberg}
\newcommand{\hsmakoerperschaften}{Ruprecht Karl University of Heidelberg}

\newcommand{\hsmaautorbib}{\hsmaautornname, \hsmaautorvname} % Autor Nachname, Vorname
\newcommand{\hsmaautor}{\hsmaautorvname \ \hsmaautornname} % Autor Vorname Nachname

\ifthenelse{\equal{\hsmasprache}{de}}%
  {\newcommand{\hsmatyp}{\hsmatypde}%
   \newcommand{\hsmathesistype}{ \hsmagrad}%
   \newcommand{\hsmakoerperschaft}{\hsmakoerperschaftde}%
   \newcommand{\hsmastudiengangname}{\hsmastudienganglangde}%
   \newcommand{\hsmaInstitut}{Institut für Wissenschaftliches Rechnen (IWR)}%
   \newcommand{\hsmastudienganglang}{\hsmastudienganglangde}%
   \newcommand{\hsmatitel}{\hsmatitelde}%
   \newcommand{\hsmatutor}{Betreuer}%
   \newcommand{\hsmafakultaetlang}{\hsmafakultaetlangde}%
   \selectlanguage{ngerman}}%
  {\newcommand{\hsmatyp}{\hsmatypen}%
   \newcommand{\hsmathesistype}{for the acquisition of the academic degree \hsmagrad}%
   \newcommand{\hsmakoerperschaft}{\hsmakoerperschaften}%
   \newcommand{\hsmastudiengangname}{Course of Studies: \hsmastudienganglang}%
   \newcommand{\hsmastudienganglang}{\hsmastudienganglangen}%
   \newcommand{\hsmatitel}{\hsmatitelen}%
   \newcommand{\hsmatutor}{Tutors}
   \newcommand{\hsmafakultaetlang}{\hsmafakultaetlangen}%
   \selectlanguage{english}}%


% Daten in die Standard-Felder von KOMA-Script eintragen
\titlehead{\hsmatyp\ in\  \hsmastudienganglang}
\subject{}
\title{\hsmatitel}
\author{\hsmaauthor}
\date{\small{\hsmadatum}}

% Daten für das fertige PDF-Dokument
\hypersetup{
  pdftitle={\hsmatitel},  % Titel des Dokuments
  pdfauthor={\hsmaautor},              % Autor
  pdfsubject={\hsmatyp\ in\ \hsmastudienganglang},                % Thema
  pdfkeywords={\hsmatitel}         % Schlüsselworte
}

\newlength{\bindekorrektur}
\newlength{\seitenanfang}
\newlength{\seitenbreite}
  
\setlength{\bindekorrektur}{-46mm}   % Korrektur der horizontalen Position
\setlength{\seitenanfang}{0mm}       % Korrektur der vertikalen Position
\setlength{\seitenbreite}{297mm}

%------------------ LOGO habe ich raus commentiert: ----------------------------------------
%\noindent\includegraphics[width=4cm]{unilogo1.jpg}\\%hsma-logo.pdf     ue-logo.png

% Uni
\begin{textblock*}{\seitenbreite}(\bindekorrektur,\seitenanfang + 40mm)
  \centering\large \sffamily
  \textbf{\hsmafirma}
\end{textblock*}

% Fakultät
\begin{textblock*}{\seitenbreite}(\bindekorrektur,\seitenanfang + 50mm)
  \centering\large\sffamily
  \textbf{\hsmafakultaetlang} %\\
  %\vspace{2mm}
  %\hsmakoerperschaft
\end{textblock*}

% Institut und typ
\begin{textblock*}{\seitenbreite}(\bindekorrektur,\seitenanfang + 60mm)
  \centering\large\sffamily
  %\hsmatyp\\
  %\begin{small}\hsmathesistype \end{small}\\
  %\vspace{2mm}
  %hier war vorher studiengang
  \textbf{\hsmaInstitut} \\
  \vspace{25mm}
  \textbf{\hsmatyp}
\end{textblock*}

% Titel der Arbeit
\begin{textblock*}{128mm}(38mm,\seitenanfang + 120mm) % 4,5cm vom linken Rand und 6,0cm vom oberen Rand
  \centering\Large\sffamily
%  \vspace{4mm} % Kleiner zusätzlicher Abstand oben für bessere Optik
  \hsmatitel \\
  \vspace{4mm}
  \hsmatitelen
\end{textblock*}%

% Name und so
\begin{textblock*}{128mm}(38mm,\seitenanfang + 210mm)%103mm
   \large \sffamily
  Name: \hspace{32mm}  \hsmaautor  
  
  Matrikelnummer: \hspace{10mm} \hsmaautornmatr
  
  Betreuer: \hspace{26mm} \hsmabetreuer
  
  Zweitkorrektorin: \hspace{10,5mm} \hsmazweitkorrektor
  
  Datum der Abgabe: \hspace{4,5mm} \hsmadatum
\end{textblock*}

%% Institut
%\begin{textblock*}{\seitenbreite}(\bindekorrektur,\seitenanfang + 50mm)
%  \centering\large\sffamily
%  %\hsmatyp\\
%  %\begin{small}\hsmathesistype \end{small}\\
%  %\vspace{2mm}
%  %hier war vorher studiengang
%  \textbf{\hsmaInstitut}
%\end{textblock*}



% Datum
%\begin{textblock*}{\seitenbreite}(\bindekorrektur,\seitenanfang + 190mm)
%  \centering\large 
%  \textsf{\hsmadatum}
%\end{textblock*}



% Betreuer
%\begin{textblock*}{\seitenbreite}(\bindekorrektur,\seitenanfang + 240mm)
%  \centering\large\sffamily
%  \hsmatutor \\
%  \vspace{2mm}
%  \hsmabetreuer\\
%  \vspace{2mm}
%  \hsmazweitkorrektor
%\end{textblock*}





% %%%%%%%%%%%%Bibliographische Informationen
\null\newpage
\thispagestyle{empty}
  
\newcommand{\hsmabibde}{\begin{small}\textbf{\hsmaautorbib}: \\ \hsmatitelde \ / \hsmaautor. \ -- \\ \hsmatypde, \hsmaort : \hsmakoerperschaftde, \hsmajahr. \pageref{lastpage} Seiten.\end{small}}

\newcommand{\hsmabiben}{\begin{small}\textbf{\hsmaautorbib}: \\ \hsmatitelen \ / \hsmaautor. \ -- \\ \hsmatypen, \hsmaort : \hsmakoerperschaften, \hsmajahr. \pageref{lastpage} pages. \end{small}}

\ifthenelse{\equal{\hsmasprache}{de}}%
  {\hsmabibde \\ \vspace{0.5cm} \\ \hsmabiben}
  {\hsmabiben \\ \vspace{0.5cm} \\ \hsmabibde}
%\hsmabibde

% Erklärung
\clearpage\setcounter{page}{1}
\thispagestyle{empty}
\textsf{\large\textbf{Erklärung}}

Ich versichere hiermit, dass ich die vorliegende Arbeit selbstständig verfasst und keine anderen als die angegebenen Hilfsmittel benutzt habe. Sowohl inhaltlich als auch wörtlich entnommene Inhalte wurden als solche kenntlich gemacht. Die Arbeit ist in gleicher
oder vergleichbarer Form noch bei keiner anderen Prüfungsbehörde eingereicht worden.

\ifthenelse{\boolean{hsmapublizieren} \and \not\boolean{hsmasperrvermerk}}%
{
\vspace{0.5cm}
Ich bin damit einverstanden, dass meine Arbeit veröffentlicht wird, d.\,h. dass die Arbeit elektronisch gespeichert, in andere Formate konvertiert, auf den Servern der Universität Heidelberg öffentlich zugänglich gemacht und über das Internet verbreitet werden darf. 
}{}%


\vspace{1cm}
\hsmaort, \hsmadatum \\

\vspace{1.2cm}						                                      
\hsmaautor

\ifthenelse{\boolean{hsmasperrvermerk}}%
{%
\vspace{11cm}
\color{red}\textsf{\large\textbf{Sperrvermerk}}

Diese Arbeit basiert auf internen und vertraulichen Daten des Unternehmens \hsmafirma.

Diese Arbeit darf Dritten, mit Ausnahme der betreuenden Dozenten und befugten Mitglieder des Prüfungsausschusses, ohne ausdrückliche Zustimmung des Unternehmens und des Verfassers nicht zugänglich gemacht werden.

Eine Vervielfältigung und Veröffentlichung der Arbeit ohne ausdrückliche Genehmigung -- auch in Auszügen -- ist nicht erlaubt.
\color{black}
}{}

\cleardoublepage

% Abstract
\chapter*{Abstract}

\ifthenelse{\equal{\hsmasprache}{de}}%
  {\subsubsection*{\hsmatitelde}\hsmaabstractde\subsubsection*{\hsmatitelen}\hsmaabstracten}
  {\subsubsection*{\hsmatitelen}\hsmaabstracten\subsubsection*{\hsmatitelde}\hsmaabstractde}
% Inhaltsverzeichnis erzeugen
\cleardoublepage
\pdfbookmark{\contentsname}{Contents}
\tableofcontents

% Korrigiert Nummerierung bei mehrseitigem Inhaltsverzeichnis
\cleardoublepage
\newcounter{frontmatterpage}
\setcounter{frontmatterpage}{\value{page}}

% Neue Seite
\cleardoublepage

% Backmatter mit normalem Zeilenabstand setzen
\singlespacing

% Römische Ziffern für die "Back-Matter", fortlaufend mit "Front-Matter"
\pagenumbering{roman}
\setcounter{page}{\value{frontmatterpage}}

% Abkürzungsverzeichnis
\chapter*{Abkürzungsverzeichnis}
\addcontentsline{toc}{chapter}{Abkürzungsverzeichnis}

\begin{acronym}
%Beispiele:
\acro{ABK}{Abkürzung}
\acro{ACM}{Association of Computing Machinery}
\acro{PDF}{Portable Document Format}
\acro{IEEE}{Institute of Electrical and Electronics Engineers}
\acro{ISO}{International Organization for Standardization}

%Mein Stuff:
\acro{UE4}{Unreal Engine 4}
\acro{GPU}{Graphics Processing Unit}
\acro{CPU}{Central Processing Unit}
\acro{HLSL}{High Level Shader Language}
\acro{GLSL}{OpenGL Shading Language}
\acro{PBD}{Position Based Dynamics}
\acro{VR}{Virtual Reality}
\acro{FPS}{Frames Per Second}
\acro{SDF}{Signed Distance Field}


%\acro{SIMD}{Single instruction, multiple data}
\end{acronym}


% Tabellenverzeichnis erzeugen
\cleardoublepage
\phantomsection
\addcontentsline{toc}{chapter}{Tabellenverzeichnis}
\listoftables

% Abbildungsverzeichnis erzeugen
\cleardoublepage
\phantomsection
\addcontentsline{toc}{chapter}{Abbildungsverzeichnis}
\listoffigures

% Listingverzeichnis erzeugen
%\cleardoublepage
%\phantomsection
%\addcontentsline{toc}{chapter}{Quellcodeverzeichnis}
%\lstlistoflistings








% Arabische Zahlen für den Hauptteil
\mainmatter

% Den Hauptteil mit vergrößertem Zeilenabstand setzen
\onehalfspacing

% ------------------------------------------------------------------
% Hauptteil der Arbeit
%% Die Arbeit besteht aus Kapiteln (chapter)
\chapter{Gliederungsebenen}

% Jedes Kapitel besteht aus Unterkapiteln (section)
\section{Zweite Ebene}

% Unterkapitel können noch einmal durch subsections untergliedert 
% werden (jetzt auf der 3. Ebene)
\subsection{Dritte Ebene}

% Mit Labels können Sie sich später im Text wieder auf diese Stelle beziehen
\label{Gliederung:EbeneDrei}

% Einträge für den Index anlegen. Ein Index wird normalerweise in einer Abschluss
% Arbeit nicht benötigt.
\index{Gliederung!Ebenen}

% Auf der 4. Ebene liegen die subsubsections. In diesem Template bekommt die
% 4. Ebene keinen Nummern mehr und erscheint auch nicht im Inhaltsverzeichnis
\subsubsection{Vierte Ebene}

% Auf der 5. Ebene werden einzelne Absätze mit Überschriften versehen.
\paragraph{Fünfte Ebene} Anders als in diesem Beispiel, darf in Ihrer Arbeit kein Gliederungspunkt auf seiner Ebene alleine stehen. D.\,h. wenn es ein 1.1 gibt, muss es auch ein 1.2 geben.\citep{Kornmeier2011}
  % Externe Datei einbinden
%\chapter{Schreibstil und Typographie}


\section{Hervorhebungen}
\label{Einleitung:Textauszeichnungen}

Achten Sie bitte auf die grundlegenden Regeln der Typographie\index{Typographie}\footnote{Ein Ratgeber in allen Detailfragen ist \cite{Forssman2002}.}, wenn Sie Ihren Text schreiben. Hierzu gehören z.\,B. die Verwendung der richtigen "`Anführungszeichen"' und der Unterschied zwischen Binde- (-), Gedankenstrich (--) und langem Strich (---).

Wenn Sie Text hervorheben wollen, dann setzten Sie ihn \textit{kursiv} (Italic) und nicht \textbf{fett} (Bold). Fettdruck ist Überschriften vorbehalten; im Fließtext stört er den Lesefluss. Das \underline{Unterstreichen} von Fließtext ist im gesamten Dokument tabu und kann maximal bei Pseudo-Code vorkommen.\index{Hervorhebungen}


\section{Anführungszeichen}

Deutsche Anführungszeichen gehen so: "`dieser Text steht in \glq Anführungszeichen\grq; alles klar?"'. Englische Anführungszeichen werden anders benutzt: ``this is an `English' quotation.''


\section{Abkürzungen}
\index{Abkürzungen}
\index{Abbreviation|see{Abkürzungen}}

Eine \ac{ABK} wird bei der ersten Verwendung ausgeschrieben. Danach nicht mehr: \ac{ABK}. Man kann allerdings die Langform explizit anfordern: \acl{ABK} oder die Kurzform \acs{ABK} oder auch noch einmal die Definition: \acf{ABK}.

Beachten Sie, dass bei Abkürzungen, die für zwei Wörter stehen, ein kleines Leerzeichen nach dem Punkt kommt: z.\,B. bzw. \zb, d.\,h. bzw. \dahe.


\section{Querverweise}

Querverweise auf eine Kapitelnummer macht man im Text mit \verb+\ref+ (Kapitel~\ref{Einleitung:Textauszeichnungen}) und auf eine bestimmte Seite mit \verb+\pageref+ (Seite~\pageref{Einleitung:Textauszeichnungen}).


\section{Fußnoten}

Fußnoten werden einfach mit in den Text geschrieben und zwar genau an die Stelle\footnote{An der die Fußnote auftauchen soll.}



\section{Fremdsprachige Begriffe}

Wenn Sie Ihre Arbeit auf Deutsch verfassen, gehen Sie sparsam mit englischen Ausdrücken um. Natürlich brauchen Sie etablierte englische Fachbegriffe, wie z.\,B. \textit{Interrupt}, nicht zu übersetzen. Sie sollten aber immer dann, wenn es einen gleichwertigen deutschen Begriff gibt, diesem den Vorrang geben. Den englischen Begriff (\textit{term}) können Sie dann in Klammern oder in einer Fußnote\footnote{Englisch: \textit{footnote}.} erwähnen. Absolut unakzeptabel sind deutsch gebeugte englische Wörter oder Kompositionen aus deutschen und englischen Wörtern wie z.\,B. downgeloadet, upgedated, Keydruck oder Beautyzentrum. 



\section{Tabellen}

Tabellen werden normalerweise ohne vertikale Striche gesetzt, sondern die Spalten werden durch einen entsprechenden Abstand voneinander getrennt.\footnote{Siehe \cite[S. 89]{Willberg1999}.} Zum Einsatz kommen ausschließlich horizontale Linien (siehe Tabelle~\ref{Kap2:Kopplungsformen}).

\begin{table}[h]
  \caption{Ebenen der Kopplung und Beispiele für enge und lose Kopplung}
  \label{Kap2:Kopplungsformen}
  \renewcommand{\arraystretch}{1.2}
  \centering
  \sffamily
  \begin{footnotesize}
    \begin{tabular}{l l l}
    \toprule
    \textbf{Form der Kopplung} & \textbf{enge Kopplung} & \textbf{lose Kopplung}\\
    \midrule
    Physikalische Verbindung	&	Punkt-zu-Punkt	& 	über Vermittler\\
    Kommunikationsstil	&	synchron		&	asynchron\\
    Datenmodell	&	komplexe gemeinsame Typen	&	nur einfache gemeinsame Typen\\
    Bindung	&	statisch		&	dynamisch\\
    \bottomrule
    \end{tabular}
  \end{footnotesize}
  \rmfamily
\end{table}

Eine Tabelle fließt genauso, wie auch Bilder durch den Text. Siehe Tabelle~\ref{Kap2:Kopplungsformen}.


\section{Aufzählungen}

Aufzählungen sind toll.

\begin{itemize}
  \item Ein wichtiger Punkt
  \item Noch ein wichtiger Punkt
  \item Ein Punkt mit Unterpunkten
    \begin{itemize}
      \item Unterpunkt 1
      \item Unterpunkt 2      
    \end{itemize}
  \item Ein abschließender Punkt ohne Unterpunkte
\end{itemize}


Aufzählungen mit laufenden Nummern sind auch toll.

\begin{enumerate}
  \item Ein wichtiger Punkt
  \item Noch ein wichtiger Punkt
  \item Ein Punkt mit Unterpunkten
    \begin{enumerate}
      \item Unterpunkt 1
      \item Unterpunkt 2      
    \end{enumerate}
  \item Ein abschließender Punkt ohne Unterpunkte
\end{enumerate}


\section{Zitate}

\subsection{Zitate im Text}

Wichtig ist das korrekte Zitieren von Quellen, wie es auch von \cite{Kornmeier2011} dargelegt wird. Interessant ist in diesem Zusammenhang auch der Artikel von \cite{Kramer2009}. Häufig werden die Zitate auch in Klammern gesetzt, wie bei \citep{Kornmeier2011} und mit Seitenzahlen versehen \citep[S. 22--24]{Kornmeier2011}.


\subsection{Zitierstile}

Verwenden Sie eine einheitliche und im gesamten Dokument konsequent durchgehaltene Zitierweise\index{Zitierweise}. Es gibt eine ganze Reihe von unterschiedlichen Standards für das Zitieren und den Aufbau eines Literaturverzeichnisses. Sie können entweder mit Fußnoten oder Kurzbelegen im Text arbeiten. Welches Verfahren Sie einsetzen ist Ihnen überlassen, nur müssen Sie es konsequent durchhalten.

In der Informatik ist das Zitieren mit Kurzbelegen\index{Zitat!Kurzbeleg} im Text (Harvard-Zitierweise) weit verbreitet, wobei für das Literaturverzeichnis häufig die Regeln der \acs{ACM} oder \acs{IEEE} angewandt werden.\footnote{Einen Überblick über viele verschiedene Zitierweisen finden Sie in der \url{http://amath.colorado.edu/documentation/LaTeX/reference/faq/bibstyles.pdf}}

Denken Sie daran, dass das Übernehmen einer fremden Textstelle ohne entsprechenden Hinweis auf die Herkunft in wissenschaftlichen Arbeiten nicht akzeptabel ist und dazu führen kann, dass die Arbeit nicht anerkannt wird. Plagiate\index{Plagiat!Bewertung} werden mit mangelhaft (5,0) bewertet und können weitere rechtliche Schritte nach sich ziehen.


\subsection{Zitieren von Internetquellen}

Internetquellen\index{Zitat!Internetquellen} sind normalerweise \textit{nicht} zitierfähig. Zum einen, weil sie nicht dauerhaft zur Verfügung stehen und damit für den Leser möglicherweise nicht beschaffbar sind und zum anderen, weil häufig der wissenschaftliche Anspruch fehlt.\footnote{Eine lesenswerte Abhandlung zu diesem Thema findet sich (im Internet) bei \cite{Weber2006}}

Wenn ausnahmsweise doch eine Internetquelle zitiert werden muss, z.\,B. weil für eine Arbeit dort Informationen zu einem beschriebenen Unternehmen abgerufen wurden, sind folgende Punkte zu beachten:

\begin{itemize}
\item Die Webseite ist auszudrucken und im Anhang der Arbeit beizufügen,
\item das Datum des Abrufs und die URL sind anzugeben,
\item verwenden Sie Internet-Seiten ausschließlich zu illustrativen Zwecken (z.\,B. um einen Sachverhalt noch etwas genauer zu erläutern), aber nicht zur Faktenvermittlung (z.\,B. um eine Ihrer Thesen zu belegen).
\end{itemize}

Wenn Sie aufgrund der Natur Ihrer Arbeit sehr viele Internetquellen benötigen, dann können Sie diese statt sie auszudrucken auch in elektronischer Form abgeben (CD/DVD). Als Abgabeformat der elektronischen Quellen ist PDF/A\footnote{Bei PDF/A handelt es sich um ein besonders stabile Variante des \ac{PDF}, die von der  \ac{ISO} standardisiert wurde.} vorteilhaft, weil es von allen Formaten die größte Stabilität besitzt.
Auf der CD/DVD geben Sie bitte auch eine HTML-Version des Literaturverzeichnisses ab, in der die Online-Quellen sowie die gespeicherten PDF-Dateien verlinkt sind.

Wikipedia\index{Zitat!Wikipedia} stellt einen immensen Wissensfundus dar und enthält zu vielen Themen hervorragende Artikel. Sie müssen sich aber darüber im Klaren sein, dass die Artikel in Wikipedia einem ständigen Wandel unterworfen sind und nicht als Quelle für wissenschaftliche Fakten genutzt werden sollten. Es gelten die allgemeinen Regeln für das Zitieren von Internetquellen. Sollten Sie doch Wikipedia nutzen müssen, verwenden Sie bitte ausschließlich den Perma-Link\footnote{Sie erhalten den Permalink über die Historie der Seite und einen Klick auf das Datum.}\index{Permalink} zu der Version der Seite, die Sie aufgerufen haben.


 % Externe Datei einbinden
%\chapter{Einbinden von Grafiken und Sourcecode}
\label{Kap2}

\section{Bilder}

Natürlich können auch Grafiken und Bilder eingebunden werden, siehe z.\,B. Abbildung~\ref{Kap2:NasaRover}.

\begin{figure}[h]
  \centering
  \includegraphics[width=6cm]{kapitel2/nasa_rover}
  \caption{Ein Nasa Rover}
  \label{Kap2:NasaRover}
\end{figure}

Man kann sich auch selber ein Makro für das Einfügen von Bildern schreiben:

\bild{kapitel2/modell_point_to_point}{6cm}{Point to Point}


\clearpage % Alle Bilder, die bisher kamen ausgeben



\section{Formelsatz}

Eine Formel gefällig? Mitten im Text $a_2 = \sqrt{x^3}$ oder als eigener Absatz (siehe Formel~\ref{Formel}):

\begin{equation}
\begin{bmatrix}         
   1 &  4 &  2 \\
   4 &  0 & -3 
\end{bmatrix}        
        \cdot
\begin{bmatrix}                
   1 &  1 &  0 \\
  -2 &  3 &  5 \\
   0 &  1 &  4 
\end{bmatrix}        
       {=} 
\begin{bmatrix}               
  -7 &  15 &  28 \\
   4 &   1 & -12 
\end{bmatrix}
\label{Formel}
\end{equation}


\section{Sourcecode}

Man kann mit Latex auch ganz toll Sourcecode in den Text aufnehmen.

\subsection{Aus einer Datei}

\lstinputlisting[firstline=2,language=Java,caption={Crypter-Interface},label=lst:CrypterInterface]{\srcloc/Crypter.java}


\subsection{Inline}

\begin{lstlisting}[language=Java,caption=Methode checkKey()]
    /**
     * Testet den Schlüssel auf Korrektheit: Er muss mindestens die Länge 1
     * haben und darf nur Zeichen von A-Z enthalten.
     *
     * @param key zu testender Schlüssel
     * @throws CrypterException wenn der Schlüssel nicht OK ist.
     */
    protected void checkKey(Key key) throws CrypterException {

        // Passt die Länge?
        if (key.getKey().length == 0) {
            throw new CrypterException("Der Schlüssel muss mindestens " +
                    "ein Zeichen lang sein");
        }

        checkCharacters(key.getKey(), ALPHABET);
    }
\end{lstlisting}

  % Externe Datei einbinden


% Die Arbeit besteht aus Kapiteln (chapter)
% strg + t -> commend selection. strg+ u -> uncommend selection
\chapter{Einleitung}
\label{chap_Einleitung}



\section{Motivation und Problemstellung}


% Motivation für Serous Games:
Die klassische Art, bestimmte Operationen praktisch zu unterrichten, ist die Beobachtung eines erfahrenen Chirurgen über mehrere durchgeführte Operationen. Dieser Ansatz wird zunehmend durch rechtliche und ethische Bedenken hinsichtlich der Patientensicherheit, des Zeitbedarfs und der Kosten für die OP-Zeit behindert \cite{SurgSim}.

Deswegen sind Serious Games in der Medizin wichtige Hilfsmittel um Personal auszubilden und zu trainieren. Vor allem in der Chirurgie werden immer häufiger medizinische Echtzeitsimulationen eingesetzt um Personal auf Operationen vorzubereiten \cite{SimRole}. So kann das Personal besser auf eine Operation am echten Patienten vorbereitet werden und die Chance, dass die Operation fehlerfrei verläuft, kann gesteigert werden \cite{VRNeuro}.

%Motivation PBD:
Die Simulationstechnik \ac{PBD} \cite{PBD} hat sich bereits in vielen Arbeiten als vielversprechende Simulationstechnik für chirurgische Simulatoren erwiesen (siehe Kapitel \ref{section_PBD} für einige Beispiele). 

%Motivation PBD und Überleitung zu Problemstellung 1:
Aufgrund der Vielseitigkeit, Stabilität, Einfachheit und Effizienz von PBD wird auch in dieser Arbeit eine PBD-basierte Physik Engine (NVIDIA FleX) genutzt um Gewebe, Flüssigkeiten und Rigid Bodies in Echtzeit zu simulieren. Zusätzlich wird eine Game Engine (Unreal Engine 4 \cite{UE4FlexDoc}) eingesetzt, die fast alle restlichen Elemente implementiert, die für Serious Games im chirurgischem Bereich benötigt werden. Dabei werden einige zusätzliche Softwareelemente methodisch entwickelt und implementiert um die vorhandene Game- und Physik-Engine um zusätzliche Funktionalitäten zu erweitern, die für chirurgische Simulatoren nützlich sein könnten (Kapitel \ref{chap_Flex_Engine}).

% Motivation VR Serous Games:
\ac{VR} spielt für Serous Games auch eine wichtige Rolle. 
%Mit einem handelsüblichen \ac{HMD} kann die virtuelle Umgebung  besonders immersiv dargestellt werden.
Mithilfe von handelsüblichen VR-Geräten und haptischem Feedback können Echtzeitsimulatoren erstellt werden, die besonders immersiv und realistisch wirken \cite{VRSim20}. 
Studien zeigen, dass VR-Anwendungen mit haptischem Feedback gut geeignet sind, um medizinische Fachkräfte effektiv für bestimmte Operationsverfahren auszubilden  \cite{VRHapticSim}.

% Überleitung problemstellung 2:
In dieser Arbeit wird hingegen eine Methode präsentiert, mit der chirurgische VR-Simulatoren ohne haptisches Feedback erstellt werden können. Der Nutzer benötigt nur eine handelsübliche VR-Anlage 
%mit Handtracking oder Motion Controllern 
um unterschiedliche Operationen trainieren zu können. 
%Für diese Arbeit wurde eine Oculus Quest 2 für 300 Euro eingesetzt. 
Dies ermöglicht es Nutzern, die keinen Zugriff auf haptische Hardware haben, Operationen in virtueller Realität zu trainieren. 

Eine besondere Herausforderung ist dabei eine Methode zu finden, mit der der Nutzer möglichst realistisch und präzise mit der virtuellen Umgebung interagieren kann, ohne haptisches Feedback zu erhalten. In dieser Arbeit wurde hierfür ein System entwickelt, das es dem Spieler ermöglicht mit virtuellen Werkzeugen simulierte Kräfte auf die virtuelle Umgebung auszuüben. So kann der Spieler beispielsweise eine Nadel in simuliertes Gewebe einführen, indem er seine echte Hand auf bestimmte Weise positioniert und beschleunigt und dadurch Interaktionskräfte zwischen Nadel und Gewebe erzeugt. 
Dabei erhält der Nutzer nur visuelles Kräftefeedback und kein haptisches.
Die Nadel-Gewebe Interaktion mit rein virtuellen Kräften ist dabei auch eine besondere Herausforderung, die in Kapitel \ref{sec_needle} behandelt wird. 
Bei der Nadel-Gewebe-Interaktion wird das Partikelmodel des Gewebes nicht verändert. 



%Um am Ende dieser Arbeit ein nützliches Framework zu erhalten, das für die Entwicklung von chirurgischen Simulatoren geeignet ist, wird analysiert welche weiteren Funktionalitäten für die Entwicklung von chirurgischen Simulatoren nützlich sind. Diese Funktionalitäten werden dann methodisch 
%um photorealistische Bilder zu zeichnen und von ihren vielseitigen Funktionalitäten zu profitieren, welche gut für die Entwicklung von Serious Games im chirurgischem Bereich geeignet sind.

%Außerdem wird in dieser Arbeit eine Game Engine, die Unreal Engine 4, genutzt um von ihrer effizienten Kollisionserkennung, ihren Renderfähigkeiten und vielseitigen Software-Bibliotheken zu profitieren. 

%In dieser Arbeit werden Konzepte vorgestellt, mit denen PBD basierte Physik Engines zusammen mit Game Engines 
%In dieser Arbeit wird mithilfe einer Game Engine und einer PBD basierten Physik Engine 



%\section{Einordnung und Abgrenzung}

%Sehr ähnlich zu \cite{VRSim20}. Unterschied: In \cite{VRSim20} wirkt sich Reaktionskraft auf echtes Werkzeug (haptische Hardware) aus. Hier wird Reaktionskraft auf virtuelles Werkzeug angewandt und Spieler erhält nur visuelles Force Feedback. 

%Ist dann zwar Weniger realistisch (das ist natürlich unbestreitbar), dafür kann Software mit jedem handelsüblichen VR-Gerät genutzt werden. Zusätzliche haptische Hardware wird nicht benötigt. Gute VR-Geräte gibts schon für 300 euro, Jede Privatperson die ein VR-Gerät mit PC oder gar PS-VR hat, könnte Trainings-Software nutzen.

%Habe Für Gewebe SelfCollision verboten, im Gegensatz zu vielen anderen Arbeiten (\cite{BreastBiopsy} \cite{PBDKidney}).

%Einfache Nadel-Gewebe-Interaktion ohne Verletzung oder schneiden, mit virtuellen Kräften auf Nadel (Reibungs und Klemmkräfte, Interaktionskraft usw)


%Falls bei präsi gefragt: Wie wird kann force feedback berrechnet werden? Siehe Fig. 7 aus \cite{VRSim20} virtuelles Lyaryngoskop (proxy) drückt Partikel ein (Partikel sehen Proxy-Werkzeug als statische Wand die sich nicht verschieben lässt). Vektoren zu Rest-positionen werden gespannt, Durchschnitt, multiplizieren mit stiffnes-faktor für das Gewebe. Genau so hab ichs am Anfang ja auch gemacht, das war aber zu aufwändig und inperformant.

\section{Aufbau dieser Arbeit}
%... mach ich am Ende...

...Mach ich zum Schluss, aber nur wenn es die Betreuer auch haben möchten...






\chapter{Grundlagen}

\section{Serious Games in der Medizin}
%Was sind Serious Games, was gibts da so bis jetzt im medizinischen Bereich?

%was wird da benutzt? meist partikel basiert, technischen stand klären...

%Übersicht: \cite{SurgSim}

Unter Serous Games versteht man Videospiele die nicht der reinen Unterhaltung dienen, sondern primär zum Lernen und Trainieren eingesetzt werden (REF???).

Vor allem in der Chirurgie können Serious Games in Form von chirurgischen Simulationen eingesetzt werden, um Personal auf bestimmte Operationen vorzubereiten. Dabei geht es nicht nur darum, Operationsabläufe zu verinnerlichen. Auch die motorischen Fähigkeiten und das motorische Gedächtnis des Spielers wird dabei trainiert. 

%par Simulatoren als Beispiel:
In \cite{VRSim20} wird beispielsweise ein VR-Simulator für die neonatale endotracheale Intubation entwickelt, weil niedrige Intubationserfolgsraten darauf hinweisen, dass das derzeitige Trainingsprogramm nicht ausreicht, um positive Erfolgsraten zu erzielen (siehe Abbildung \ref{VRSim20}).

\bild{VRSim20}{14cm}{Aus \cite{VRSim20}. Neonatale endotracheale Intubation als VR-Simulation mit einem \ac{HMD} (\textit{HTC Vive Pro}). Der Nutzer hält ein reales Werkzeug (Lyaryngoskop) in der Hand, welches auch in der virtuellen Realität sichtbar ist. Trifft das virtuelle Lyaryngoskop auf simuliertes Gewebe, wird das haptische Gerät so angesteuert, dass der Spieler einen entsprechenden Widerstand wahrnehmen kann.}

Aufgrund des wachsenden Bedarfs an chirurgischen Simulatoren sind einige Firmen entstanden die sich auf dieses Gebiet spezialisiert haben. Beispielsweise stellt die Firma \textit{VRmagic} \cite{VRmagic} Simulatoren für die intraokulare Chirurgie her. In Abbildung \ref{SurSim2} ist der Simulator \textit{Eyesi Surgical} von VRmagic zu sehen. Wie man in der Abbildung rechts sieht, werden hierfür virtuelle Werkzeuge und ein virtuelles Auge fotorealistisch gerendert.

\bild{SurSim2}{14cm}{Aus \cite{VRmagic}. Links: Die Trainingsstation. Mitte: Position und Orientierung des Werkzeugs wird mittels entsprechender Sensorik erfasst. Rechts: Die Interaktion zwischen dem virtuellen Auge und dem Werkzeug wird in Echtzeit simuliert und fotorealistisch dargestellt.}

Die wichtigste Anforderung von solchen Simulatoren ist die Echtzeitfähigkeit. LESEZEICHEN!!! HIER WEITER!!!...

%...In \cite{SimRole} wird zwischen Physical models (tote tiere, puppen) Virtual reality simulators (stereoscpische goggles, haptic feedback) und Mixed-reality simulator ...



\section{Anforderungen}
\label{sec_Anforderungen}
Folgende Anforderungen werden an das zu entwickelnde System und dieser Arbeit gestellt: ...

\subsubsection{Interaktion und Echtzeit}
...

\subsubsection{Performance}

\subsubsection{Simulationsgenauigkeit}
Muss nicht so gut sein.

\section{Simulationstechniken und Physik Engines für chirurgische Simulatoren}
Engines: meist partikelbasiert und gpu-beschleunigt, PhysX, FleX, CUDA erwähnen, SOFA
\subsection{Position Based Dynamics}
\label{section_PBD}

Beliebt im medizinischem bereich weil:

- Einfach (einfacher zu implementieren und zu erweitern auch weil wirs mit punktmassen zutun haben, siehe kapitel xxx!), hohe Effizienz, Stabilität, Echtzeitleistung, 

im Gegensatz zu Kraftbasierte methoden: PBD projiziert Position direkt als Lösung (lösung von constraints) statt Beschleunigungen zu integrieren (um equilibrium confi zu erreichen). Dadurch beseitigt PBD das Überschießproblem. PBD ist nicht so genau wir kraftbasierte Methoden

Kann visuell plausible Ergebnisse liefern \cite{PBDKidney} \cite{BreastBiopsy}, obwohl ungenau .

Wurde schon erfolgreich eingesetzt für:

Simulation von chirurgischen Fäden (nähen, mit haptic feedback) \cite{PBDThread}, 
der Simulation von laparoskopischen Operationen in VR \cite{VRLaparoscop}, 
dem chirurgischen Schneiden \cite{PBDCutting}
und der Simulation von Roboteroperationen \cite{VRRobSim} (auch in VR).


\subsection{Erweiterungen von PBD}
\subsubsection{HPBD}

\subsubsection{Unified Particle Dynamics und NVIDIA Flex}
zu erwähnen über \cite{UPP} und Flex: 

Partikel sind über constraints miteinander verbunden.

Unified solver gut weil es einfacher zu handeln ist. Nur ein solver muss optimiert werden und vor allem weil dann objekte unterschiedlicher Art (soft, rigid, fluid) miteinander interagieren (collision constraint gillt für alle) können, aut of the box, in Echtzeit.

single particle radius sorgt für effizentere kollisionserkennen (wieso? nochmal nachlesen!)

Constaints: Distance(cloth) shape-matching(rigid, plastic/elastic deformation, siehe paper fig 5) density(fluid) volume(inflatables), friction(wichtiges constraint für nadel, flex is stolz auf friction constraint, siehe sand-teepot-demos), different collisions...

solver: PBD-solver, \cite{UPP} constraint-projektion als optimization problem, siehe kap 4.2, Gleichung (7). Bei jedem sim-step PBD versucht die Distanz zur Constraint-Verteilung (constraint manifold) zu minimieren unter berücksichtigung der partikel massen. Das tolle an dieser optimization sichtweise ist, das wir den PBD solver mit anderen solvern veergleichen können. So erkennt man das implizit euler recht ähnlich zu PBD ist. Unterschied: euler macht die minimization with respect zur initialen condition. bei BPD with respect to last iteration.

Parallelisierung mit gauss-jacobi solver (statt gaus-seidel aus PBD, der sequenziel arbeitet), der convergiert dann allerdings häufig nicht, z.b. wenn 2 partikel über identisches constraint verbunden sind. lösung: constraints paralell lösen, constraint deltas für jedes partikel summieren und durch n=anzahl constraints teilen. Convergiert gut. 





%\section{Verwendete Software}
%
%In diesem Kapitel werden alle Softwareelemente, die im Rahmen dieser Arbeit genutzt wurden, vorgestellt.
%
%\subsection{NVIDIA Flex}
%Das waren die verwendeten Software Werkzeuge. Während es zuvor über flex-theorie ging, wird hier kurz auf software-architektur von flex eingegangen (worauf achten, wie benutzen, Kompabilitäten, lizenz?...)
%
%\subsection{Unreal Engine 4 mit NVIDIA FleX-Integration}
%un joa... par worte über ue4.19 mit flex integration

%\section{Virtual Reality}
%
%Geräte werden immer günstiger usw ... finden wir vielleicht noch ein paper wo ein haptisches Gerät an einen Motion Controller befestigt wird? Klappt sowas gut?



%%%%%%%%%%%%%%%%%%%%%%% aus bachelor %%%%%%%%%%%%%%%%%%%%%%%%%%%%%%%%%%%%%%%%%%%%

%\subsection{Grundverhaltensregeln}
%
%Die Grundlage des Boid-Modells von Reynolds bilden drei einfache Verhaltensregeln, die jedes Boid befolgt \cite{Reynolds99steeringbehaviors}. Allein mit diesen Grundregeln können erste Bewegungen erzeugt werden, welche an Tierschwärme erinnern. 
%
%Jedes Boid besitzt dabei einen Wahrnehmungsbereich. Der Wahrnehmungsbereich ist ein kreis- oder kugelförmiger Bereich, in dessen Zentrum das zugehörige Individuum sitzt.
%
%\bild{BoidWahrnehmung}{5cm}{Ein möglicher Wahrnehmungsbereich eines Boids aus \cite{Reynolds99steeringbehaviors}. Hier besitzt jedes Individuum einen toten Winkel im Wahrnehmungsbereich.}
%
%Das Verhalten eines Boids wird durch die Zustände seiner benachbarten Schwarmmitglieder, die sich innerhalb seines Wahrnehmungsbereichs befinden, beeinflusst. Boids, die sich außerhalb des Wahrnehmungsbereichs befinden, werden nicht berücksichtigt.  Abbildung \ref{BoidWahrnehmung} zeigt einen solchen Wahrnehmungsbereich.
%
%
%\subsubsection{Kohäsion \textit{(Cohasion)}}
%Ein Boid hat stets das Bedürfnis, dicht in seiner Gruppe zu bleiben um vor Jägern geschützt zu sein. Dieses Verhalten wird mit der Grundregel \textit{Kohäsion} realisiert. In Abbildung \ref{Kohaesion} wird dieses Verhalten illustriert.
%
%\bild{Kohaesion}{7cm}{Das Grundverhaltensgesetz Kohäsion: Ein Individuum (grün) wird vom Schwerpunkt der Nachbarschaft angezogen \cite{Reynolds99steeringbehaviors}.} 
%
%Hierfür wird der relative Schwerpunkt aller benachbarten Boids ermittelt, um so einen Bewegungsvektor zu spannen der beschreibt, in welche Richtung sich das Individuum hingezogen  \glqq fühlt\grqq{}. Je weiter die Nachbarn entfernt desto größer wird der Vektor und auch die korrigierende Bewegungsgeschwindigkeit. Ein Individuum bewegt sich umso schneller zu seiner Gruppe hin, je größer die Gefahr ist, den Anschluss zum Schwarm zu verlieren.
%
%\subsubsection{Separation}
%
%Die Separation ist dazu da, um Kollisionen unter den Individuen zu vermeiden. Ein Individuum separiert sich, indem zuerst alle Nachbarn ermittelt werden, die sich innerhalb einer kritischen Zone befinden. Diese Nachbarn gelten dann als \glqq{}zu nah\grqq{}. In Abbildung \ref{Separation} wurden zum Beispiel drei Nachbarn gefunden, die dem grünen Boid zu nahe sind:
%
%\bild{Separation}{7cm}{Ein Boid separiert sich von seiner Nachbarschaft in dem ein abstoßender Vektor generiert wird, der von den Nachbarn weg zeigt \cite{Reynolds99steeringbehaviors}.}
%
%Danach wird für jeden dieser drei Nachbarn, ein Vektor gespannt, der von Ihnen weg und zum grünen Boid hin zeigt. Diese Vektoren werden anschließend normiert. Je kleiner der Abstand zu einem Nachbarn, desto größer sollte der abstoßende Vektor sein. Dies kann erreicht werden, indem jeder abstoßende Vektor umgekehrt proportional zum Abstand gewichtet wird.
%
%
%\subsubsection{Ausrichtung \textit{(Alignment)}}
%
%Die letzte Grundverhaltensregel von Reynold sorgt dafür, dass sich eine Gruppe aus Boids, als eine Einheit, in eine Richtung fortbewegt. Die Individuen einer Gruppe müssen dafür ihre Bewegungsrichtungen aneinander angleichen. Neben der Angleichung der Bewegungsrichtung kann auch zusätzlich eine Geschwindigkeitsangleichung erfolgen \cite{Reynolds99steeringbehaviors}.
%
%Um eine Ausrichtung zu erzeugen, sucht ein Individuum zuerst all seine Nachbarn. Dann wird der durchschnittliche Geschwindigkeitsvektor der Nachbarn ermittelt. Dieser Vektor beschreibt die derzeit angestrebte Richtung und Geschwindigkeit der Gruppe. Jedes Individuum berücksichtigt den Geschwindigkeitsvektor der Gruppe und versucht sich diesem anzupassen.
%
%\bild{Ausrichtung}{7cm}{Ein Boid richtet sich nach seinen Nachbarn aus \cite{Reynolds99steeringbehaviors}.} 
%
%Der finale Bewegungsvektor eines Individuums ist dann die Differenz zwischen Geschwindigkeitsvektor der Gruppe und des eigenen Bewegungsvektors. 
%
%
%
%\subsection{Zusatzvorschriften}
%\label{subsec_Zusatzvorschriften}
%
%Meist werden neben den drei Grundverhaltensregeln noch weitere Zusatzvorschriften implementiert um dem Schwarm noch mehr Fähigkeiten zu verleihen. Zu diesen Zusatzvorschriften zählen beispielsweise \glqq Hindernis vermeiden\grqq{}, \glqq Pfad folgen\grqq{} oder \glqq Flussfeld folgen\grqq{} \cite{Reynolds99steeringbehaviors}. 
%
%Das Verhalten der Boids lässt sich aber auch mit beliebigen, selbst entwickelten Zusatzvorschriften erweitern. So wurde beispielsweise in \cite{Oliver14Tschesche} dafür gesorgt, dass sich alle Individuen zu einem zentralen Punkt hingezogen fühlen, damit der Schwarm nicht \glqq davon fliegt\grqq{}.
%
%Im Rahmen dieser Arbeit wurden die folgenden Zusatzverhaltensregeln umgesetzt:
%
%\subsubsection{Position folgen}
%Es kann ein Punkt im Raum definiert und umher bewegt werden. Der Schwarm wird von diesem Punkt angezogen. Sobald dieser Punkt erscheint, vernachlässigen die Individuen die drei Grundverhaltensregeln ein wenig um den Punkt anzufliegen. Der anziehende Punkt könnte beispielsweise Futter repräsentieren und wird \textit{Kohäsionspunkt} genannt.
%
%
%\subsubsection{Position meiden}
%Mit dieser Verhaltensregel kann ein menschlicher Spieler den virtuellen Schwarm \glqq jagen\grqq{}. Im Gegensatz zur Verhaltensregel \textit{Position folgen} wird hier ein Abstoßender Punkt im Raum definiert. Die Boids meiden diesen Ort und bewegen sich von ihm weg. Dieser Punkt wird \textit{Separationspunkt} genannt.
%
%
%\subsection{Kombinierung von Verhaltensregeln}
%
%Jede eingesetzte Verhaltensregel generiert einen Vektor, der beschreibt, in welche Richtung sich das betrachtete Individuum bewegen möchte. Diese Vektoren werden auch als \textit{Steuersignale} oder \textit{Steuervektoren} bezeichnet \cite{Ulm01Maschke}.  Jede Verhaltensregel, $i$, generiert  für jedes Individuum ein Einzelsteuersignal, $\vec{v_i}$. Aus $N$ Verhaltensregeln entstehen also $N$ Einzelsteuersignale. Alle Einzelsteuervektoren, $\vec{v_i}$, werden am Ende zu einem finalen Bewegungsvektor, $\vec{V}$, zusammengefasst (siehe Gleichung \ref{V}). Die finalen Steuervektoren, $\vec{V}$, bewegen schlussendlich die Individuen. 
%
%$\vec{V} = f(\vec{v_i})$ kann auch wie die finale Entscheidung, in welche Richtung es gehen soll, verstanden werden, während $\vec{v_i}$ das Nachdenken repräsentiert. Jedes $\vec{v_i}$ schätzt während des  \glqq Überlegungs-Prozesses\grqq{} eine bestimmte Situation ein. Beispielsweise wie hoch die Gefahr einer Kollision
%ist (Separation) oder ob die Gefahr besteht, den Anschluss zum Schwarm zu verlieren (Kohäsion). Je größer ein Steuervektor, $\vec{v_i}$, desto dringlicher ist der Handlungsbedarf.
%
%Bei der Berechnung des finalen Bewegungsvektors, $\vec{V}$, wird mit Prioritätsgewichten, $w_i$, gearbeitet, um einstellen zu können, wie stark die jeweiligen Verhaltensregeln, $i$, ausgeprägt sein sollen. $\vec{V}$ wird dann durch Bildung des gewichteten Mittelwerts aus allen Steuersignalen gebildet:
%
%\begin{equation}
%\vec{V} = \dfrac{1}{N} \cdot \sum_{i=1}^N(w_i \cdot \vec{v_i})
%\label{V}
%\end{equation}
%
%So entsteht ein finaler Steuervektor, $\vec{V}$, der für die Bewegung eines einzelnen \mbox{Boids}  verantwortlich ist. Formel \ref{V} kann auch anders gestaltet werden. Wichtig ist dabei nur, dass die Einzelsteuersignale mithilfe von Gewichtungen parametrisiert und miteinander kombiniert  werden. Es kann auch lediglich die gewichtete Summe aus $\vec{v_i}$ gebildet werden. Reynolds schlägt außerdem vor, die Steuervektoren, $\vec{v_i}$, zu normieren. Diese Variante kann für eine einfachere Kontrolle des Schwarms sorgen und die Parametrisierungsarbeit erleichtern \cite{Reynolds99steeringbehaviors}.
%
%Mit dieser Methodik kann ein Boid im übertragendem Sinne auch mehrere Entscheidungen zugleich treffen, weil die finale Entscheidung, $V$, immer von allen Verhaltensregeln, $i$, abhängig ist. Dafür wird der algorithmische Aufwand für das Treffen von Entscheidungen minimiert. Statt auf einen bedinungsabhängigen, stark verzweigten Algorithmus zu setzen, wird hier eine Entscheidung mithilfe einer einfachen Formel gefällt. Für den Entscheidungsprozess werden keine Programmschleifen oder Abzweigungen benötigt, nur Formel \ref{V}. Daher ist diese Methode gut für Grafikkarten geeignet, weil diese am besten als reine Rechenwerke genutzt werden sollten (siehe Kapitel \ref{subsec_CDMR}).
%
%\subsection{Bewegung der Boids}
%
%Die Bewegung eines Boids, das sich an einer Position $P$ befindet, kann beispielsweise nach Formel \ref{move} erfolgen \cite{Ulm01Maschke}. 
%
%\begin{equation}
%P(t+dt) = P(t) + \vec{V(t)} \cdot dt
%\label{move}
%\end{equation}
%
%Gleichung \ref{move} lässt sich allerdings nicht ohne weiteres von Shadern umsetzen. Das liegt vor allem an der Funktionsweise von Shadern. Ausführliche Erläuterungen über diese Problematik folgen in Kapitel \ref{section_BewDerBoids}. In Kapitel \ref{section_BewDerBoids} wird auch eine alternative Methode für die Bewegung der Boids konzipiert.
%
%
%
%
%
%
%
%
%%=====================UE4===============================UE4=================
%\section{Die Unreal Engine 4}
%\label{section_ue4}
%Für die Umsetzung dieser Arbeit wurde die Unreal Engine 4 von Epic Games verwendet. 
%Sie bietet viele Hilfestellungen in Form von fertigen Framesworks, Softwarebibliotheken und 3D-Editoren. 
%
%Die Unreal Engine 4 leistet zudem Hilfe bei der plattformübergreifenden Entwicklung. Sie unterstützt einige Betriebssysteme, wie Linux, Windows, IOS oder Android sowie alle aktuellen Videospielkonsolen von Microsoft, Sony und Nintendo.
%
%Mithilfe der \ac{UE4} können auch sogenannte Compute-Shader-Pipelines erstellt werden. So können Simulationsberechnungen auf die Grafikkarte verlagert werden, um von ihrer parallelisierten Rechenleistung zu profitieren. In Kapitel \ref{subsec_CDMR} wird diese Thematik näher erläutert.
%
%
%
%\section{Shader-Entwicklung mit Unreal Engine 4}
%\label{sec_UE4UndShader}
%
%Ein Shader-Programm ist meistens dafür vorgesehen die Oberfläche eines 3D-Modells wie ein Material aussehen zu lassen. Daher wird ein Shader häufig als  \glqq Material\grqq{} bezeichnet. Ein Material ist ein Shader-Programm, das dafür vorgesehen ist, 3D-Modelle in unterschiedlichster Weise zu zeichnen \cite{Max17Legnar}.
%
%\bild{MaterialShader}{7cm}{Ein Shader, oder auch Material genannt, wird auf ein 3D-Modell angewandt, damit dessen gezeichnete Oberfläche an einen Felsen erinnert \cite{Max17Legnar}.}
%
%Shader-Programme werden in der Unreal Engine 4 mithilfe des sogenannten Material Editors erstellt. Dieser besteht aus einer rein grafischen Programmierumgebung, welche auf das Verbinden unterschiedlicher \glqq Material-Knoten\grqq{}  basiert \cite{UE4MatNodes}.
%
%Ein mit der Unreal Engine erstellter Shader wird von der Engine automatisch in \ac{HLSL} übersetzt. \ac{HLSL} ist die Shader-Sprache der Grafikkartenschnittstelle DirectX. Der übersetzte HLSL-Programmcode kann dem Programmierer bei Bedarf auch angezeigt werden. Ein in der Unreal Engine integrierter HLSL Cross Compiler sorgt für eine erweiterte Kompatilität eines Shaders \cite{CrossCompiler}. Er  optimiert generierten \ac{HLSL}-Programmcode  und übersetzt ihn in GLSL-Programmcode (\ac{GLSL}). GLSL ist die Shader-Sprache von OpenGL und Vulkan. Diese Grafikkartenschnittstellen unterstützen unter anderem das Betriebssystem Linux und Sonys Playstation 4, was DirectX nicht tut. Somit muss der Programmierer nur einen Shader programmieren. Dieser kann dann von besonders vielen unterschiedlichen Plattformen ausgeführt werden.
%
%Im Material-Editor der UE4 können auch neue Material-Knoten per \ac{HLSL} erstellt werden \cite{CustomNode}. Diese Knoten werden als \textit{Custom Nodes} bezeichnet und halten dem Programmierer weitere Möglichkeiten offen. Es können nämlich nicht alle Funktionalitäten, die HLSL bietet, mit den zur Verfügung gestellten Material-Knoten umgesetzt werden. Daher wurden die Custom Nodes eingeführt, um alle Möglichkeiten offen zu halten. Nichts desto trotz rät Epic Games von der Verwendung der Custom Nodes ab. Nicht nur weil diese unsicherer sind, sondern auch weil ihr beinhalteter HLSL-Programmcode nicht weiter von der Unreal Engine optimiert werden kann \cite{CustomNode}.
%
%Mithilfe eines Custom Nodes können beispielsweise fortgeschrittene Renderverfahren implementiert werden, so wie Ryan Brucks in \cite{Ryan18Brucks} zeigt.
%
%
%
%\subsection{Umgang mit Texturen in der Shader-Programmierung}
%\label{subsec_BilFiltUV}
%
%\subsubsection{UV-Koordinaten}
%
%Wesentlicher Bestandteil der Shader-Programmierung ist die Verarbeitung von Texturen. Häufig muss dabei auf ein bestimmtes Pixel, das sich in einer Textur befindet, zugegriffen werden. Der Zugriff auf ein Pixel erfolgt über eine Positionsangabe. Die Positionen der Pixel werden von sogenannten UV-Koordinaten beschrieben.
%
%\bild{UV1}{6cm}{Eine Textur kann man sich wie ein Gitter-Netz aus Zellen vorstellen. Jede Zelle besitzt einen Farbvektor, der sich genau in ihrem Zentrum befindet \cite{BilTexFilt}.}
%
%In Abbildung \ref{UV1} ist eine Textur, $T$, mit einer Auflösung von $R=4$ mal $R=4$ Pixeln abgebildet. Eine Textur kann man sich wie ein Gitternetz vorstellen, wobei jedes Quadrat dieses Netzes häufig als \textit{Zelle} oder \textit{Texel} bezeichnet wird. Jede dieser Zellen besitzt einen Farbwert, der in Abbildung \ref{UV1} als farblicher Punkt dargestellt wird. 
%
%Farbwerte werden mit 3D-Vektoren, in Form von $(R \quad G \quad B)$, beschrieben. Die Komponenten eines Farbvektors beschreiben den Rotanteil, $R$, Grünanteil, $G$ und Blauanteil, $B$, der Farbe. Jede Komponente hat meist einen Wert zwischen 0 und 255, wobei 255 ein voller Farbanteil ist. Es gibt allerdings auch einige andere Farbformate.
%
%Der Ursprung des \textit{UV}-Koordinatensystems liegt immer in der oberen linken Ecke der Textur. An dieser Stelle befindet sich die Position $UV=(0 \quad 0)$. Unabhängig von der Auflösung, $R$, der Textur, wird die untere rechte Ecke immer mit $UV=(1 \quad 1)$ beschrieben. Dementsprechend ist beispielsweise das Zentrum der Textur immer bei $UV=(0,5 \quad 0,5)$.
%
%Das UV-Koordinatensystem einer Textur reicht also immer von 0 nach 1 und es ist linear. Jede Zelle einer quadratischen Textur besitzt immer eine Breite und Höhe von $1/R$. Die Größe $1/R$ wird häufig auch als \textit{Texel-Size} bezeichnet. An dieser Größe kann man sich gut orientieren um gezielt durch eine Textur zu navigieren. 
%
%Hierzu ein Beispiel: Mit dem Wissen, dass jede Zelle $1/R$ hoch und breit ist, kann man schlussfolgern, dass sich der Mittelpunkt der oberen linken Zelle, von $T$, genau bei $UV = \begin{pmatrix} \dfrac{1}{R} \cdot \dfrac{1}{2} \quad \dfrac{1}{R} \cdot \dfrac{1}{2}  \end{pmatrix}$ befindet (Siehe Abbildung \ref{UV1}).
%
%
%\subsubsection{Bilineare Filterung}
%
%Mit dem UV-Koordinatensystem kann ohne weiteres auf einen Ort, der genau zwischen vier benachbarten Pixeln liegt, zugegriffen werden. Dabei stellt sich die Frage, was für ein Wert in diesem Fall ausgelesen werden würde. Die Antwort auf diese Frage verbirgt sich hinter der bilinearen Filterung \cite{BilTexFilt}. 
%
%Jeder ausgelesene Farbwert wird von der Unreal Engine 4 automatisch bilinear gefiltert.
%Ein bilinear gefilterter Farbwert ist prinzipiell nichts anderes als ein gewichteter Mittelwert von 4 benachbarten Pixeln. Hierzu ein Rechenbeispiel:
%
%Angenommen es soll auf Position $\vec{x}=(0,5  \quad 0,5)$ von $T$ zugegriffen werden. Dieser Ort liegt exakt in der Mitte von 4 benachbarten Pixeln (Siehe Abbildung \ref{UV1}). Der Auslesepunkt hat also zu allen 4 Pixeln die selbe Entfernung. Deswegen wird der gleichmäßig gewichtete Mittelwert aus den Farbwerten der 4 Pixel ausgelesen. Der bilinear gefilterte Wert ist dann:
%
%\begin{equation}
%T(\vec{x}) = 
%  0,25 \cdot \begin{pmatrix} 255 \\ 0 \\ 0 \end{pmatrix}
%+ 0,25 \cdot \begin{pmatrix} 0 \\ 255 \\ 0 \end{pmatrix} 
%+ 0,25 \cdot \begin{pmatrix} 0 \\ 0 \\ 255 \end{pmatrix}
%+ 0,25 \cdot \begin{pmatrix} 255 \\ 255 \\ 255 \end{pmatrix}
%\label{Form_LinFilterMitte}
%\end{equation}
%
%%\begin{equation}
%%\vec{e}_1=\left(\begin{array}{c} 1 \\ 0 \end{array}\right) \qquad
%%\vec{e}_1=\left(\begin{array}{c} 1 \\ 0 \end{array}\right)
%%\label{Formel}
%%\end{equation}
%
%Formel \ref{Form_LinFilterMitte} bildet eine gewichtete Summe aus allen Farbvektoren. Dabei ist die Summe aller Gewichte immer 1. In diesem Fall wurden alle Farbwerte gleich stark gewichtet, weil $x$ zu allen 4 Pixeln dieselbe Entfernung hatte. Wenn $x$ beispielsweise etwas näher bei der roten Zelle wäre, würde dieser Farbwert stärker gewichtet werden, als die anderen 3 Farbwerte. Wenn $x$ genau auf den Mittelunkt einer Zelle zeigt, wird Ihr Farbwert mit 1 und alle anderen Farbwerte mit 0 gewichtet. Das heißt: Ist der exakte Farbwert eines Pixels gefragt, muss genau auf dessen Mittelpunkt zugegriffen werden. Genau auf den Mittelpunkt einer Zelle zuzugreifen ist allerdings nicht immer möglich, weil eine digital arbeitende Maschine (die GPU) keine unendlich genauen Zahlen darstellen kann.
%
%Die in Abbildung \ref{UV1} dargestellte Textur sieht, nachdem sie vollständig bilinear gefiltert wurde, wie folgt aus:
%
%\bild{gefilterteTextur}{4cm}{Eine vollständig bilinear gefilterte Textur \cite{BilTexFilt}.}
%
%Wie man in Abbildung \ref{gefilterteTextur} sieht, ergibt sich somit eine gewisse Unschärfe.
%
%\section{Echtzeitsimulationen mit Content-Driven Multipass Rendering}
%\label{subsec_CDMR}
%
%\subsection{Content-Driven Multipass Rendering}
%
%\ac{CDMR} ist eine Technik, mit der rechenintensive Simulationen von herkömlichen Shadern ausgeführt werden können. Im Zusammenhang mit der Unreal Engine 4, wurde \ac{CDMR} erstmals auf der \textit{Game Development Converence 2017}, von Ryan Brucks vorgestellt  und erklärt \cite{GDC2017}. 
%
%Bei \ac{CDMR} werden Shader nicht genutzt um 3D-Modelle darzustellen. Stattdessen werden sie als reine Rechenwerke eingesetzt, die nichts anderes tun, als Texturen miteinander zu verrechnen. Ein Shader, der nur zum Lösen von Rechenaufgaben da ist, wird auch als \textit{Compute-Shader} bezeichnet und nicht als \textit{Material}. Ein Compute-Shader sollte im Idealfall die GPU als reines Rechenwerk nutzen und dementsprechend nichts anderes tun, als eine große Rechenaufgabe zu lösen. Herkömmliche Programmschleifen oder Verzweigungen sollten im Programmcode eines Shaders vermieden werden \cite{CustomNode}.
%
%Abbildung \ref{CDMRexplain2} zeigt über welche Objekte eine Zusammenarbeit zwischen \ac{CPU} und GPU zustande kommt. Alle im Schaubild grün eingefärbten Daten und Programme befinden sich im Speicher der GPU und werden von ihr verarbeitet. Die blau markierte Raute repräsentiert hingegen eine Klasse, die von der CPU ausgeführt wird. Der Begriff \textit{Blueprint} kommt aus der Unreal Engine-Programmierung und ist eine auf C++ aufbauende Programmiersprache.
%\clearpage
%
%\bild{CDMRexplain2}{12cm}{Dem Programmierer werden von der Unreal Engine einige Objekte zur Verfügung gestellt, mit dessen Hilfe interaktive Shader-Effekte realisiert werden können. Dieser Graph zeigt schematisch, wie die CPU über diese Objekte mit der GPU zusammenarbeiten kann.}
%
%
%Wie in Abbildung \ref{CDMRexplain2} visualisiert, kann ein Shader ein oder mehrere Texturen einlesen. Ein Shader verrechnet eingelesene Texturen in irgendeiner Form, sobald dieser von der CPU ein Ausführungssignal bekommt. Die durch die Verrechnung entstehenden Ergebnisse werden dann in eine Ergebnis-Textur geschrieben. 
%
%Neben Texturen kann ein Shader auch Vektor- oder Skalar-Parameter einlesen. Mit ihnen lassen sich die vom Shader durchgeführten Berechnungen entsprechend anpassen.
%
%Die \ac{CPU} kann einen Compute-Shader steuern, indem sie ihn mit den in Abbildung \ref{CDMRexplain2} gezeigten Mitteln beeinflusst. Alle Einflussmöglichkeiten werden von blauen Feilen repräsentiert. Eine Blueprint-Klasse kann genau bestimmen, wann ein Compute-Shader ausgeführt werden soll. Wird ein Shader aufgerufen, führt er seine Berechnungen durch und schreibt seine Ergebnisse in seine zugeteilte Ergebnis-Textur. Um eine Echtzeitsimulation für Videospiele zu realisieren, sollte ein Compute-Shader mindestens 30 mal pro Sekunde ausgeführt werden können, damit für den Betrachter flüssige Bewegungen zustande kommen. Dabei muss natürlich berücksichtigt werden, dass in einem Computerspiel noch einige weitere Shader und Unterprogramme Ressourcen und Rechenzeiten in Anspruch nehmen.
%
%Des weiteren hat die CPU die Möglichkeit beliebige Texturen, die ein Shader nutzt, mit anderen Texturen auszutauschen. Auch die Werte der Shader-Parameter können jederzeit verändert werden. All dies kann zur Laufzeit der Anwendung geschehen. So können interaktive Shader-Effekte realisiert werden.
%
%Alle in Abbildung \ref{CDMRexplain2} dargestellten Zusammenhänge stellen alle nötigen Voraussetzung für \ac{CDMR} dar. Mit diesen Möglichkeiten können nämlich sogenannte Shader-Pipelines erstellt werden, was die Grundlage von \ac{CDMR} ist.  Hierzu ein Beispiel aus \cite{Max17Legnar}:
%
%
%\bild{ShaderPipelineBsp}{14cm}{Eine Shader-Pipeline aus \cite{Max17Legnar} als Beispiel. Mehrere Shader können in einer Pipeline-Struktur zusammenarbeiten, indem sie ihre Ergebnisse über Texturen miteinander teilen.}
%
%
%Die in Abbildung \ref{ShaderPipelineBsp} dargestellte Shader-Pipeline entstand in \cite{Max17Legnar} um eine interaktive Wasseroberfläche zu simulieren. Alle darin zu sehenden Rechtecke repräsentieren Texturen und jeder Kreis ist ein Shader. Die Shader sind nummeriert, um zu markieren in welcher Reihenfolge die jeweiligen Shader ausgeführt werden. Abbildung \ref{ShaderPipelineBsp} soll nur das grundlegende Prinzip einer Shader-Pipeline veranschaulichen. Ein Shader beginnt mit seiner Arbeit und verrechnet Texturen miteinander um seine Rechenergebnisse in eine Ergebnistextur zu speichern. Danach kann ein weiterer Shader diese Ergebnistextur einlesen und so das Ergebnis weiterverarbeiten. Mit diesem Grundprinzip kann eine Verarbeitungskette aus mehreren hintereinander geschalteten Shadern gebildet werden. Solch eine Kette wird in dieser Arbeit als Shader-Pipeline bezeichnet und beschreibt das Grundprinzip von \ac{CDMR} auf einfache Weise.
%
%Eine  Shader-Pipeline lässt sich besonders flexibel von der CPU steuern. Mithilfe einer Blueprint-Klasse kann zur Laufzeit der Simulation, die Reihenfolge, in der die Shader ausgeführt werden, verändert werden. Genauso können auch jeder Zeit zusätzliche Shader zur Shader-Pipeline hinzu geschaltet werden. 
%
%\ac{CDMR} bringt allerdings auch Einschränkungen mit sich. Ein Shader kann niemals dieselbe Textur beschreiben, die er gerade ausliest. Daher muss in einigen Fällen mit zwischenkopierten Texturen gearbeitet werden. Zwar profitiert die Shader-Pipeline von der starken Rechenleistung der GPU, dafür kann sie aber nur Texturen verarbeiten. Das Endergebnis der Pipeline liegt immer in Texturform vor. 
%
%Um ein effizientes System zu erhalten, sollte die Ergebnis-Textur einer Shader-Pipeline nicht den Grafikkartenspeicher verlassen. Das Ergebnis der Pipeline sollte am besten bei der GPU bleiben und direkt in irgendeiner Form gezeichnet werden. Wie die Darstellung einer Ergebnis-Textur erfolgen kann, wird im nächstem Kapitel erläutert.
%
%\subsection{Scatter oder Gather}
%\label{subsec_ScatterVSGather}
%
%Die sogenannte \textit{Scatter VS Gather}-Problematik wurde von Ryan Brucks in  \cite{GDC2017}, im Zusammenhang mit \ac{CDMR} thematisiert.
%
%Bei dieser Problematik geht es vor allem darum, zu verstehen, wie eine Grafikkarte arbeitet, wenn sie Shader ausführt:
%
%Wenn ein Shader eine Textur verarbeitet, kann man sich vorstellen, dass dabei aus jedem Pixel der verarbeiteten Textur ein Thread entsteht. Alle Threads werden dabei von der GPU parallel ausgeführt. Jedes Pixel, hier auch meist als Zelle bezeichnet, ist also ein eigenständiger Thread. 
%
%Aus diesem Grund wird in dieser Arbeit häufig mit Begriffen wie  \glqq Pixel-Thread\grqq{} oder  \glqq Zellen-Thread\grqq{} gearbeitet.
%
%Natürlich entsteht in der Praxis nicht immer aus jedem Pixel ein eigenständiger Thread, weil nicht immer genügend Hardware-Ressourcen zur Verfügung stehen. Dennoch kann man sich vereinfacht vorstellen, alle Pixel einer Textur arbeiten gleichzeitig und jedes Pixel habe dabei seinen  \glqq eigenen Kopf\grqq{}. 
%
%Jeder Pixel-Thread führt Berechnungen durch um dann sein Rechenergebnis zurückzugeben, indem er seine eigene Farbe ändert. Der Farbwert eines Pixels repräsentiert also das Rechenergebnis eines Threads. Dabei gilt die Regel, dass ein Pixel-Thread sein Rechenergebnis immer nur \glqq in sich selbst\grqq{} hinein schreiben kann. Ein Pixel-Thread kann immer nur seinen eigenen Zustand (bzw. Farbwert) ändern. 
%
%Also gibt es bei der Shader-Programmierung die folgende Einschränkung:
%Ein Pixel-Thread kann niemals benachbarte Pixel beschreiben. 
%
%Diese Einschränkung ist allerdings auch notwendig. Wenn alle Pixel-Threads aus einer Textur, Werte in benachbarte Pixel schreiben würden, könnte eine Situation entstehen, in der mehrere Pixel versuchen in dasselbe Pixel zu schreiben. Hinzu kommt, dass alle Pixel-Threads gleichzeitig arbeiten. So würde ein unlösbarer Konflikt entstehen.
%
%Wegen der Funktionsweise von Shadern sollten Compute-Shader-Algorithmen stets mithilfe von sogenannten Gather-Methoden umgesetzt werden. Ein klassischer zellulärer Automat ist ein gutes Beispiel eines Systems, das nach der Gather-Methode arbeitet. Eine Zelle, die ihren eigen Zustand in Abhängigkeit von ihren eigenen und den Zuständen ihrer Nachbarschaft ändert, arbeitet mit einer Gather-Methode. Genau dieses Verhalten kann von Shader-Programmen nachgeahmt werden. Abbildung \ref{ScatterVSGather} veranschaulicht dieses Prinzip.
%
%\bild{ScatterVSGather}{13cm}{Der Unterschied zwischen der Scatter- und Gather-Methode. Bei der Scatter-Methode beeinflusst eine Zelle die Zustände von anderen Zellen. Bei der Gather-Methode liest eine Zelle andere Zellen aus, um in Abhängigkeit von diesen Informationen ihren eigenen Zustand zu überschreiben. Shader arbeiten eher nach der Gather-Methode. \cite{GDC2017}}
%
%Was jedoch nicht von Shadern nachgeahmt werden kann, ist die sogenannte Scatter-Methode. Bei  dieser Methode ändert eine Zelle die Zustände benachbarter Zellen. Diese Methode sollte vermieden werden, weil ein Shader nicht nach diesem Prinzip arbeitet.
%
%Die \textit{Scatter oder Gather}-Problematik ist bei der Konzipierung eines Compute-Shaders stets zu berücksichtigen. Manch eine Lösungsmethode, die einem zuerst in den Sinn kommt, lässt sich im Nachhinein als Scatter-Methode identifizieren. In solchen Fällen muss umgedacht werden, um die gegebene Problemstellung mithilfe einer Gather-Methode zu lösen. Nur Gather-Methoden lassen sich problemlos und effizient von Shadern umsetzen.
%
%
%
%
%
%
%\section{Texture-Based Volume Rendering}
%\label{subsec_VolumetricRayMarching}
%
%Weil die Simulation mit Content-Driven Multipass Rendering durchgeführt wird, wird der Schwarm in Texturform vorliegen. Dabei stellt sich die Frage wie der simulierte Schwarm dargestellt werden könnte. Eine empfohlene Möglichkeit ist \textit{Texture-Based Volume Rendering} \cite{GPUGems1} \cite{Drebin:1988:VR}. 
%
%Mit dieser Technik kann eine 3D-Textur verwendet werden um dreidimensionale Objekte zu zeichnen. Der in dieser Arbeit simulierte Schwarm könnte mithilfe von \textit{Texture-Based Volume Rendering} gezeichnet werden.
%
%Diese Render-Technik wird auch \textit{Volumetric Ray Marching} genannt, weil das Funktionsprinzip darauf basiert, Strahlen durch eine 3D-Textur hindurch \glqq marschieren\grqq{} zu lassen, um ein Bild zu zeichnen. Dieses Prinzip wird in Abbildung \ref{VolRendering} dargestellt \cite{Ryan18Brucks}.
%
%
%\bild{VolRendering}{11cm}{Das Ray Marching-Verfahren in der Draufsicht. Mehrer Strahlen  \glqq marschieren\grqq{} durch eine 3D-Textur, wobei die Schrittgröße der Strahlen konstant ist. Die 3D-Textur beschreibt ein Objekt mithilfe von Skalaren. Jeder Wert repräsentiert die Dichte des Objekts an einer Stelle im Raum \cite{Ryan18Brucks}.}
%
%Die von Strahlen durchwanderte 3D-Textur repräsentiert dabei einen Körper, der an unterschiedlichen Stellen unterschiedlich dicht ist. Je größer ein Skalar-Wert in der 3D-Textur, desto dichter ist das gezeichnete Objekt an dieser Stelle und desto mehr Licht wird an dieser Stelle absorbiert. In der Spieleprogrammierung werden mit dieser Technik häufig Wolken- oder Rauch-Effekte gezeichnet.
%
%\bild{Smokeball}{7cm}{Mit dem volumetrischen Renderer von Ryan Brucks können unter anderem dreidimensionale Wolken gezeichnet werden \cite{Ryan18Brucks}.}
%
%Ryan Brucks hat bereits einen volumetrischen \textit{Ray Marcher} mithilfe der unreal Engine 4 implementiert. Dieser Renderer ist frei verfügbar, außerdem wird in \cite{Ryan18Brucks} ausführlich erklärt, wie dieser funktioniert.
%
%

\chapter{NVIDIA FleX und Unreal Engine 4 für Serious Games}
\label{chap_Flex_Engine}

In Kapitel \ref{chap_Flex_Engine} wird die verwendete Unreal Engine 4 mit FleX Integration evaluiert und wenn nötig verbessert. Die Evaluation dient dazu, einen Eindruck davon zu bekommen, wie gut die Software für die Entwicklung von Serious Games im medizinischen Bereich geeignet ist und ob noch fehlende Funktionalitäten hinzugefügt werden müssen. 

Wenn eine fehlende Funktionalität entdeckt wird, die für die Entwicklung von Serious Games unverzichtbar ist, wird direkt gezeigt, wie die fehlende Funktionalität implementiert wurde.

%nicht ändern, nur erweitern:
%Im Rahmen dieser Arbeit wird stets darauf geachtet, den Source Code der Unreal Integration nicht zu verändern. 
In dieser Arbeit wurden fehlende Funktionalitäten nur durch Hinzufügen neuer Softwareelemente realisiert, nicht durch Veränderung der zugrunde liegenden Software. So bleibt die Software modular und neue Versionen der Unreal Integration können problemlos verwendet werden, durch einfaches Austauschen. 
%So soll die Qualitätsanforderung \textit{Modularität} eingehalten werden (siehe Kapitel \textbf{XXX}).


\section{Performance}
\label{sec_performance}

%\subsubsection{NVIDIA FleX}
%NVIDIA FleX zeichnet sich laut diversen Quellen (\cite{FlexD3D}, \cite{UPP}) durch eine besonders gute Performance aus.
%, wobei in \cite{FlexD3D} genauer darauf eingegangen wird, wie schnell FleX unterschiedliche Constraints lösen kann.

%In diesem Kapitel wird die Performance von NVIDIA FleX, zusammen mit Unreal Engine 4, untersucht.

Laut \cite{FlexD3D} wirkt sich die Partikelanzahl weniger stark auf die Performance aus wie die Anzahl und Komplexität der verwendeten Constraints, wobei in \cite{FlexD3D} genauer darauf eingegangen wird, wie schnell FleX unterschiedliche Constraints lösen kann.
%Befinden sich beispielsweise nur Soft Bodies in einer Szene, muss der Solver deutlich weniger unterschiedliche Constraints lösen als in einer Szene, in der sich Soft Bodies, Flüssigkeiten und Stoffe befindet.

%Die Constraints \textbf{X, X und X} sind beispielsweise deutlich rechenaufwändiger, als ein einfaches \textbf{X-Constraint} (siehe \cite{FlexD3D}).% und wirken sich daher stärker auf die Rechenzeit aus.

Im Rahmen dieser Arbeit wurde die Performance von NVIDIA FleX in Verbindung mit Unreal Engine 4 ausführlich untersucht. In diesem Kapitel (\ref{sec_performance}) werden die Einflussgrößen präsentiert, die den größten Einfluss auf die Performance hatten.

\subsubsection{Verwendete Hardware}

Für diese Arbeit und für alle durchgeführten Messungen und Tests stand die Hardware aus Tabelle \ref{table_hardware} zur Verfügung.

%\clearpage

\begin{table}[hbt!]
\centering
\caption{Die verwendete Hardware.}
\label{table_hardware}
\begin{tabular}{l|l}
 \hline
CPU & Intel Core i7-9700K  \\ \hline
GPU & NVIDIA GTX 1070 - 8GB GDDR5      \\ \hline
RAM & 32GB DDR4            \\ \hline
\end{tabular}
\end{table}

%Die verwendete Grafikkarte gehört zur vorletzten Generation der NVIDIA Grafikkarten und ist somit weniger aktuell und nicht übermäßig stark.

%% Please add the following required packages to your document preamble:
%% \usepackage{booktabs}
%\begin{table}[]
%\centering
%\caption{Einige Boid-Erhaltungsmessungen. Der Verlust von Boids konnte eindeutig nachgewiesen werden. Vor allem bei der Verteilung in zwei Zellen traten massive Verluste auf. }
%\label{Erhaltungsmessungen}
%\begin{tabular}{@{}llll@{}}
%\toprule
%\begin{tabular}[c]{@{}l@{}}Anzahl der\\ Wahrnehmungsvektoren\end{tabular} & Transportmethode                                                               & \begin{tabular}[c]{@{}l@{}}absoluter \\ Verlust\end{tabular} & \begin{tabular}[c]{@{}l@{}}prozentualer\\ Verlust\end{tabular} \\ \midrule

\subsection{Performance Tests}
\label{subsec_perf_tests}
\subsubsection{Gemischte Simulationen}

Um einen ersten Überblick von der Performance zu bekommen, wird zunächst eine Szene simuliert, in der sich mehrere unterschiedliche Objektarten befinden. Dadurch wird der FleX-Solver stark belastet, weil er viele unterschiedliche Constraints lösen muss. Dazu gehören beispielsweise Shape Constraints (Soft- und Rigid Bodies), Spring-Constraints (Stoff-Objekte), Flüssigkeiten (Density-Constraints), sowie Kollisions-Constraints und Friction-Constraints.

%\bild{WaterInteraction}{10cm}{Mehrere Objektarten interagieren auf realistische Weise miteinander. Werkzeuge (Skalpelle) werden von der Flüssigkeit weggespült und Körper, die eine geringere Dichte als die Flüssigkeit haben, schwimmen.}

\bild{WaterShaderPerformance}{14cm}{Mehrere Objektarten interagieren auf realistische Weise miteinander. Wird ein aufwändiger Shader für die Fluid-Partikel verwendet, kommt die GPU schnell an ihre Grenzen und die \ac{FPS} werden deutlich verringert.}

Die Simulationsparameter werden außerdem so gewählt, dass die Simulation möglichst akkurat ausgeführt wird, wodurch der FleX-Solver ebenfalls stärker belastet wird. Die gewählten Simulationsparameter wurden so wie in \cite{BreastBiopsy} gewählt. Diese Parametrisierung hat sich nämlich auch in \cite{PBDKidney} und \newline \cite{VRSim20} bewährt, um Gewebe akkurat zu simulieren.

Die verwendeten Simulationsparameter sowie alle Messergebnisse sind in Tabelle \ref{table_mixed_sim} zu sehen.

\begin{table}[hbt!]
\centering
\caption{Simulationsparameter und Messergebnisse der simulierten Szene.}
\label{table_mixed_sim}
\begin{tabular}{l|l}
 \hline
Partikelradius & 1 cm  \\ \hline
Simulation Substeps & 3      \\ \hline
Substep Iterations & 9            \\ \hline
Anzahl Partikel & 32021            \\ \hline
Davon Fluid Partikel & 16384	\\ \hline
Anzahl Shape-Constraints & 1214		\\ \hline
Anzahl Federn & 15360 		\\ \hline
Self Collision & aktiviert \\ \hline%Nur bei Cloth Assets aktiviert, also für 2562 Partikel		\\ \hline
\textbf{Solve Sync Time (CPU)} & \textbf{0,35 ms}
\end{tabular}
\end{table}

Der Wert \textit{Solve Sync Time} gibt an, wie lange die CPU gewartet hat, bis die FleX Simulation für einen Frame von der GPU abgearbeitet wurde. Das heißt, dass der FleX Solver wahrscheinlich noch schneller gearbeitet hat als angegeben, weil die Datentransferzeiten zwischen CPU und GPU in den Messungen enthalten sind. Dennoch ist dies ein gutes Maß dafür, wie schnell die FleX Simulation ist.

Das Ergebnis von 0,35 ms ist zufriedenstellend. Die GPU hat dadurch noch genügend Zeit, um die Szene mit 60-120 \ac{FPS} bei 1080p zu rendern, solange einfache Rendertechniken und mittelmäßig komplexe Shader eingesetzt werden. 

Was in dieser Arbeit besonders auffiel, ist dass das Rendern von simulierten Flüssigkeiten besonders rechenintensiv werden kann. Wird die Flüssigkeit mit aufwendigen Wasser-Shadern gerendert, so wie in Abbildung \ref{WaterShaderPerformance} rechts, kann dies zu deutlich weniger \ac{FPS} führen. Es ist also ratsam, besonders effiziente Rendertechniken für simulierte Flüssigkeiten einzusetzen.

Da für chirurgische Simulatoren vor allem das Simulieren von Weichkörpern relevant ist, werden die nächsten Tests nur noch mit Soft Bodies simuliert, die aus \textit{clustered Shape-Constraints} bestehen.

%\begin{table}[]
%\centering
%\caption{Neue Objektarten wurden Schrittweise hinzugefügt}
%\label{table_mixed_res}
%\begin{tabular}{l|rrrr|l}
%\textbf{Hinzugefügt} & \multicolumn{1}{l}{\textbf{Partikelanzahl}} & \multicolumn{1}{l}{\textbf{Shape Constraints}} & \multicolumn{1}{l}{\textbf{Fluid Partikel}} & \multicolumn{1}{l|}{\textbf{Distance Constraints}} & \textbf{Solve Sync Time} \\ \hline
%Nur Soft Bodies      & 12548                                       & 1209                                           & 0                                           & 0                                                  & 0.16                     \\
%+ Flüssigkeit        & 28932                                       & 1209                                           & 16384                                       & 0                                                  & 0.17                     \\
%+ Stoff              & 31494                                       & 1209                                           & 16384                                       & 15360                                              & 0.31                     \\
%+ Rigid Bodies       & 32021                                       & 1214                                           & 16381                                       & 15360                                              & 0.35                    
%\end{tabular}
%\end{table}
\clearpage
\subsubsection{Simulationsparametrisierung}

Für alle Messungen, die in diesem Kapitel präsentiert werden, wurden die Einstellungen aus Tabelle \ref{table_soft_param} verwendet.

\begin{table}[hbt!]
\centering
\caption{Verwendete Parameter für diesen Versuch}
\label{table_soft_param}
\begin{tabular}{l|l}
 \hline
Partikelradius & 1 cm  \\ \hline
Simulation Substeps & variierend (2-3)      \\ \hline
Substep Iterations & variierend (3-9)            \\ \hline
Simulierte Objektarten &  nur Soft Bodies mit Shape Constraints\\ \hline
Self Collision & deaktiviert \\ \hline%Nur bei Cloth Assets aktiviert, also für 2562 Partikel		\\ \hline
\end{tabular}
\end{table}

%TODO: Auf Substeps und Substep iterations eingehen!!!

In Abbildung \ref{figure33} ist zu sehen, wie schnell unterschiedlich viele Soft Bodies simuliert werden konnten, bei unterschiedlichen Simulationsparametern. %Für alle Simulationen wurde ein Partikelradius von 1 cm verwendet. 
Mit den Simulationsparametern der blauen Kurve können bereits erste optisch realistisch wirkende Simulationen ausgeführt werden, wobei sich damit keine akkuraten Gewebeverformungen simulieren lassen. Wie man sieht, steigt die blaue Kurve nur sehr schwach an (Steigung ca. $10^{-5}$). Selbst wenn mehr als 100.000 Partikel (26 Soft Bodies, 10478 Cluster) simuliert werden, benötigt der FleX-Solver nur 1,12 Millisekunden für einen Simulationsschritt. Die Szene konnte daher immer mit 120 FPS gerendert werden.
%\clearpage
\bild{figure33}{14cm}{Die blaue Kurve zeigt die Simulationsgeschwindigkeit von FleX bei einfachen (kleinen) Simulationsparametern. Die rote Kurve zeigt die Simulationsgeschwindigkeit von FleX bei einer akkurateren Simulation mit größeren Simulationsparametern.}

Wird die Simulation allerdings akkurater ausgeführt, mit den Parametern der roten Kurve (gut für akkurate Gewebesimulationen), wird die Lösungsgeschwindigkeit ab einer kritischen Partikelanzahl deutlich langsamer. In Abbildung \ref{figure33} erkennt man sehr deutlich, wie die rote Kurve ab einer Partikalanzahl von etwa 60.000 nach oben schießt. Davor läuft die Simulation nahezu genauso schnell wie zuvor, bei geringeren Simulationsparametern.
%, mit einer Steigung von ca. 0,28.

Abbildung \ref{figure4} zeigt welche Auswirkung die steigende Simulationszeit auf die Renderfrequenz hatten. Hierzu sei gesagt, dass in dieser Szene mittelmäßig komplexe Shader und Grafikeinstellungen eingesetzt wurden.

\bild{figure4}{10cm}{Blau: Slove Sync Sime bei der akkuraten Simulation. Die rote Kurve zeigt, wie die FPS mit steigender Partikelanzahl abgefallen sind.}

Das plötzliche Ansteigen der Simulationszeiten bei der akkurateren Simulation weist auf eine Hardwaregrenze hin, bei der nicht mehr genügend Hardwareressourcen zur Verfügung stehen, um weiterhin effizient arbeiten zu können. Der Grafikspeicher der Grafikkarte wurde allerdings niemals voll ausgenutzt (maximal 4 von 8 GB wurden belegt). 
Die Grenze von 60.000 Partikeln konnte im Anwendungsteil dieser Arbeit nie überschritten werden. Im finalen chirurgischen Simulator mussten lediglich 7.000-20.000 Partikel simuliert werden. Dadurch konnten problemlos zusätzliche Simulationen ausgeführt und komplexere Rendertechniken eingesetzt werden, ohne negativen Einfluss auf die FPS.



% ------ HIER EVT FIGURE 6 präsentieren, sprengt aber evt den Rahmen!!! -----------------------

%Als nächstes wurde eine Szene mit konstanter Partikelanzahl simuliert, während einzelne Simulationsparameter stufenweise angehoben wurden.
%figure6


\subsubsection{Cluster Konfigurationen}


%Zeige: Viele kleine Cluster vs wenig große Cluster -> Viele kleine Cluster ist deutlich performanter als wenige große Cluster. WIESO? -> FleX löst Constraints parallelisiert. Die Lösung eines einzelnen Constraints ist eine sequenzielle Operation. Je größer ein Cluster, desto mehr Partikel beinhaltet. Je mehr Partikel in einem Cluster, desto rechenintensiver ist die Lösung des Shape Constraints!

%In diesem Kapitel wird untersucht welchen Einfluss unterschiedliche Cluster Konfigurationen auf die Simulationsperformance von Soft Bodies, die aus Shape-Constraints bestehen, haben. 

In diesem Kapitel wird untersucht, welchen Einfluss unterschiedliche Cluster Konfigurationen auf die Performance von Shape-Constraint-basierten Soft Body Simulationen haben. 

Hierfür wurden zwei sehr ähnliche FleX Assets modelliert, die sich nur durch ihre Cluster Konfigurationen unterscheiden (siehe Abbildung \ref{Clusters} und Tabelle \ref{table_clusters}). Der wesentliche Unterschied zwischen den beiden Assets ist, dass Asset A aus vielen kleinen Clustern besteht, während Asset B aus wenigen großen Clustern besteht. 

\bild{Clusters}{14cm}{Dasselbe FleX-Asset mit zwei unterschiedlicher Cluster Anordnungen. Die Cluster werden als Koordinatensysteme visualisiert. Das linke Asset hat viele kleine Cluster, während das Rechte wenige große Cluster besitzt. Mit beiden Assets können optisch realistisch wirkende Deformationen simuliert werden, allerdings gibt es große Performance-Unterschiede.}

Die erste intuititive Vermutung wäre nun, dass Asset B performanter simuliert werden kann, weil hier weniger Shape-Constraints gelöst werden müssen. Aus den Messergebnissen (siehe Tabelle \ref{table_clusters}) geht allerdings das Gegenteil hervor: Bei Asset B wurde eine 70-mal größere Solve Sync Time gemessen als bei Asset A. 
%Dadurch konnte Asset B nurnoch mit 46 FPS gerendert werden, während Asset A bei 120 FPS gemessen wurden.

%\begin{table}[]
%\centering
%\caption{Verwendete Parameter für diesen Versuch}
%\label{table_soft_assets}
%\begin{tabular}{l|l}
% \hline
%Partikelradius & 1 cm  \\ \hline
%Simulation Substeps & 3     \\ \hline
%Substep Iterations & 9            \\ \hline
%Simulierte Objektarten &  nur Soft Bodies mit Shape Constraints\\ \hline
%Self Collision & deaktiviert \\ \hline%Nur bei Cloth Assets aktiviert, also für 2562 Partikel		\\ \hline
%\end{tabular}
%\end{table}

\begin{table}[hbt!]
\centering
\caption{Messergebnisse des Cluster Versuchs. Assets mit mehr und dafür kleineren Clustern performen deutlich besser.}
\label{table_clusters}
\begin{tabular}{lcc}
                                         & \textbf{Asset A}      & \textbf{Asset B}      \\ \hline
\multicolumn{1}{l|}{Partikelradius}      & 1                     & 1                     \\
\multicolumn{1}{l|}{Simulation Substeps} & 3                     & 3                     \\
\multicolumn{1}{l|}{Substep Iterations}  & 9 					 & 9 				     \\
\multicolumn{1}{l|}{Partikelanzahl}      & 23815                 & 23815                 \\
\multicolumn{1}{l|}{Clusterradius}       & 6                     & 20                    \\
\multicolumn{1}{l|}{Clusteranzahl}       & 133                   & 18                    \\
\multicolumn{1}{l|}{Partikel / Cluster}  & 179                   & 1323                  \\ \hline
\multicolumn{1}{l|}{Solve Sync Time}     & 0,24 ms               & 16,83 ms              \\
\multicolumn{1}{l|}{FPS}                 & 120                   & 46   
\end{tabular}
\end{table}

Die GPU kommt also offenbar bei der Verarbeitung von wenigen großen Clustern schnell an ihre Grenzen.

Hierfür gibt es zwei Ursachen: Zum Einen ist die Verarbeitung der Shape Constraints laut \cite{FlexD3D} aufwändig, weil hierbei die Momentenmatrizen der Cluster polar zerlegt werden müssen. Das ist besonders rechenintensiv.

Zum anderen verarbeitet NVIDIA FleX Shape Constraints nach der sogenannten Constraint orientierten Methode. Um die Simulation zu verarbeiten, wird in \newline \cite{UPP} zwischen zwei unterschiedlichen Verarbeitungsansätzen unterschieden: Partikel-orientiert und Constraint-orientiert. 

Beim Partikel-orientierten Ansatz werden alle Partikel parallel verarbeitet und für jedes Partikel werden die Constraints, die das Partikel beeinflussen, sequenziell verarbeitet. 

Beim Constraint-orientierten Ansatz werden alle Constraints parallel verarbeitet und alle Partikel, die ein Constraint beeinflusst, werden sequenziell verarbeitet, beziehungsweise verschoben. 

NVIDIA FleX setzt beide Verarbeitungsansätze ein, je nachdem um welche Art von Constraint es sich handelt. Shape-Constraints werden mit dem Constraint-orientierten Ansatz verarbeitet \cite{FlexD3D}. Weil für jedes Cluster ein entsprechendes Shape-Constraint gelöst wird, kann man sich also vorstellen, dass alle Cluster parallel verarbeitet werden, während alle Partikel, die sich innerhalb eines Clusters befinden, sequenziell verarbeitet werden. Je größer ein Cluster, desto mehr Partikel beinhält es und desto mehr Partikel müssen sequenziell verarbeitet werden. 

Deswegen ist es besser, viele kleine Cluster zu verwenden. So wird mehr Arbeit parallelisiert (viele Cluster) und es muss weniger sequenziell verarbeitet werden (wenig Partikel pro Cluster). Außerdem bringt dies den zusätzlichen Vorteil, dass sich ein Soft Body, das aus vielen Clustern besteht, vielfältiger verformen lässt.

Bei der Modellierung von größeren Rigid Bodies kann es allerdings schwierig sein, die Partikelanzahl pro Cluster klein zu halten. Jedes Rigid Body besteht nämlich aus nur einem Cluster. Deswegen kommt die GPU schnell an ihre Grenzen, wenn größere Rigid Bodies simuliert werden, die aus vielen Partikeln bestehen. Daher sollten Rigid Bodies mit so wenig Partikeln wie möglich modelliert werden, wenn möglich. %Beispielsweise indem der Partikelradius auf einen großen Wert gesetzt wird oder indem die Partikel nur auf die Oberfläche des Polygonmodells gesetzt werden. 
Lässt sich die Partikelanzahl pro Rigid Body nicht ausreichend reduzieren, kann es sich lohnen, alle Rigid Bodies mit einer traditionellen Rigid Body Physics Engine zu simulieren, beispielsweise mit NVIDIA PhysX. Solche Engines können auch größere Rigid Bodies effizient simulieren. Wie NVIDIA PhysX und NVIDIA FleX gleichzeitig verwendet werden können, wird in Kapitel \ref{physXFleX} behandelt.

\subsubsection{Self Collision und Partikelradius}

%Fazit: Wenig Partikelüberschneidung bei Self Collision. Wichtig zu erwähnen: Führt auch zu instabilen, zitternden Zuständen! Außerdem: Je mehr Kollisionen/Kontackte existieren, desto mehr Collisions und friction constraints müssen gelöst werden.

In diesem Kapitel wird untersucht, wie stark sich unterschiedliche Partikelradi auf die Simulationsperformance von Soft Bodies auswirkt. Hierfür wurde eine Szene mit den Einstellungen aus Tabelle \ref{table_self_collision} simuliert. 

\begin{table}[hbt!]
\centering
\caption{Verwendete Parameter für diesen Versuch}
\label{table_self_collision}
\begin{tabular}{l|l}
 \hline
Partikelradius & variierend  \\ \hline
Simulation Substeps & 3      \\ \hline
Substep Iterations & 9            \\ \hline
Simulierte Objekte &  10 Soft Bodies \\ \hline
Self Collision & variierend \\ \hline%Nur bei Cloth Assets aktiviert, also für 2562 Partikel		\\ \hline
Ignore Rest Collision & True\\ \hline
Shape-Constraints & 850            \\ \hline
Partikelanzahl & 26010            \\ \hline
\end{tabular}
\end{table}

Dabei wurde der Partikelradius stufenweise erhöht. 

\bild{RadiusOctupus}{14cm}{Je größer der Partikelradius, desto mehr Selbstkollisionen kommen in einem einzelnen Asset vor. Sind die Partikel zu klein, entstehen Lücken zwischen den Partikeln (links). Das sollte vermieden werden, sonst können Objekte durch solche Lücken hindurch fallen.}

Dieser Test wurde einmal mit Assets durchgeführt, bei denen \textit{Self Collision} aktiviert war und einmal mit Assets, bei denen \textit{Self Collision} deaktiviert war. Ist \textit{Self Collision} aktiviert, können die Partikel, die sich innerhalb eines Assets befinden, miteinander kollidieren, wobei bei der Ruhestellung keine Kollisionen stattfinden können. Ist Self Collision deaktiviert, können die Partikel desselben FleX Objekts nicht miteinander kollidieren. Nur mit den Partikeln von anderen Objektinstanzen.

Die Messergebnisse sind in Abbildung \ref{figure5} zu sehen.

\bild{figure5}{10cm}{Nur wenn Self Collision aktiviert ist, können zu große Partikelradi zu einer schlechteren Performance führen.}

Wie man sieht, entsteht ab einem Partikelradius von etwas 15 cm ein starker Performanceverlust, allerdings nur, wenn Self Collision aktiviert ist. Das liegt vermutlich daran, dass sich mit wachsendem Partikelradius benachbarte Partikel immer stärker überlappen (siehe Abbildung \ref{RadiusOctupus}). Dabei entstehen immer mehr Partikelkollisionen. Ab einem bestimmten Punkt stehen nicht mehr genügend Hardwareressourcen zur Verfügung, um die vielen Kollisions-Constraints effizient zu lösen.

Solange sich die Partikel nicht zu stark überlappen (so wie in Abbildung \ref{RadiusOctupus}, Mitte), kommt FleX allerdings gut mit der Kollisionsverarbeitung zurecht, auch wenn Self Collision aktiviert ist.

\subsection{Zusammengefasste Erkenntnisse}
\label{subsec_perf_zusa}

Mithilfe der durchgeführten Tests konnten wichtige Erkenntnisse über die Performance von FleX Simulationen gewonnen werden. Zusammengefasst sind dies die wichtigsten Punkte, auf die zu achten sind, um effiziente FleX Simulationen zu erhalten:

    \begin{itemize}
    
      \item \textbf{Hardwaregrenzen:} \\Die Performance von FleX Simulationen skaliert nur flach mit der Partikel- und Constraint-Anzahl, solange alle Cluster klein sind. Je nach Hardware gibt es jedoch eine Grenze, bei der die Performance deutlich stärker mit wachsender Partikel- und Constraint-Anzahl abfällt. Je nach Simulationsparametrisierung tritt diese Grenze früher oder später auf. Die Hardwaregrenze sollte bekannt sein und nicht überschritten werden, um noch genügend Ressourcen für zusätzliche Simulationen und fotorealistisches Rendern zu haben.
      
      \item \textbf{Partikeldichte:} \\ Bei der Modellierung von FleX Assets sollten sich benachbarte Partikel berühren, aber nicht zu stark überlappen. Vor allem wenn Self Collision aktiviert ist, können zu viele Kollisions-Constraints entstehen, wodurch die Performance negativ beeinflusst wird.
      
      \item \textbf{Geringe Partikelanzahl pro Cluster:} \\ Bei der Modellierung von Soft Assets sollten viele kleine Cluster verwendet werden. Das heißt, die Cluster sollten nur wenige Partikel, etwa maximal 200, umfassen. Je mehr Partikel die Cluster verarbeiten müssen, desto schneller kommt die Hardware an ihre Grenzen. Weil ein Rigid Body immer aus einem Cluster besteht, sollten alle Rigid Bodies aus so wenig Partikeln wie möglich bestehen.
      

      
      \item \textbf{Simulationsparametrisierung:} \\ In einigen Arbeiten und auch in dieser Arbeit erwiesen sich die Simulationsparameter aus Tabelle \ref{table_perf_zusa} als geeignet, um performante und akkurate Simulationen zu erhalten:
      
    \end{itemize}

\begin{table}[hbt!]
\centering
\caption{Empfohlene Simulationseinstellungen}
\label{table_perf_zusa}
\begin{tabular}{l|l}
 \hline
Partikelradius & \begin{tabular}[c]{@{}l@{}}So dass die Form der simulierten Körper\\ gut von Partikeln repräsentiert werden kann.\end{tabular}  \\ \hline
Simulation Substeps & 3      \\ \hline
Substep Iterations & 9            \\ \hline
Self Collision & Je nach Bedarf. \\ \hline%Nur bei Cloth Assets aktiviert, also für 2562 Partikel		\\ \hline
Ignore Rest Collision & True\\ \hline
Maximale Partikelanzahl & Bei GTX 1070: Maximal 60.000\\ \hline
Relaxation type& Local\\ \hline
Fixed Time Step& True\\ \hline
\end{tabular}
\end{table}


\clearpage
\section{Funktionalitäten}
\label{sec_functional}
In der Einleitung von \cite{UE4FlexDoc} wird bereits darauf hingewiesen, dass die Unreal Integration noch nicht weit genug entwickelt ist, um Gameplay-beeinflussende Physikszenarios zu erstellen. In diesem Kapitel wird genauer untersucht, welche   Funktionalitäten  die Unreal Integration bereits bietet und welche noch fehlen. 
%Hierfür gehen wir nach und nach alle funktionalen Anforderungen aus Kapitel \textbf{XXX} durch und kontrollieren, ob Anforderungen fehlen oder nur teilweise vorhanden sind.

%\subsection{Synchronisation mit Flex Simulation}

%... kommt evt raus, ist ja nicht so wichtig....

\subsection{Transformationsmanipulationen}


Game Engines verfügen in der Regel über eine umfangreiche Bibliothek für Transformationsberechnungen. So können Punkte oder Richtungen mit homogenen Transformationen manipuliert und die Transformationsmatrizen von unterschiedlichen 3D-Objekten miteinander verrechnet werden.
Dies ist eine wichtige Grundlage, um 3D-Objekte bewegen, rotieren und skalieren zu können. 

Die Unreal Engine bietet hierfür den Datentyp \textit{FTransform} und die Bibliothek \textit{TransformCalculus.h} an. Jedes 3D-Objekt besitzt eine eigene Transformation, die zur Laufzeit manipuliert werden kann. Häufig werden die beiden Funktionen \newline \textit{SetWorldTransform($T_{world}$)} und \textit{SetRelativeTransform($T_{relative}$)} genutzt, um 3D-Objekte relativ zur Welt oder relativ zum derzeitigen Eltern-Objekt zu bewegen.

Bei der Unreal Engine mit Flex Integration wurde das Transformationsframework nur teilweise für FleX-Objekte angepasst. Obwohl alle \textit{FlexComponents} eine eigene Transformation vom Typ \textit{FTransform} besitzen, kann diese weder relativ noch absolut gesetzt werden. Die beiden Funktionen \textit{SetWorldTransform()} und \textit{SetRelativeTransform()} haben keine Auswirkung auf FlexComponents.

Diese wichtige Funktionalität muss also noch hinzugefügt werden. Dabei wird das bereits bestehende Transformationsframework der Unreal Engine benutzt.

\subsubsection{\textit{SetWorldTransform}}

Um die Welt-Transformation eines Flex-Objekts zu ändern, müssen die Positionen all seiner Partikel mithilfe einer homogenen Transformationsmatrix, $T_{world}$, transformiert werden. Über die NVIDIA FleX Bibliothek können die lokalen Positionen, $\vec{p}_{local}$, der Partikel eines Flex-Assets ermittelt werden. Das sind die lokalen Partikelpositionen des Flex-Assets, während seiner Ruheposition. Die neue Position, $\vec{p}_{new}$ eines Partikels ergibt sich dann mit Gleichung \ref{form_worldtransf}.

\begin{equation}
\vec{p}_{new} = T_{world} \cdot \vec{p}_{local}
\label{form_worldtransf}
\end{equation}

Gleichung \ref{form_worldtransf} ist im Grunde nichts anderes als die Transformation eines Punktes vom Objektkoordinatensystem ins Weltkoordinatensystem aus der klassischen Computergraphik. 

Im Rahmen dieser Arbeit hat es sich als hilfreich erwiesen, die Geschwindigkeiten der transformierten Partikel auf null zu setzen, um das versetzte Flex-Objekt sofort in einen ruhigen Simulationszustand zu versetzten. 

Soll ein Flex-Objekt  über einen längeren Zeitraum entlang einer Trajektorie bewegt werden, ist es sinnvoll, die bewegten Partikel von der Flex-Simulation auszuschließen, indem ihre invertierten Massen auf null gesetzt werden. Dadurch sind die bewegten Partikel unendlich schwer und können nicht mehr von der Flex-Simulation bewegt werden \cite{PBD}. Wenn bewegte Partikel von der Simulation ausgeschlossen werden, erscheinen sie allerdings starr. Um bewegte Weichkörper weiterhin elastisch erscheinen zu lassen, wird nur eine Teilmenge der Partikel bewegt. Dann erscheint nur ein Teil des bewegten Körpers starr und der Rest bleibt elastisch und folgt den bewegten Partikeln.

\subsubsection{\textit{SetRelativeTransform}}
\label{subsubsec_reltrans}
%Um ein Partikel relativ zu einer Parent-Szene zu bewegen (\textit{SetRelativeTransform()}), wird zunächst die allgemeine Situation aus Abbildung \ref{Transform1} betrachtet:

% eventuel streichen!
%\bild{Transform1}{10cm}{Um ein Partikel mit $T_{rel}$ relativ zu \textit{ParentScene} zu transformieren, muss das Matrixprodukt aus der Welt-Transformation von \textit{ParentScene}, $T_{parent}$, und $T_{rel}$ gebildet werden.}

Gegeben sei die Welt-Transformation einer Eltern-Szene, $T_{parent}$ und die gewünschte Transformation, $T_{rel}$, die das Partikel relativ zur Eltern-Szene einnehmen soll.
%(siehe Abbildung \ref{Transform1}). 
Dann kann mit Gleichung \ref{form_reltransf} die neue Weltposition, $\vec{p}_{new}$, des Partikels berechnet werden:

\begin{equation}
\vec{p}_{new} = T_{parent} \cdot T_{rel} \cdot \vec{p}_{local}
\label{form_reltransf}
\end{equation}

Gleichung \ref{form_worldtransf} und \ref{form_reltransf} sind einfache, homogene Matrix-Vektor-Multiplikationen und sollten daher effizient von heutiger Hardware gelöst werden können. Außerdem kann die Berechnung der Partikelpositionen parallelisiert werden.

\subsection{Befestigung von FleX-Objekten}
\label{subsec_attach}

%Für gewöhnlich arbeiten Game Engines mit sogenannten Scene Trees um 3D-Objekte hierarchisch zu organisieren (REFERENZ). Dabei repräsentiert jedes 3D-Objekt einen Knoten im Scene Tree. 3D-Objekte werden aneinander befestigt, so das sich eine Baumstruktur bildet, in der die Kinderobjekte an ihre Elternobjekte befestigt sind. Ändert sich die Transformation eines Elternknotens, bewegen sich all seine Kinderobjekte mit. 
%
%\bild{SceneTree}{8cm}{Alle Szenen und 3D-Objekte werden in einem Szene Tree organisiert}
%
%Die Struktur eines Scene Trees lässt sich zur Laufzeit eines Spiels ändern, indem 3D-Objekte per Funktionsaufruf von ihren Eltern gelöst (\textit{Detach()}) oder an andere 3D-Objekte befestigt (\textit{AttachTo(NewParent)}) werden. Diese beiden Funktionen werden in Computerspielen häufig benötigt, beispielsweise damit der Spieler 3D-Objekte aufheben (\textit{AttachTo(Hand)}) und wieder loslassen (\textit{Detach()}) kann.

\subsubsection{Befestigung von Flex-Partikeln an Szenen}

In der Unreal Integration können mehrere Flex-Partikel an eine Szene befestigt werden. Hierfür stehen allerdings nur Funktionen zur Verfügung, bei denen alle Partikel, die sich innerhalb einer Kugel befinden, befestigt werden. Auf Grundlage dieser Funktion wurden daher weitere Funktionen angelegt, um einzelne Partikel via Partikelindex an Szenen befestigen zu können.

Um ein Flex-Partikel an eine Szene zu befestigen, muss dessen inverse Masse, $w$, auf 0 gesetzt werden. Dann hat das Partikel eine unendlich große Masse und kann nicht mehr von der FleX-Simulation bewegt werden. 
Gegeben sei außerdem die Position, $\vec{p}_{local}$, die das Partikel relativ zur Szene beibehalten soll. Immer wenn sich die Szene bewegt, wird die Weltposition, $\vec{p}_{new}$, des befestigten Partikels so verändert, dass die relative Position $\vec{p}_{local}$ unverändert bleibt.
$\vec{p}_{new}$ kann mithilfe von Gleichung \ref{form_reltransf} aus Kapitel \ref{subsubsec_reltrans} berechnet werden.

%Des weiteren musste eine Funktion zum Lösen von befestigten Partikel angelegt werden (\textit{DetachAllFlexParticles()}), weil auch solch eine Funktionalität fehlte.


%%% EVT noch genauer beschreiben wie genau das Befestigen funktioniert...

\subsubsection{Befestigung von Szenen an Szenen}

Erfreulicherweise unterstützt die verwendete Unreal Integration bereits das Befestigen von Szenen an ein Flex-Objekt. Allerdings drehen sich Szenen nicht mit Soft Bodies mit. Nur die Position der befestigten Szenen wird angepasst, während die Rotation unverändert bleibt. 

%Diese Problemstellung wurde im Rahmen dieser Arbeit nicht weiter verfolgt, weil hier keine Szenen an Soft Bodies befestigt werden mussten. Es wurden nur Szenen an Rigid Bodies befestigt, was problemlos funktionierte.

Diese Problemstellung wurde in dieser Arbeit aus Zeitgründen nicht weiter verfolgt, sie könnte für nachfolgende Arbeiten interessant sein.

\subsubsection{Befestigung von Flex-Partikel an Flex-Partikel}

Grundsätzlich kann ein FleX-Partikel an einem anderes FleX-Partikel befestigt werden, indem die inverse Masse des befestigten Partikels auf 0 gesetzt wird und dessen Position stets auf die Position des anderen Partikels gesetzt wird. 
Dann wird allerdings keine Physik mehr für das befestigte Partikel simuliert. %satz bisl doof...
Im Rahmen dieser Arbeit war dies allerdings nicht erwünscht. In vielen Fällen sollten zwar zwei Flex-Objekte aneinander befestigt werden, allerdings sollten sich beide physikalisch realistisch verhalten. 

Für genau diesen Anwendungsfall bietet NVIDIA FleX bereits eine Lösung an. Mithilfe von NVIDIA FleX können physikalisch realistisch wirkende Verbindungen zwischen mehreren Flex-Objekten hergestellt werden. Solche Verbindungen heißen \textit{Soft Joints} (siehe \textit{NvFlexExtSoftJoint} aus \textit{NvFlexExt.h}).

\subsection{Kollisionserkennung}


Eine wichtige Funktionalität, die die meisten Game Engines bieten, ist ein Framework für die Erkennung von Kollisionen zwischen 3d-Modellen. 
Ein Kollisionserkennungsframework wird für Computerspiele häufig benötigt, beispielsweise um zu erkennen, ob der Spieler in tödliche Lava fällt, ob der Spieler ein Zielgebiet erreicht hat oder ob eine geschwungene Klinge ein Ziel getroffen hat. Die Programmierschnittstelle für solche Kollisionserkennungen existiert meist in der Form von \textit{Event-Bindings}. 

So verhält es sich auch bei der Unreal Engine 4. Hier bieten 3D-Objekte unterschiedliche Kollisions-Events an. Zum Beispiel wird ein \textit{BeginOverlap}-Event ausgelöst, wenn sich zwei 3D-Modelle berühren oder überlappen. Beliebige Klassen können auf das \textit{BeginOverlap}-Event eines 3D-Modells hören (Event-Binding), damit Sie immer darüber informiert werden, wenn das 3D-Modell etwas berührt hat.

Für chirurgische Simulatoren wird ein Kollisionserkennungsframework benötigt, beispielsweise um zu erkennen, ob ein Skalpell ein empfindliches Organ berührt oder ob die Spitze einer Kanüle eine Blutader getroffen hat.

In der Unreal Integration können Flex-Objekte zwar grundsätzlich miteinander kollidieren, allerdings gibt es nicht die Möglichkeit, auf das Kollisions-Event eines Flex-Objekts zu hören (\textit{Event-Binding}) oder eine Kollision zu erkennen.

Um diese fehlende Funktionalität zu ergänzen, wird das bereits bestehende Kollisionserkennungsframework der Unreal Engine 4 genutzt, welches auf NVIDIA PhysX aufbaut \cite{ue4physics}. Dieses Kollisionserkennungsframework ist nämlich bereits sehr umfangreich und effizient.

%In \cite{UE4FlexDoc} wird empfohlen, fehlende Kollisionserkennungen mit einer traditionellen Rigid Body Physics Engine zu realisieren. Weil die Unreal Engine 4 bereits solch eine Physics Engine besitzt (\cite{ue4physics}), bietet es sich an, die fehlende Kollisionserkennung mit dieser zu realisieren.

Hier kommt es uns zugute, dass bei der Unreal Integration 3D-Modelle an Flex-Objekte befestigt werden können (siehe Kapitel \ref{subsec_attach}). So können auch PhysX-Kollisionskörper für Flex-Objekte modelliert und an diese befestigt werden, so wie in Abbildung \ref{Scalpel1} zu sehen.

\bild{Scalpel1}{14cm}{Ein passender PhysX-Kollisionskörper (grünes Polygonmodel) wird an ein Flex-Asset befestigt, um so die umfangreiche Kollisionserkennung der Unreal Engine nutzen zu können. Hier reicht ein ungenauer Kollisionskörper aus, weil in diesem Fall nur erkannt werden soll, ob eine virtuelle Hand nach dem Skalpell greift.}

Die separaten Kollisionskörper werden dann so konfiguriert, dass sie keine Physik simulieren. Nur die Überlappungserkennung wird genutzt, um Überlappungen und Berührungen zu erkennen.

An ein Flex-Asset können auch mehrere Kollisionskörper befestigt werden, um zwischen unterschiedlichen Berührungen unterscheiden zu können. 
%Beispielsweise um zu erkennen ob die Klinge oder der Griff eines Skalpells etwas berührt hat.

%Hier evt nochmal ein Bild mit komplexeren Kollisionskörpern für Klinge, Griff, Kanülenspitze, Hauptader usw...

%  - FlexCollisionComponent kann nur emitter-partikel zählen oder löschen
  
  %%%%%%%%%%%%%%%%%%%% Flex-PhysX Interaktion %%%%%%%%%%%%%%%%%%%%%%%%%%%%%%%%%%%%%%%%%%%%%%%%
\subsection{Interaktion zwischen FleX- und PhysX-Objekte}
\label{physXFleX}

%Flex-Objekte sehen Physics-Kollisionen als World-Kollision, also als ein statisches Objekt das nicht weg schiebbar ist. Dadurch überrollen Physics-Objekte andere Flex-Objekte, so als ob ihre Masse unendlich groß wäre.
 
Die Unreal Engine 4 besitzt bereits eine traditionelle Physics Engine für Rigid Bodies, die auf NVIDIA PhysX basiert \cite{ue4physics}. Mithilfe der Unreal Integration können also  FleX-Simulationen und PhysX-Simulationen gleichzeitig ausgeführt werden. Nun stellt sich die Frage, ob und wie elastische  FleX-Objekte mit PhysX-Objekten interagieren können.

Bei der verwendeten Unreal Integration können PhysX-Objekte grundsätzlich mit FleX-Objekten kollidieren. Allerdings üben Flex-Objekte keine Kräfte auf PhysX-Objekte aus. FleX-Objekte sehen PhysX-Objekte als statische Kollisionskörper, die nicht verschoben werden können. Dies führt zu Situationen, wie in Abbildung \ref{PhsXInteraction} links abgebildet: Kollidiert ein PhysX-Objekt (weiße Kugel) mit einem Flex-Objekt (blaue Kugel), wird das FleX-Objekt vom PhysX-Objekt ungehindert weggeschoben oder gar überrollt, wie in der Abbildung zu sehen. Weil ein FleX-Objekt keine Reaktionskraft auf das PhysX-Objekt ausüben kann, bewegt sich das PhysX-Objekt ungehindert weiter, als würde überhaupt keine Kollision mit einem anderen Objekt stattfinden.

\bild{PhsXInteraction}{10cm}{Links: Das PhysX-Objekt (weiße Kugel) überrollt ungehindert das elastische FleX-Objekt (blaue Kugel), weil FleX-Objekte keine Kraft auf PhysX-Objekte ausüben. Rechts: Die fehlende Reaktionskraft auf das PhysX-Objekt wird anhand der Verformung des Flex-Objekts ausgerechnet und angewandt. Nun prallt das PhysX-Objekt vom Flex-Objekt ab und verursacht nur eine geringe Verformung.}

Um dieses Problem zu lösen, muss die fehlende Reaktionskraft, die auf das PhysX-Objekt wirkt, berechnet und angewandt werden. 
%Im Rahmen dieser Arbeit wurden hierfür eine vereinfachte Methode entwickelt. 
%Der Resultierende Kraftvektor ergibt sich aus der Verformung, die das PhysX Objekt auf das Flex Objekt
Im Rahmen dieser Arbeit wurde der Vektor der fehlenden Reaktionskraft so berechnet, wie in \cite{VRSim20}, Kapitel 6 (\textit{Haptic Feedback}). Dieser einfache Ansatz basiert auf der folgenden Idee: 

Wenn ein PhysX-Objekt mit einem elastischen FleX-Objekt kollidiert, verursacht das PhysX-Objekt eine Verformung beim FleX-Objekt. Vereinfacht lässt sich sagen: Je stärker die Verformung und je größer die Fläche der Verformung, desto stärker die Reaktionskraft auf das PhysX-Objekt. Für die Flex-Partikel des Flex-Objekts heißt das: Je stärker die Partikel verschoben werden und je mehr Partikel verschoben werden, desto stärker ist die Reaktionskraft. Also müssen nur die Partikelverschiebungsvektoren ermittelt, umgedreht und summiert werden, um einen finalen Kraftvektor zu erhalten, der sich wie gewünscht verhält. 
%Dabei werden die Partikelverschiebungsvektoren mit den Massen der verschobenen Partikel gewichtet. Dadurch werden schwerere FleX Objekte geringer verschoben.

Zusätzlich wird der Kräftevektor mit dem \textit{Stiffness} Parameter des Flex-Objekts gewichtet.
Der Kräftevektor wird dann mithilfe eines einfachen Funktionsaufrufs an die PhysX Simulation übergeben.
In dieser Arbeit konnten mit dieser einfachen Methode optisch realistisch wirkende Ergebnisse erzielt werden.
%Dieser Ansatz ist nur für weiche Soft Bodies geeignet. Handelt es sich um steifere Soft Bodies, könnte noch zusätzlich der Impuls mitberücksichtigt werden. 
%Dank des gegebenen PhysX Frameworks wirkt sich weiterhin auch die Masse des PhysX-Objekts auf die Simulation aus: Je schwerer ein PhysX-Objekt, desto stärker kann es Flex-Objekte eindrücken.

  %%%%%%%%%%%%%%%%%%%% Flex-PhysX Interaktion ENDE %%%%%%%%%%%%%%%%%%%%%%%%%%%%%%%%%%%%%%%%%%%%%%%%

\subsection{Kräfte und Momente}
\label{subsec_forces}

Eine Funktionalität, die einige Game Engines anbieten, ist die Möglichkeit per Funktionsaufruf Kräftevektoren auf einzelne Körper wirken zu lassen. 
%Diese Funktionalität steht bei Game Engines als Funktion in der Form \textit{AddForce(FVector* ForceVector)} zur Verfügung, wie auch bei der Unreal Engine 4 (REEEEF). 

Diese Funktionalität wird meist genutzt, um kräftebasierte Effekte oder Systeme zu programmieren. Beispielsweise können Kräfte auf Körper, die sich innerhalb einer Explosion befinden, angewandt werden, damit sie von der Explosion weggeschleudert werden. 
In \cite{Studienarbeit} wird diese Funktionalität genutzt, um Auftriebskräfte auf schwimmende Objekte wirken zu lassen. 
Außerdem können mit dieser Funktionalität Regelsysteme implementiert werden, die mithilfe von simulierten Kräften Körper balancieren oder zu einer Zielkonfiguration bewegen. In dieser Arbeit wird solch ein kräftebasiertes Regelsystem implementiert und dient als wichtige Grundlage für die Interaktion mit Soft Bodies in virtueller Realität ohne haptisches Feedback (siehe Kapitel \ref{sec_FSystem}).

%Bei der verwendeten Unreal Engine mit Flex Integration gibt es nicht die Möglichkeit, Kräfte auf einzelne Flex-Objekte oder Flex-Partikel wirken zu lassen.

Die NVIDIA FleX Bibliothek bietet zwar die Möglichkeit, Kraftfelder zu erstellen (\textit{NvFlexExtForceField}), allerdings gibt es nicht die Möglichkeit, Kräfte oder Momente gezielt auf einzelne Flex-Objekte oder Partikel wirken zu lassen. %Die Kräftefelder definieren nur einen Bereich, in dem konstante Kräftevektoren im Raum verteilt sind. Dann wirkt auf alle Partikel, die sich innerhalb eines Kräftefeldes befinden, eine Kraft, abhängig von den Partikelpositionen.

%Idee und Ansatz:
Eine einfache Möglichkeit, eine Kraft auf ein Partikel wirken zu lassen, ist die Manipulation der Partikelgeschwindigkeiten, so wie es auch in \cite{PBD} getan wird. Um eine Kraft, die auf ein Partikel wirken soll, zu simulieren, wird eine Geschwindigkeitsänderung, $\Delta\textbf{v}$, zur aktuellen Geschwindigkeit, $\textbf{v}$, des Partikels hinzuaddiert. So kann die Flex-Simulation ungestört weiterlaufen und alle Constraints
%, wie das Kollisions-Constraint, 
wirken sich weiterhin auf die simulierten Partikel aus.

%Methode
\subsubsection{Kräfte}
Gegeben sei eine Kraft, $\textbf{F}$, die auf ein Flex-Objekt, das aus $n$ Partikeln besteht, wirken soll. Dann wird $\textbf{F}$ in $n$ gleichgroße Partikelkräfte, $\textbf{F}_i = \textbf{F} / n$ unterteilt. Jede Partikelkraft, $\textbf{F}_i$, wird dann auf ein Partikel, $i$, mit einer Punktmasse von $m_i$, angewandt.

Weil wir es mit Punktmassen zu tun haben, gilt für die Partikelkraft, $\textbf{F}_i$:

\begin{equation}
\textbf{F}_i = m_i \cdot \textbf{a} = m_i \dfrac{\Delta \textbf{v}_i}{\Delta t}
\label{form_F}
\end{equation}

Also ergibt sich die Geschwindigkeit, $\Delta \textbf{v}_i$, die zur Geschwindigkeit des Partikels addiert werden muss durch:

\begin{equation}
\Delta \textbf{v}_i = \dfrac{\textbf{F}_i}{m_i}\Delta t
\label{form_dv}
\end{equation}

Wobei $\Delta t$ der derzeitige Zeitschritt der Flex-Simulation ist. 

%Um eine Kraft auf ein komplettes Flex-Objekt zu wirken, wird $\Delta \textbf{v_i}$ auf all seine Partikel angewandt. 
%Um die gesamte Kraft, die auf ein Flex-Objekt wirkt, zu berechnen, müssen alle Partikelkräfte, $F_i$ summiert werden.

Mit dieser einfachen Methode können Flex-Objekte realistisch umher geschoben werden. Wird eine Kraft von $(0 \quad 0 \quad 9,81 )$ Newton auf ein Flex-Objekts mit einer Masse von 1 kg angewandt, wirkt es schwerelos (bei $\textbf{g}=(0 \quad 0 \quad -9,81 )$). Das zeigt, dass alle Einheiten korrekt berechnet wurden.

%Die Masse eines Flex-Objekts ergibt sich durch die Summe all seiner $n$ Partikelmassen, beziehungsweise durch $n \cdot m_p$, wenn alle Partikelmassen, $m_p$, gleich groß sind.

\subsubsection{Momente}

Wirkt während eines Zeitschritts ein Moment, $M$, auf ein Flex-Objekt, müssen auf all seine Partikel Winkelgeschwindigkeiten, $\Delta \omega_i$, hinzu addiert werden. Wie bei der Kraft auch, zerlegen wir $M$ in $n$ gleichgroße Partikelmomente, $M_i=M/n$ und betrachten nun kleinere Momente, die auf Punktmassen wirken.

%Zunächst wird die Winkelgeschwindigkeit, $\Delta \omega$, berechnet. 
Die Berechnung von $\Delta \omega_i$ erfolgt wie zuvor nur im rotatorischem Sinne (siehe Gleichung \ref{form_omega}):

\begin{equation}
\begin{split}
M_i = m_i   \dot{\omega}_i = m_i \dfrac{\Delta \omega_i}{\Delta t} \\
\Rightarrow  \Delta \omega_i = \dfrac{M_i}{m_i} \Delta t
\end{split}
\label{form_omega}
\end{equation}

%Um die Rechenkomplexität zu senken, gehen wir davon aus, dass alle Partikel die selbe Masse haben, wodurch alle Partikel mit der selben Winkelgeschwindigkeit ($\Delta \omega_i$) rotiert werden müssen. 

Neben der Größe des Moments, $M_i$, sei auch eine örtlich gebundene Achse definiert, die angibt, in welche Richtung das Drehmoment wirkt. Diese Achse besteht aus einem normalisierten Richtungsvektor, $\textbf{e}$, der sich an der Position, $\textbf{c}$, befindet (siehe Abbildung \ref{moment}). $\textbf{e}$ und $\textbf{c}$ definieren also die örtlich gebundene Achse, um die sich das Flex-Objekt drehen wird, aufgrund des wirkenden Drehmoments. Das bedeutet, alle $n$ Partikel sollen sich nun zusätzlich mit einer Winkelgeschwindigkeit von $\Delta \omega_i$ um diese Drehmomentachse drehen.

Nun muss berechnet werden, in welche vektorielle Richtung, $\Delta \textbf{v}_i$ sich ein Partikel $i$ bewegen muss, damit es sich mit der gewünschten Winkelgeschwindigkeit, $\Delta \omega_i$, um die Drehmomentachse dreht. Dabei habe das Partikel die Position $\textbf{p}_i$. Abbildung \ref{moment} visualisiert diese Problemstellung.

\bild{moment}{6cm}{Das Partikel $i$ muss sich zusätzlich in die Richtung $\textbf{d}$ bewegen, damit es damit beginnt, sich um die Drehmomentachse zu drehen. Die Drehmomentachse habe die Position $\textbf{c}$ und Ausrichtung $\textbf{e}$.}


$\Delta \textbf{v}_i$ besteht aus einer Geschwindigkeit, $s_i=|\Delta \textbf{v}_i|$ und einer Bewegungsrichtung, $\textbf{d}$, mit $|\textbf{d}|=1$. Das heißt, es gilt $\Delta \textbf{v}_i = s_i \cdot \textbf{d}$.

$s$ ist die Umfangsgeschwindigkeit des Partikels um die Drehmomentachse und ist abhängig von der Winkelgeschwindigkeit $\Delta \omega_i$ (siehe Gleichung \ref{form_speed}).

\begin{equation}
s_i = \Delta \omega_i \cdot r
\label{form_speed}
\end{equation}

$r$ ist dabei die kürzeste Entfernung, die das Partikel zur Drehmomentachse hat.
Weil $|\textbf{e}|=1$, kann $r$ aus dem folgenden Skalarprodukt berechnet werden (siehe Gleichung \ref{form_r}):

\begin{equation}
\begin{split}
r = ||\textbf{e} \cdot \textbf{u}_i|| \\
\text{mit } \textbf{u}_i=\textbf{p}_i-\textbf{c}
\end{split}
\label{form_r}
\end{equation}

Der Richtungsvektor, $\textbf{d}$, muss orthogonal zu $\textbf{u}_i$ und $\textbf{e}$ sein und eine Länge von 1 haben. Also kann $\textbf{d}$ mit Gleich \ref{form_dir} berechnet werden:

\begin{equation}
\textbf{d} = unit(\textbf{u}_i \times \textbf{e} )
\label{form_dir}
\end{equation}

Wobei mit $unit(...)$ der Einheitsvektor aus einem Vektor gemeint ist.

Schlussendlich ergibt sich die Geschwindigkeit, die zur Partikelgeschiwindigkeit hinzuaddiert werden muss, durch Formel    \ref{form_momentfinal}:

\begin{equation}
\Delta \textbf{v}_i =  \dfrac{M_i}{m_i} \Delta t ||\textbf{e} \cdot \textbf{u}|| \cdot unit(\textbf{u}_i \times \textbf{e} )
\label{form_momentfinal}
\end{equation}

Auch mit diesem Ansatz konnten optisch realistisch wirkende Ergebnisse erzielt werden.

%TODO: Kurz erläutern wie das funktioniert und was Probleme sind (Partikel wollen sich immer bisl weg von Achse bewegen, wodurch innere Spannungen auftreten, die con constraints aber gelöst werden...

%Mit diesem Ansatz können FleX-Objekte mithilfe von Drehmomenten in Rotationen versetzt werden. Weil hierbei die Partikel pro Zeitschritt, $\Delta t$, auch einen kleinen Schritt von der Rotationsachse weg bewegt werden, neigen besonders weiche Soft Bodies dazu, sich geringfügig auszudehnen. Dieser Effekt kann mit kleineren Zeitschritten verringert werden.

\chapter{Implementation von Chirurgischen Szenarien}
\label{chap_Szenarien}

%Welche Use Cases habe ich implementiert und wie? Und wie klappte das?
In Kapitel \ref{chap_Flex_Engine} wurde die Unreal Engine 4 mit Flex Integration um einige Funktionalitäten erweitert, damit sie sich besser für die Implementation von chirurgischen Simulatoren eignet. 

Um nun die Eignung der erweiterten Unreal Integration für chirurgische Simulatoren zu testen, werden in diesem Kapitel einige Szenarien implementiert, die für chirurgische Simulatoren interessant sein könnten.

%In Diesem Kapitel werden Testszenarien implementiert, die für chirurgische Simulatoren interessant sein könnten. Dabei kann experimentell untersucht werden, ob die erweiterte Unreal Integration für Serious Games im medizinischen Bereich geeignet ist.

\section{Soft Body Interaktion mit simulierten Kräften}
\label{sec_FSystem}
%Idee und Umsetzung von “VR-Physics2.0”, wie ich es genannt habe

Im Rahmen dieser Arbeit wurde eine Demoversion eines Serious Games implementiert, das in VR gespielt wird. Hierfür wurde eine VR Anlage mit Motion Controllern verwendet (Oculus quest 2), damit der Spieler mit virtuellen Händen nach chirurgischen Werkzeugen greifen und diese benutzen kann.

Für gewöhnlich werden die virtuellen Hände des Spielers stets auf die selbe Position und Rotation gesetzt, wie die des Motion Controllers. Dadurch befinden sich die virtuellen Hände immer genau da, wo sich die echten Hände des Spieler befinden, weil der Spieler die Motion Controller in den Händen hält.
Diese einfache Herangehensweise führt allerdings dazu, dass die virtuellen Hände des Spielers nicht mit statischen Objekten kollidieren können, weil statische Objekte nicht verschiebbar sind. Mit einem Flex-Objekt kann eine virtuelle Hand hingegen kollidieren, indem Sie Flex-Objekte eindrücken oder weg schiebt. Abbildung \ref{VHand1} zeigt wie dies in der virtuellen Welt aussieht. Wie man Sieht, kann der Spieler hier durch eine Mauer (weißer Block) hindurch greifen, wohingegen sich ein Flex-Objekt (grüne Kugel) deformieren und bewegen lässt. Eine Kollision zwischen virtueller Hand und der statischen Mauer ist nicht realisierbar, weil ein statisches Objekt nicht verschiebbar ist.

\bild{VHand1}{6cm}{Die grüne Kugel ist ein Flex-Objekt und lässt sich von eine virtuellen Hand berühren und beliebig stark eindrücken. Die Mauer (weißer Block) kann nicht berührt werden. Die virtuelle Hand durchdringt die Mauer ungestört, weil
%Sie immer auf die Position des Motion Controllers (orange) gesetzt wird und 
die statische Mauer nicht zur Seite gedrückt werden kann.}

Um die Immersion der VR-Anwendung zu erhöhen wäre es allerdings besser, wenn eine virtuelle Hand auch mit statischen Objekten kollidieren kann. Außerdem ist eine virtuelle Hand sozusagen unendlich stark: Soft Bodies können vom Spieler beliebig stark eingedrückt werden, egal wie steif oder schwer es ist. Das liegt daran, dass die virtuelle Hand immer auf die Position und Rotation des Motion Controllers gesetzt wird. Daher kann eine virtuelle Hand keine Physik simulieren und schiebt Objekte, die Phyik simulieren problemlos zur Seite.

Weil der Spieler unendlich viel Kraft hat und durch statische Objekte hindurch fassen kann, wirkt die VR-Anwendung weniger realistisch.

\subsubsection{Lösungsansatz: Soft Body Interaktion mit simulierten Kräften}
Um für mehr Realismus zu sorgen, wird ein kräftebasiertes Interaktions-System für die virtuellen Hände entwickelt. Nun ist die virtuelle Hand ein Flex-Objekt, das Physik simuliert. Statt die Hand direkt zur Position des Motion Controllers zu setzen, wird sie mithilfe von simulierten Kräften (und Moment) zum Motion Controller bewegt, beziehungsweise "gedrückt". Befindet sich ein Hindernis zwischen Hand und Motion Controller, kollidiert die Hand mit dem Hindernis, weil sie Physik simuliert. 

In Abbildung \ref{VHand2} ist dieses Prinzip zu sehen. Wie man sieht, kann sich der Motion Controller (orange gezeichnet) zwar in einem Gegenstand befinden, allerdings nicht die virtuelle Hand. Sie folgt nur dem Motion Controller, wenn es die Umgebung zulässt.

\bild{VHand2}{12cm}{Ein Moment und eine Kraft (roter Pfeil) "drücken" die virtuelle Hand zum Motion Controller. Links: Befindet sich der Motion Controller in einer statischen Mauer (blau), kollidiert die Hand mit der Mauer und drückt mit einer simulierten Kraft dagegen. Rechts: Weil die simulierte Kraft begrenzt groß ist, kann der Spieler den Weichkörper nur zu einem gewissen Grad eindrücken.}

%Werkzeuge
Ersetzt man die virtuelle Hand durch ein anderes Flex-Objekt, wirkt es für den Spieler, als habe er ein Werkzeug, beispielsweise ein Skalpell, in seiner Hand (siehe Abbildung \ref{Scalpel2}). 
%Ersetzen wir beispielsweise die virtuelle Hand durch ein Skalpell, wirkt es für den Spieler so, als habe er ein Skalpell in seiner Hand. 
Dank des kräftebasierten Interaktions-Systems kollidiert das Skalpell mit der Umgebung und der Spieler kann mit wenig oder viel Kraft schneiden, indem er seinen Motion Controller unterschiedlich bewegt und positioniert. Je größer der Unterschied zwischen Soll- und Ist-Position, desto stärker drückt das Skalpell auf das Gewebe und desto stärker wird es eingedrückt. Ist die Interaktionskraft stark und das modellierte Gewebe weich genug, drückt die Spitze des Skalpells die elastischen Gewebe-Partikel zur Seite und dringt in das Gewebe ein (siehe Abbildung \ref{Scalpel2}).

\bild{Scalpel2}{10cm}{ Die virtuelle Hand kann durch beliebige Werkzeuge ersetzt werden, beispielsweise durch ein Skalpell. Dünne Flex-Objekte können bei genügend Druckkraft in Soft Bodies eindringen. }

%im Ausblick nochmal drauf eingehen
Über die Verformung des Gewebes oder über die Größe der Druckkraft können entsprechende Ereignisse ausgelöst werden, beispielsweise das Schneiden des Gewebes oder haptisches Feedback.

Mit diesem Framework können beliebige Werkzeuge implementiert werden, die der Spieler in seine Hände nehmen, benutzen und wieder fallenlassen kann. Greift der Spieler mit einer virtuellen Hand nach einem Werkzeug, wird das Flex-Objekt der virtuellen hand mit dem Flex-Objekt des Werkzeugs ersetzt. Dadurch entsteht der Eindruck, der Spieler habe das Werkzeug aufgehoben.

Im Rahmen dieser Arbeit wurde eine Pinzette und eine Nadel implementiert. Mithilfe der Pinzette können Fremdkörper aus Gewebe gezogen und nach Flex-Partikeln gegriffen werden. Eine Nadel ist ein besonders dünnen Flex-Objekt und kann daher gut in Gewebe eindringen und darin stecken bleiben (siehe Kapitel \ref{sec_needle}).

\subsubsection{Regelung des kräftebasierten Interaktions-Systems}
Das kräftebasierte Interaktions-System kann als Regelsystem betrachtet werden. Dieses Regelsystem drückt mit Kräften und Momenten ($U = (\vec{F} , \vec{M}$) auf die virtuelle Hand, um sie in eine gewünschte Soll-Konfiguration, $W$, zu versetzen. Die Soll-Konfiguration ist die Position und Rotation des Motion Controllers, $W = (\vec{P_s} , \vec{R_s})$. Der Regelfehler, $E$, ergibt sich aus der Differenz zwischen Soll-Konfiguration, $W$, und der derzeitigen Ist-Konfiguration, $X = (\vec{P_i} , \vec{R_i})$ der virtuellen Hand.

\begin{equation}
E = W - X
\label{form_error}
\end{equation}

Dieses System kann mit unterschiedlichen Reglern geregelt werden. Im Rahmen der Arbeit wurde ein PD-Regler genutzt, der proportional zum Positions- und Rotations-Fehler ($E$) reagiert und außerdem proportional zur derzeitigen Geschwindigkeit reagiert ($\Delta E = E_0 - E_1$). Mit diesem Regler konnten schnell passende P- und D-Parameter gefunden werden, so dass die Hand dem Motion Controller sehr schnell folgt und ohne zu schwingen zum Stillstand kommt.

!!! hier evt die parameter angeben die gut geeignet waren !!!

Eine wichtige Eigenschaft, die der Regler haben sollte, ist eine Kraft, $F_{max}$ die er maximalen ausüben darf. $F_{max}$ repräsentiert dann nämlich die maximale Kraft, die der Spieler mit einer Hand ausüben kann. Wenn der Spieler beispielsweise maximal 50 kg heben können soll, wird die maximale Kraft auf $F_{max} = 50kg \cdot g$ gesetzt.

!!! hier evt noch bild indem 50kg und 60 kg aufgehoben werden. 60kg bleibt am boden liegen usw... !!!

\subsubsection{Experimentelle Parametrisierung für Gewebe}

Hier evt nochmal tipps rein wie gewebe-weichheit experimentell eingestellt werden kann

\section{Nadel-Gewebe Interaktion}%%%%%%%%%%%%%%%%%%%%%%%%%%%%%%%%%%%%%%%%%%%%%%%
\label{sec_needle}

%Wie umgesetzt, welche Tricks angewandt? Constraints? Geeignete Simulationsparameter?

%In dieser Arbeit wurde auch untersucht ob mit Flex und UE4 Nadel-Gewebe Interaktionen realisiert werden können. %Grundsätzlich gelingt dies, allerdings kann es viel Zeit in Anspruch nehmen, die korrekten Simulationsparameter zu finden. 
In diesem Kapitel wird ein einfaches Nadel-Gewebe Szenario implementiert und experimentell untersucht. Dabei wird das in Kapitel \ref{sec_FSystem} entwickelte Interaktionsystem genutzt um in VR Nadeln in ein Gewebe hinein zu stecken. 

Die Simulation wird dabei möglichst trivial gehalten um die Echtzeitfähigkeit zu gewährleisten. Das simulierte Gewebe kann nicht von der Nadel verletzt, beziehungsweise geschnitten werden. Daher muss die Partikelstruktur des Gewebemodels nicht verändert werden, so wie es bei virtuellen Schnitten getan werden muss (BEISPIEL PAPER REFERENZIEREN).

Das Flex-Asset für das simulierte Gewebe wird als Soft Body und die Nadel als Rigid Body modelliert, wobei die Verformbarkeit des Gewebes mithilfe von mehreren \textit{Shape Matching Constraints} realisiert wird, so wie in \cite{UPP}, \cite{BreastBiopsy} und \cite{PBDKidney}. 

Als Ausgangspunkt dienen die Werke \cite{PBDKidney} und \cite{BreastBiopsy} in denen bereits gute Simulationsparameter für unterschiedliche Gewebearten ermittelt wurden.

\subsubsection{Simulationsparameter und Einstellungen: }    

In \cite{PBDKidney} und \cite{BreastBiopsy} wurden bereits sehr gute Simulationseinstellungen gefunden um unterschiedliche Gewebearten realitätsnahe zu simulieren. Die in diesen Arbeiten empfohlenen Simulationseinstellungen wurden daher zum Großteil übernommen.

\begin{itemize}

  \item \textit{Sparse Signed Distance Field Collision} sollte deaktiviert sein. \\ 
  Diese Technik aus wird von NVIDIA Flex eingesetzt um zu vermeiden dass sich zwei Flex-Objekte ineinander verkeilen (siehe \cite{UPP} Kapitel 5.1). Befindet sich ein fremdes Partikel in einem Flex-Objekt, wird es mithilfe eines \ac{SDF} aus dem Flex-Objekt herausbewegt. In unserem Fall ist es allerdings erwünscht, dass Nadel-Partikel von Gewebe-Partikel festgehalten werden. Durch Anpassung der Steifheit, Partikeldichte und Partikelmasse des Gewebes kann eingestellt werden wie fest die Nadel-Partikel eingeklemmt werden.    
    
     \item Partikelradius, $r$: \\ Der Radius eines Partikels bestimmt seinen Interaktionsradius. Er wurde auf 0,5 bis 2 cm gestellt. Obwohl die Nadel eigentlich deutlich dünner als die Partikel war, konnten dennoch realistisch wirkende Ergebnisse erzielt werden. Kleinere Partikel wurden aus Performance Gründen nicht verwendet. Die benötigte Partikelmenge steigt nämlich kubisch mit abnehmendem Partikelradius.
     
      \item Reibungskoeffizient: \\ Der Parameter \textit{ParticleFriction} bestimmt wie stark die Reibung zwischen Partikeln sein soll. Mit ihm kann gut angepasst werden wie gut sich die Nadel im Gewebe bewegen lässt und wie stark sich das Gewebe heraus- oder herein-wölbt, wenn die Nadel bewegt wird.
      
      %in \cite{PBDKidney} und \cite{BreastBiopsy} wird substep 3 und iterations 9 genommen
      \item Simulation Substeps:  3%, so wie in \cite{PBDKidney} und \cite{BreastBiopsy} empfohlen.
      
      \item Substep Iterations: 9%, so wie in \cite{PBDKidney} und \cite{BreastBiopsy} empfohlen
      
      \item Gravitation: Wurde auf $(0 , 0 , -980) cm/s^2$ gesetzt, entsprechend der Erde.
      
      \item Partikel Dämpfung: \\ Mithilfe des Simulationsparameters \textit{Damping} können alle Partikelbewegungen gedämpft werden, als währen diese in einer Flüssigkeit. Es erwies sich als hilfreich etwas Dämpfung (0,1 bis 1) zu verwenden um das Auftreten von zitternden Bewegungen zu verringern.
      
\end{itemize}
    

\subsubsection{Modellierung der Nadel}
Für die Nadel-Gewebe Interaktion wird als erstes eine stark vereinfachte Nadel als Flex-Objekt modelliert, so wie in Abbildung \ref{Needle3} zu sehen. Hierbei wurde darauf geachtet, dass sich die Partikel der Nadel berühren und das Partikelmodel möglichst dünn, beziehungsweise Nadelförmig ist.

\bild{Needle3}{9cm}{Oben: Das Partikel-Modell der Nadel. Unten: Die hier grün gezeichnete Nadel wird mit dem kräftebasierten Interaktions-System in ein Gewbe-Modell eingeführt.}

Eine Besonderheit ist, dass der VR-Spieler die Nadel loslassen kann, nachdem sie in das Gewebe eingeführt wurde. In günstigen Situationen bleibt die Nadel dann im Gewebe stecken, weil Nadel-Partikel von Gewebe-Partikeln umgeben sind und sich darin verklemmen. 
Ob die Nadel im Gewebe stecken bleibt hängt vor allem mit der Gravitationskraft, die auf die Nadel wirkt, ab.
Ist die Gravitationskraft auf die Nadel-Partikel zu groß, fällt die Nadel nahezu ungehindert durch das Gewebe hindurch und schiebt dabei alle Gewebe-Partikel zur Seite. Daher ist darauf zu achten, die Partikelmassen der Nadel entsprechend klein zu halten, damit eine kleine Gravitationskraft auf die Nadel wirkt. Allerdings hat es eine leichte Nadel deutlich schwerer im Gewebe einzudringen. Wenn die Nadel-Partikel deutlich leichter sind als die Gewebe-Partikel, lassen sich die Gewebe-Partikel kaum von den Nadel-Partikeln verschieben.

Um dieses Problem zu lösen wird das Gewicht der Nadel je nach Situation angepasst: Während der Spieler die Nadel in der Hand hält, werden die Massen der Nadel-Partikel erhöht, damit Sie besser in das Gewebe eindringen können. Lässt der Spieler die Nadel los, wird die Masse der Nadel auf einen geringen Wert gesetzt, damit die Gravitationskraft auf die Nadel so klein ist, dass Sie nicht durch das Gewebe hindurch fällt. 

Alternativ kann die Nadel auch an das Gewebe befestigt werden (wie in Kapitel \textbf{XXX} beschrieben), wenn sie sich im Gewebe befindet und losgelassen wird.

%\subsubsection{Parametrisierung für Nadel und Gewebe}


%Um das soeben beschriebene Verhalten zwischen Nadel und Gewebe zu erhalten, wurde bei der Modellierung und Parametrisierung auf die folgenden Punkte geachtet: 

%Um das zuvor beschriebene Verhalten zu erreichen müssen bei der Parametrisierung und Modellierung die folgenden Punkte beachtet werden:

%\begin{itemize}%%%%%%%%%%%%%%%%%%%%%%%%%%%%%%

\subsubsection{Modellierung des Gewebes:}
Das Partikelmodel des Gewebes ist in Abbildung \ref{Tissue1}, rechts zu sehen. Die folgenden Punkte erwiesen sich bei der Modellierung als wichtig:

    \begin{itemize}
    
      \item Partikelabstand, $d$: \\ Die Gewebe-Partikel sollten sich nicht berühren ($d>2r$) sonst können instabile Zustände entstehen, bei denen einzelne Partikel keine Ruheposition (=lokales Minimum im Kräftefeld) mehr finden und dauerhaft umher springen und zittern. Gleichzeitig sollten die Gewebe-Partikel nah genug beisammen sein, damit Nadel-Partikel nicht durch Lücken fallen können, also $d<4r$. 
      %Ein Partikelabstand zwischen 2r und 3r erwies sich dabei als optimal.
      
      \item Irreguläre Anordnung der Partikel: \\ Die Partikel werden chaotisch/irregulär verteilt, wobei die Partikeldichte weitestgehend homogen bleiben sollte. Werden die Partikel beispielsweise als reguläres Gitter angeordnet (Abbildung \ref{Tissue1} links), neigt die Nadel dazu sich Entlang der Hauptachsen des regulären Gitters auszurichten. Außerdem ist dann der Reibungswiederstand entlang der Hauptachsen geringer als in diagonale Richtungen. Solch ein Verhalten ist unrealistisch und ist nicht erwünscht.
      
      \bild{Tissue1}{8cm}{Links: Partikelanordnung als reguläres 3D Gitter. Rechts: Chaotische Verteilung der Partikel, wobei die Partikeldichte weitestgehend homogen verteilt ist.}
      
      \item Keine Eigenkollision im Gewebe: \\ Die Partikel des Gewebes sollten nicht miteinander kollidieren können (Attribut \textit{SelfCollide} auf \textit{false} setzen). Wenn die Gewebe-Partikel miteinander kollidieren können, entstanden häufiger instabile Zustände bei denen Partikel keine Ruheposition fanden.
      
    \end{itemize}
    
%\end{itemize}%%%%%%%%%%%%%%%%%%%%%



\subsubsection{Erzieltes Verhalten}
%Mithilfe des gegebenen Frameworks konnte nach längerem Experimentieren das folgende Verhalten erreicht werden:

Dank des  kräftebasierten Interaktions-Systems (aus Kapitel \ref{sec_FSystem}) kann die Nadel mit dem Gewebe kollidieren und drückt dann mit einer simulierten Kraft dagegen. 
%Durch gezielte Positionierung des Motion Controllers kann die Stärke und Richtung der ausgeübten Kraft gezielt justiert werden. 
%Weil das Gewebe verformbar ist, kann die starre Nadel in das Gewebe eingeführt werden, wenn die aufgewendete virtuelle Kraft groß genug ist. Dann drücken die Nadel-Partikel die Gewebe-Partikel zur Seite. 
Ist die Interaktionskraft stark genug, drücken die Nadel-Partikel Gewebe-Partikel zur Seite und die Nadel dringt in das Gewebe ein.
Aufgrund der dünnen Form kann die Nadel nur mit axialen Bewegungen in das Gewebe eingeführt werden. Wird sie in einem ungünstigen Winkel eingeführt, kann sie seitlich wegrutschen, ohne in das Gewebe einzudringen (siehe Abbildung \ref{Needle4}). Ein ungünstiger Winkel ist ein besonders flacher Winkel zur Gewebeoberfläche. Dann ist die Richtung der simulierten Interaktionskraft ungünstig und der Kraftanteil, der in das Gewebe hineingeht, zu gering.

!!! Zu diskutieren: Ist dieses weg rutschen überhaupt realistisch??? Eine Nadel kann sehr flach angesetzt werden!

%Reibung
Befindet sich die Nadel im Gewebe, sorgt die simulierte Partikel-Reibung und Partikel-Kollision dafür, dass der Spieler noch mehr Kraft aufwenden muss, um die Nadel zu bewegen. Die Gewebe-Partikel, die in Kontakt mit der Nadel sind, haften an der Nadel und bewegen sich mit ihr mit. Dadurch verformt sich das Gewebe wenn die Nadel bewegt wird. Wird die Nadel herausgezogen, wölbt sich das Gewebe nach außen, wird sie hineingeschoben, wölbt sich das Gewebe nach innen (siehe Abbildung \ref{Needle4}).

\bild{Needle4}{10cm}{Oben: Verformung des Gewebes bei axialer Bewegung der Nadel. Unten links: Abrutschen der Nadel bei ungünstigem Stechwinkel. Unten rechts: Mehrere Nadeln wurden in unterschiedlichen Tiefen und Winkeln in das Gewebe gesteckt.}

%Steckenbleiben der Nadel:
Wenn die Nadel tief genug im Gewebe steckt und losgelassen wird, bleibt sie im Gewebe stecken (siehe Abbildung \ref{Needle4}). Wobei die Nadel in einigen Situationen kurz nach dem Loslassen eine kurze Bewegung vollzieht um darauf hin stehen zu bleiben. Das liegt daran, dass die Nadel-Partikel im übertragendem Sinne von mehreren Kraftfeldern umgeben sind, wegen der Partikel-Partikel-Kollision zwischen Gewebe- und Nadel-Partikeln. Wird die Nadel losgelassen, werden die Nadel-Partikel einen kurzen Moment von den Kraftfeldern bewegt bis sie ein lokales Minimum finden und dort zur Ruhe kommen.

\section{Weitere Chirurgische Szenarien}

Erwähnen: Nicht nur Nadeln können sich in Gewebe verklemmen, auch andere Fremdkörper wie Projektile oder Spreisel.



%\section{ALT}
% ---------------------------- ALT -----------------------------------

%Das konzipierte Simulationsmodell wird im Rahmen dieser Arbeit mithilfe der Unreal Engine 4 prototypisch implementiert. Bei dieser Implementierung wird der Fokus darauf gelegt, den Prototypen möglichst umfangreich und flexibel zu gestalten. Die Software bietet besonders viele Einstellungsmöglichkeiten und so können einige unterschiedliche Methoden und Variationen miteinander verglichen werden.
%
%Der Prototyp simuliert den Schwarm  ausschließlich im zweidimensionalem Raum. Die Darstellung des Schwarms erfolgt daher mit einer einfachen, zweidimensionalen Textur. Die Simulation eines dreidimensionalen Schwarms konnte aus Zeitgründen nicht mehr implementiert werden, zumal die verwendete Unreal Engine 4.18 zu diesem Zeitpunkt noch keine dreidimensionalen Texturen unterstützte.
%
%
%\section{Softwareübersicht}
%\label{sec_SoftwareuebersichtUmsetzung}
%
%Der in Kapitel \ref{sec_ShaderUebersichttheorie} gezeigte Aufbau für die Simulations-Shaderpipeline ist lediglich ein exemplarisches Beispiel. Die Shaderpipeline kann auch anders strukturiert werden. Wichtig ist nur, dass dabei alle Arbeitsschritte, aus Kapitel \ref{chap_Theorie}, ausgeführt werden. 
%
%Im Rahmen dieser Bachelorarbeit wurden möglichst viele Arbeitsschritte in einem Shader zusammengefasst, um weniger sogenannte \glqq Draw-Calls\grqq{}, also \glqq Zeichnungs-Aufrufe\grqq{}, zu generieren. So nennt man den Prozess des Beschreibens einer Textur. Das Beschreiben einer Textur nimmt signifikant Zeit in Anspruch, unabhängig von der Komplexität des zeichnenden Shaders. Daher macht es Sinn, die Anzahl der Schreibprozesse in der Shader-Pipeline klein zu halten.
%
%Wie Abbildung \ref{UebersichtUmsetzung} zeigt, besteht die implementierte Shader-Pipeline aus zwei Haupt-Shadern. Es werden also nur zwei Draw-Calls pro Simulationsschritt generiert. Der Shader \textit{SBoidBrush} wird nur dann ausgeführt, wenn der Nutzer der Software neue Individuen, mithilfe eines digitalen Pinsels, in $D$ zeichnet.
%
%\bild{UebersichtUmsetzung}{14cm}{Eine vereinfachte Softwareübersicht des entwickelten Prototyps. Im wesentlichen besteht dieser aus einer Shader-Pipeline, die von der Blueprint-Klasse \textit{FlockingSimBP} kontrolliert wird. Während die beiden Shader \textit{SGenerateControlSignals} und \textit{SBoidTransport} in gleichmäßigen Intervallen nacheinander ausgeführt werden, wird der  \textit{SBoudBrush}-Shader nur dann ausgeführt, wenn der Nutzer neue Boids erstellt.}
%
%Zuerst wird der Shader \glqq \textit{SGenerateControlSignals}\grqq{} ausgeführt. Dieser wendet alle eingesetzten Verhaltensregeln an und generiert direkt die finalen Steuervektoren, $\vec{V}$, durch gewichtete Summierung aller Einzelsteuersignale, $\vec{v}_i$. 
%
%\textit{SGenerateControlSignals} generiert die folgenden Einzelsteuersignale:
%	
%
%\begin{itemize}
%  \item Reynolds drei Grundverhaltensregeln:
%    \begin{itemize}
%      \item Kohäsion, $\vec{v}_c$
%      \item Separation, $\vec{v}_s$    
%      \item Ausrichtung, $\vec{v}_a$
%    \end{itemize}  
%  \item Zusatzverhaltensregeln:
%    \begin{itemize}
%      \item Punkt verfolgen, $\vec{v}_{cp}$
%      \item Punkt meiden, $\vec{v}_{sp}$      
%    \end{itemize}
%\end{itemize}
%
%Die Berechnungen der jeweiligen Einzelsteuersignale erfolgte genau so, wie in Kapitel \ref{chap_Theorie} beschrieben.
%
%All diese Einzelsteuersignale werden dann gewichtet und zu $\vec{V}$ summiert. Danach löscht \textit{SGenerateControlSignals} alle Steuervektoren, die zu einer leeren Zelle gehören. Dieser nachträgliche Löschvorgang wurde in Kapitel \ref{subsec_leereZellen} empfohlen um Fehlinterpretationen zu vermeiden. 
%
%Nach dem Löschvorgang ist die Steuervektoren-Textur bereit für die Übergabe in den zweiten Shader der Shader-Pipeline:  \glqq \textit{SBoidTransport}\grqq{}.
%
%Dieser Shader übernimmt den Transport der Boids. Hierfür stehen alle Verteilungsmethoden bereit, die in Kapitel \ref{section_BewDerBoids} vorgeschlagenen wurden.
%
%Die Parameterübergabe an die beiden Shader-Programme erfolgt von der CPU aus. Eine Klasse namens  \glqq \textit{FlockingSimBP}\grqq{} übernimmt diese Aufgabe. Diese Klasse erstellt und initialisiert alle Shader und Texturen. Danach steuert sie, wann welcher Shader mit welchen Parametern ausgeführt werden soll. Sie regelt auch, wie oft die Shader-Pipeline pro Sekunde abgearbeitet werden soll. 
%
%\section{Interaktivität}
%
%Die Klasse \textit{FlockingsSimBP} stellt auch die Schnittstelle zum Menschen dar und lässt einen Nutzer mit dem virtuellen Schwarm interagieren. Jegliche Interaktion wird ermöglicht, indem \textit{FlockingsSimBP} zur Simulationslaufzeit bestimmte Parameter der Shader-Pipeline ändert. 
%
%Der Nutzer kann mit der Maus entweder einen Kohäsions- oder einen Separations-Punkt in die Schwarm-Textur zeichnen. Diese Punkte wurden in Kapitel \ref{subsec_Zusatzvorschriften} eingeführt. Der Schwarm kann so in Echtzeit gejagt oder angelockt werden. Wenn gerade keiner der beiden Punkte eingesetzt wird, setzt \textit{FlockingsSimBP} die Gewichtung der Kohäsions- und Separations-Punkte auf 0. Wird einer der beiden Punkte vom Nutzer verwendet, wird dieser mit einer Gewichtung, die größer als 0 ist, versehen.  Wenn der Nutzer einen aktiven Punkt bewegt, übergibt \textit{FlockingsSimBP} die neuen Koordinaten an \textit{SGenerateControlSignals} weiter.
%
%Mithilfe des Kohäsionspunktes können zwei Boid-Gruppen zusammengeführt werden. Zudem kann mit diesem anziehenden Punkt vermieden werden, dass eine Gruppe den Simulationsbereich verlässt.
%
%Mit dem Separationspunkt können Gruppen gejagt werden. Bewegt der Nutzer einen Separationspunkt langsam durch einen Schwarm hindurch, kann eine Schwarm-Teilung erzwungen werden.
%
%\bild{Schwarmteilung}{8cm}{Eine vom Nutzer erzwungene Schwarmteilung. Die rötlich leuchtende Markierung stellt den eingesetzten Separationspunkt dar. Die beiden neu entstandenen Gruppen nehmen recht zügig abgerundete Formen an.}
%
%Abbildung \ref{Schwarmteilung} zeigt eine solche Schwarmteilung.
%
%Außerdem kann der Nutzer mit der Maus neue Boids zum Schwarm hinzufügen, indem er direkt in die Schwarm-Textur, $D$, zeichnet.
%
%\bild{Pinsel}{12cm}{Egal ob ungleichmäßig, gleichmäßig oder punktuell verteilt. Mithilfe eines implementierten Pinsel-Tools können unterschiedlichste Schwärme gezeichnet und editiert werden. Auch während der Simulationslaufzeit.}
%
%Hierfür wurden unterschiedliche \textit{Pinsel}-Shader angelegt, damit der Schwarm mithilfe unterschiedlicher Zeichenwerkzeuge  \glqq gezeichnet\grqq{} werden kann. In Abbildung \ref{Pinsel} sind Beispiele der zur Verfügung gestellten Pinsel-Muster zu sehen.
%
%\section{Darstellung des Schwarms}
%
%Bei der Darstellung des simulierten Schwarms ging es primär darum, interessante Zustände und Zwischengrößen erkennen zu können. Auf optisch aufwändigere Render-Techniken wurde vorerst verzichtet. 
%
%Der Nutzer hat die Möglichkeit die Schwarm-Textur, $D$ und die derzeitigen Steuervektoren, $\vec{V}$, zu beobachten. 
%
%\bild{Schwarmdarstellung}{10cm}{Links ist eine leicht nachbearbeitete Version der Boid-Textur, $D$, zu sehen. Rechts davon lässt sich beobachten in welcher Richtungen die Steuervektoren der Zellen zeigen.}
%
%In Abbildung \ref{Schwarmdarstellung} ist ein Screenshot davon zu sehen. Die Steuersignal-Textur wird unverändert angezeigt. Das Skalarfeld, $D$, wird hingegen leicht  \glqq aufgehübscht\grqq{}. Der Schwarm wird mit einer türkis-leuchtenden Grundfarbe gezeichnet. Dabei gilt: Je heller ein Pixel, desto mehr Boids enthält es. Wenn sich besonders viele Boids auf engem Raum aufhalten, wird dies durch hellere Farben erkenntlich. Dabei kann die Grundfarbe so stark aufgehellt werden, dass manche Pixel weiß strahlen. Weiße Pixel stellen also eine besonders volle Zelle dar. In Abbildung \ref{Schwarmdarstellung} können im Schwarm unterschiedliche Helligkeitsstufen erkannt werden. Solche Bilder entstehen meistens nur dann, wenn der Schwarm in Bewegung ist. Bewegt sich eine Gruppe Boids nicht, oder wenig, erscheinen alle Boids in etwa im selben Farbton. Eine stehende Gruppe befindet sich in einem stabilen Zustand, was voraussetzt, dass benachbarte Zellen ähnliche Zustände und somit auch den selben Farbton besitzen.
%
%Wenn mit starken Kohäsions-Gewichten gearbeitet wird, verdichten sich die Schwärme stärker, wodurch immer stärker leuchtende Gruppen entstehen. 
%
%%\clearpage
%\bild{KomprimierteGruppe}{4cm}{Wenn mit starker Kohäsion gearbeitet wird, nehmen die Zellen besonders viele Boids auf. Dadurch steigt die Boid-Dichte, die Schwärme werden kleiner und heller.}
%
%
%
%
%\section{Verhaltensänderungen}
%
%Die in Kapitel \ref{subsec_Verhaltensaenderungen} vorgeschlagenen Verhaltensänderungen wurden auch im Prototypen implementiert. So kann jederzeit per Knopfdruck  eine signifikante Verhaltensänderung im Schwarm ausgelöst werden, indem gezielt Gewichtungen verändert werden. Wie erhofft kann der Schwarm mit dieser einfachen Technik noch lebendiger und eigenwilliger wirken. Eine Änderung des Verhaltens erinnert auch an eine Änderung der Gemütszustände der Individuen. Vor allem die beiden Gegensätze  \glqq Aufregung\grqq{} und  \glqq Ruhe\grqq{} können gut erzeugt werden, indem das in Kapitel \ref{subsec_Verhaltensaenderungen} erwähnte \textit{Konkurrenzverhalten} gesteigert oder abgeschwächt wird. 
%
%Um Verhaltensänderungen für den Betrachter noch auffälliger zu gestalten, wird mit dem Gemütszustand auch die Farbe des Schwarms geändert. Der Schwarm ändert langsam seine Farbe von türkis nach rot, wenn sein Gemütszustand von ruhig nach aufgeregt wechselt.
%
%
%
%\section{Variierter Einsatz von Wahrnehmungsvektoren}
%
%In Kapitel \ref{subsec_nachbarsucheMitTestvektoren} wurden die sogenannten Wahrnehmungsvektoren, $\vec{s}$, eingeführt. Je nachdem wie viele und in welchen Konstellationen Wahrnehmungsvektoren eingesetzt werden, wirkt sich das unterschiedlich auf die Simulation aus.
%
%Im Rahmen dieser Arbeit wurde mit einigen unterschiedlichen Konstellationen experimentiert, um herauszufinden, wann und wie die Wahrnehmungsvektoren am effizientesten funktionieren. Es folgen nun einige Erfahrungsberichte über diese Thematik. Jegliche Beurteilungen über Simulationsqualitäten sind an dieser Stelle allerdings weitgehend subjektiv.
%
%%\subsection{gesonderte Separations-Vektoren}
%%Es könnte sinnvoll sein, kürzere Wahrnehmungsvektoren für die Separation einzusetzen, weil sich ein Boid bei diesem Gesetz mehr für die zu nah befindlichen Nachbarn  \glqq interessiert\grqq{}. Denn bei der Separation geht es darum, zu nah kommenden Nachbarn auszuweichen und so Kollisionen zu vermeiden. Daher könnten gesondert Wahrnehmungsvektoren für die Separation eingesetzt werden. Separations-Wahrnehmungsvektoren sind kürzer als die Wahrnehmungsvektoren für die restlichen Verhaltensregeln, weil für die Separation im Grunde nur die nahen Nachbarn interessant sich. 
%%
%%\bild{SeparationsWahrnehmungsvektoren}{6cm}{slkfdj sdf j}
%%
%%In Abbildung \ref{SeparationsWahrnehmungsvektoren} wird verdeutlicht, was gemeint ist. Auf der linken Seite ist ein Boid zu sehen, das zusätzliche Wahrnehmungsvektoren (, in Form von grünen Pfeilen,) besitzt. Sie sind ausschließlich für die Generierung von Separations-Steuersignalen, $\vec{v}_s$, gedacht. Rechts ist ein Boid dargestellt, welches für alle Verhaltensregeln die selben Wahrnehmungsvektoren einsetzt. 
%%
%%Aus Abbildung \ref{SeparationsWahrnehmungsvektoren} geht hervor, dass für diese Methode insgesamt mehr Wahrnehmungsvektoren eingesetzt werden müssen. Dann werden auch mehr Stichproben gezogen, wodurch ein insgesamt größerer Arbeitsaufwand entsteht. Dennoch wurde mit gesonderten Separations-Vektoren experimentiert, in der Hoffnung das sich der Mehraufwand mit einer besseren Simulationsqualität bezahlt macht.
%%
%%Im Test hat sich allerdings gezeigt, dass der Einsatz von gesonderten Separations-Vektoren keinen Mehrwert mit sich bringt. Werden für alle drei Grundverhaltensregeln die selben Wahrnehmungsvektoren eingesetzt, können gleichwertige Simulationsergebnisse erzielt werden, wie bei der Verwendung von gesonderten Separations-Vektoren. Hierfür müssen nur unterschiedliche Gewichtungen eingesetzt werden.
%
%%Eine interessante Erkenntnis ist hierbei die Tatsache, dass in diesem Simulationsmodell keine Separation entsteht, weil benachbarte Boids einen zu geringen Abstand zueinander haben. In diesem System separieren sich Boids, wenn in einer benachbarten Zelle zu viele Boids enthalten sind. Hier ist also weniger die Entfernung, sondern vielmehr die Menge ausschlaggebend. Nichts desto trotz können auch in diesem System besonders realistisch wirkende Sachwarmbewegungen entstehen.
%
%
%\subsection{Anzahl der eingesetzten Wahrnehmungsvektoren}
%
%Abbildung \ref{schwarmUnterschPara4VS8Vecs} zeigt einige Bilder aus Schwarmsimulationen mit unterschiedlich vielen Wahrnehmungsvektoren. Für die blau eingefärbten Schwärme wurden 4 Wahrnehmungsvektoren eingesetzt. Die rot gezeichneten Schwarmformationen entstanden hingegen beim Einsatz von 8 Wahrnehmungsvektoren. Auf den ersten Blick können nur schwer signifikante Unterschiede ausgemacht werden.
%
%\bild{schwarmUnterschPara4VS8Vecs}{10cm}{Die beiden oberen Schwärme verwendeten 4 Wahrnehmungsvektoren, während  die roten Schwärme mit 8 Wahrnehmungsvektoren ausgestattet waren. Es sind keine großen Unterschiede erkennbar.}
%
%Sowohl mit 4 als auch mit 8 Wahrnehmungsvektoren können abgerundete und individuelle  Schwarmformationen gebildet werden. Die Anzahl an eingesetzten Wahrnehmungsvektoren schien geringe Auswirkungen auf die Formbildung von Gruppen zu haben.
%
%Es hat sich schnell herausgestellt, dass auch mit wenigen Wahrnehmungsvektoren zufriedenstellende Bewegungen erzeugen werden können.
%
%
%
%
%
%
%
%\clearpage
%\subsection{Länge der Wahrnehmungsvektoren}
%\label{subsec_laengeDerWahrnehmungsvektoren}
%Abbildung \ref{SchwarmUnterschVLength} zeigt wie sich die Länge der eingesetzten Wahrnehmungsvektoren auf die Optik des simulierten Schwarms auswirkt.
%
%\bild{SchwarmUnterschVLength}{14cm}{Für diese Bilderreihe wurden unterschiedliche lange Wahrnehmunsvektoren eingesetzt. Von links nach rechts betragen die Vektorlängen 1, 2, 3, 5 und 10. Nicht nur die Körnung steigt an. Es entstehen auch ungleichmäßigere Schwarmformationen mit steigenden Vektorlängen.}
%
%Für diese Bilderreihe wurden immer 4 gleichmäßig verteilte Wahrnehmungsvektoren eingesetzt. Außerdem wurden für alle 5 Bilder die selben Gewichtungen verwendet. Von links nach rechts steigen die Längen der Wahrnehmungsvektoren.
%
%Aus Abbildung \ref{SchwarmUnterschVLength} geht deutlich hervor, dass die  \glqq Körnung\grqq{} von $D$ mit der Länge der Wahrnehmungsvektoren steigt. Mit Körnung ist gemeint, dass im Schwarm eindeutiger voneinander separierte Induviduen erkennbar sind. Es bilden sich sozusagen eindeutig erkennbare Untergruppen in einer Schwarm-Gruppe. Je länger die Wahrnehmungsvektoren sind, desto heller werden die Untergruppen, was bedeutet, dass dort die Boid-Dichten gestiegen sind.
%
%Die Längen der Wahrnehmungsvektoren scheint laut Abbildung \ref{SchwarmUnterschVLength} auch einen Einfluss auf die äußeren Formen des Schwarms zu haben. Sie wirken mit steigender Vektorlänge unförmiger und weniger abgerundet.
%
%Das Bewegungsverhalten der Schwärme verändert sich auch mit steigenden Vektorlängen. Die Boids können im übertragendem Sinne  \glqq weiter sehen\grqq{}, wenn sie mit längeren Wahrnehmungsvektoren ausgestattet werden. Das macht sich bemerkbar. Wenn die Individuen mit einer  \glqq weiteren Sicht\grqq{} ausgestattet werden, finden sie sich auch schneller zu Gruppen zusammen.
%
%Generell sollte ein Wahrnehmungsvektor immer länger als 1 sein. Ein kürzerer Vektor sorgt dafür, dass sich eine Zelle teilweise selbst ausliest.
%
%Je länger die Wahrnehmungsvektoren sind, desto stärker werden auch die generierten Steuersignale der Boids. Die Länge der Wahrnehmungsvektoren beeinflusst also auch das Bewegungsverhalten des Schwarms. Dem kann mithilfe von entsprechender Gewichtung entgegengewirkt werden.
%
%
%\subsection{Anordnung und Ausrichtung der Wahrnehmungsvektoren}
%Durch Experimente mit unterschiedlich angeordneten Wahrnehmungsvektoren ergaben sich die folgenden Erkenntnisse:
%
%\subsubsection{Asymmetrische Anordnung}
%Erwartungsgemäß entstehen durch asymmetrisch angeordnete Wahrnehmungssensoren auch asymmetrisch formierte Schwarmformationen. 
%
%Dies hat auch Auswirkungen auf die Bewegung des Schwarms. Der Schwarm bewegt sich schneller in die Richtungen, in die die Wahrnehmungsvektoren am häufigsten zeigen. Je weniger x-Anteile die Vektoren besitzen, desto seltener und langsamer bewegen sich die Individuen auch in diese Richtung. Existiert kein x-Anteil, finden auch keine Bewegungen in diese Richtung statt. 
%
%\subsubsection{Aymmetrische und gleichmäßige Anordnung}
%Um realistischere Bewegungen und Formationen zu erhalten, sollten die Wahrnehmungsvektoren symmetrisch und gleichmäßig um ein Boid herum angeordnet werden.
%
%Nur mit solchen Anordnungen bilden sich besonders abgerundete Gruppen. Außerdem bewegen sich die Schwärme dann in alle Richtungen gleich schnell, was vermutlich in den meisten Fällen erwünscht ist.
%
%\subsubsection{Bewegte Wahrnehmungsvektoren}
%
%In Kapitel \ref{subsec_nachbarsucheMitTestvektoren} wurde vorgeschlagen,  die Wahrnehmungsvektoren in Abhängigkeit von der Zeit zu bewegen. Eine Besonderheit ist hierbei, dass die Bewegung der Wahrnehmungsvektoren auch von der \ac{CPU} übernommen werden kann. Die Bewegung muss also nicht unbedingt die \ac{GPU} belasten, die ohnehin bereits von der Schwarmsimulation in Anspruch genommen wird.
%
%Die Vektoren dehnen und schrumpfen zu lassen, wirkt sich vor allem auf die Körnung der Schwarm-Textur aus. Dadurch entstehen nämlich die in Kapitel \ref{subsec_laengeDerWahrnehmungsvektoren} erwähnten Effekte. So ändert sich die Körnung der Schwarm-Textur mit der Zeit. Dieser Effekt kann bei Bedarf aber auch gezielt hervorgerufen werden, um den Schwarm optisch interessanter zu gestalten. Bei gezielter Parametrisierung können  pulsierende Effekte im Schwarm entstehen.
%
%Die Wahrnehmungsvektoren mit einer konstanten Rotationsgeschwindigkeit um einen Boid herum rotieren zu lassen, erwies sich im Test als lohnenswert. Dies wirkte sich vor allem positiv auf das Erscheinungsbild des Schwarms aus, wie Abbildung \ref{SchwarmRotierenNichtRotieren} zeigt.
%
%
%\bild{SchwarmRotierenNichtRotieren}{10cm}{Bei fixierten Wahrnehmungsvektoren entstehen ausgefranste Ausläufe an den Rändern der Schwärme (links). Dieser Effekt kann geglättet werden, indem rotierende Wahrnehmungsvektoren eingesetzt werden (rechts). }
%
%In \ref{SchwarmRotierenNichtRotieren} ist auf der rechten Seite ein Schwarm abgebildet, der mir rotierenden Wahrnehmungsvektoren arbeitet. Dieser wirkt geordneter und seine Ränder sind weniger  \glqq ausgefranst\grqq{} im Vergleich zum linken Schwarm. Der links abgebildete Schwarm arbeitet mit herkömmlichen, fixierten Wahrnehmungsvektoren. Mithilfe rotierender Vektoren konnten außerdem ruhigere und stabilere Schwärme erzeugt werden. Damit sind Schwärme gemeint, die sich ruhiger bewegen und rundere Formen bilden.
%
%
%
%\section{Einsatz der drei Grundverhaltensregeln}
%
%Mit den beiden Verhaltensregeln Kohäsion und Separation können bereits erste Erfolge erzielt werden. Mit diesen Regeln versetzen sich bereits erste Boids in Bewegung, um stabile Gruppen zu bilden. Sobald sich aber eine Gruppe zusammengetan hat, stabilisiert sich das System und die Gruppe bewegt sich nicht mehr. Solch ein stiller Schwarm wirkt nicht sehr \glqq lebendig\grqq{}. Deswegen hat die Verhaltensregel Ausrichtung einen hohen Stellenwert in der Simulation. Diese Verhaltensregeln kann nämlich bei optimaler Parametrisierung einem Schwarm mehr \glqq Leben\grqq{} einhauchen. Die Ausrichtung kann eine Eigendynamik in einer Gruppe auslösen, die zu selbstständigen Bewegungen führt. Solche Bewegungen können einen Schwarm mehr wie ein lebendiges und eigenwilliges Wesen aussehen lassen.
%
%In Kapitel \ref{subsec_Ausrichtung} wurden zwei unterschiedliche Formeln für die Berechnung des Steuersignals für die Ausrichtung, $v_a$, vorgeschlagen.
%
%Zuerst wurde eine exakte Berechnung angestrebt, wodurch Formel \ref{Form_Align} entstand. Danach wurde Formel \ref{Form_Align} vereinfacht, wodurch Formel \ref{Form_AlignCa} hergeleitet wurde. 
%
%Im Test zeigte sich, das mit beiden Rechenmethoden gute Simulationsergebnisse erzielt werden können. Es konnten keine signifikanten Unterschiede erkannt werden. 
%
%
%
%
%\section{Boid-Transporttechniken}
%\label{sec_LoseBoids}
%
%%-------------------messunge------------------------
%%fixe vektoren, r=2, 8 T, 4 wahrnehmungsvektoren: 	2997.4307	1662.9481   1661.5632
%%
%%rot vektoren, r=2, 8 T, 4 wahrnehmungsvektoren: 	2997.4307	2457.9767
%%
%%rot vektoren, r=1,2,  8 T , 8 wahrnehmungsvektoren: 	2997.4307	2702.9867
%%
%%rot vektoren, r=2,3,  8 T , 8 wahrnehmungsvektoren: 	2997.4307	2567.455
%%
%%fix vektoren, r=1,2,  8 T , 8 wahrnehmungsvektoren: 	2997.4307	2512.7107
%%
%%
%%fix vektoren, r=2, 4 T, 4 wahrnehmungsvektoren: 	2997.4307	382.9346
%%
%%rot vektoren, r=2, 4 T, 4 wahrnehmungsvektoren: 	2997.4307	1686.3750
%%
%%rot vektoren, r=1,2,  4 T , 8 wahrnehmungsvektoren: 	2997.4307	2426.6205 <-
%%
%%fix vektoren, r=1,2,  4 T , 8 wahrnehmungsvektoren: 	2997.4307	2603.3807
%%
%%
%%fix vektoren, r=2,  4 T smooth, 4 wahrnehmungsvektoren: 	2997.4307	.2362077
%%
%%rot vektoren, r=2,  4 T smooth, 4 wahrnehmungsvektoren: 	2997.4307	1.288783
%%
%%fix vektoren, r=2,  4 T smooth, 8 wahrnehmungsvektoren: 	2997.4307	47.120405
%%
%%rot vektoren, r=2,  4 T smooth, 8 wahrnehmungsvektoren: 	2997.4307	269.78772
%
%
%In Kapitel \ref{sec_Zellenbesuch} wurden zwei unterschiedliche Methoden vorgestellt, wie die Fortbewegung der Boids erfolgen kann. Dies wurde auch als Boid-Transport bezeichnet. Hierfür wurden die beiden Verteilungsgesetze  \glqq \textit{Verteilung in eine Zelle}\grqq{} und  \glqq \textit{Verteilung in zwei Zellen}\grqq{} definiert. 
%
%Die beiden Methoden werden nun miteinander verglichen. Hierbei ist nicht nur darauf zu achten, wie sich der Schwarm bewegt und verteilt, es wird auch kontrolliert, ob die Anzahl der im System befindlichen Boids konstant bleibt (Einhaltung des Gesetzes \ref{Form_BoidFlussKonst}).
%
%\subsection{Erhaltungsmessungen}
%
%Es wurden Erhaltungsmessungen vorgenommen, um einen möglichen Boid-Verlust nachzuweisen. 
%
%Hierfür wurde ein Shader erstellt, der alle Boids, die sich im Zellensystem befinden, zählt. Dieser Shader summiert alle Farbwerte von allen Pixeln aus der $D$-Textur. Für die Messungen war immer die selbe Anfangsbedingung gegeben, indem die Boid-Dichte-Textur wie folgt initialisiert wurde (siehe Abbildung \ref{InitBoids}):
%
%\bild{InitBoids}{6cm}{Die Boid-Dichte-Textur wurde mithilfe eines ungleichmäßigen Gradienten initialisiert. Alle Messungen wurden also mit der selben Menge Boids (2997,4307 an der Zahl) durchgeführt und die Menge war immer auf gleiche Weise verteilt.}
%
%Daher waren bei allen Messungen exakt die selben Anfangsbedingungen gegeben. Mithilfe eines Kohäsionspunkts, der sich genau im Zentrum des Systems befand, wurde dafür gesorgt, dass der Schwarm während einer Messung nicht aus der Textur heraus fliegt. 
%
%Die Messungen erfolgten nicht besonders detailliert, es ging nur darum, Boid-Verluste anhand von Stichproben nachzuweisen. Ein Messvorgang verlief folgendermaßen: Das Messprogramm speicherte während der Initialisierung die Anfangsmenge der Boids. Danach wurden 18000 Simulationsschritte durchgeführt. Dies entspricht einer Laufzeit von 5 Minuten, wenn 60 Simulationsschritte pro Sekunde abgearbeitet werden. Danach wurde gezählt, wie viele Boids übrig geblieben sind. Ein Messvorgang bestand also aus zwei Stichproben. Eine Anfangsmenge und eine Endmenge. Wie viele Boids in 5 Minuten verloren gegangen sind, kann durch Vergleich der Anfangsmenge und Endmenge in Erfahrung gebracht werden.
%
%In
%Tabelle \ref{Erhaltungsmessungen} 
%werden die Messergebnisse präsentiert:
%\clearpage
%
%% Please add the following required packages to your document preamble:
%% \usepackage{booktabs}
%\begin{table}[]
%\centering
%\caption{Einige Boid-Erhaltungsmessungen. Der Verlust von Boids konnte eindeutig nachgewiesen werden. Vor allem bei der Verteilung in zwei Zellen traten massive Verluste auf. }
%\label{Erhaltungsmessungen}
%\begin{tabular}{@{}llll@{}}
%\toprule
%\begin{tabular}[c]{@{}l@{}}Anzahl der\\ Wahrnehmungsvektoren\end{tabular} & Transportmethode                                                               & \begin{tabular}[c]{@{}l@{}}absoluter \\ Verlust\end{tabular} & \begin{tabular}[c]{@{}l@{}}prozentualer\\ Verlust\end{tabular} \\ \midrule
%4                                                                         & \begin{tabular}[c]{@{}l@{}}Verteilung in\\ eine Zelle,\\ $N_c=8$\end{tabular}  & 539,454                                                      & 17,99                                                          \\
%8                                                                         & \begin{tabular}[c]{@{}l@{}}Verteilung in\\ eine Zelle,\\ $N_c=8$\end{tabular}  & 294,444                                                      & 9,82                                                           \\ \midrule
%4                                                                         & \begin{tabular}[c]{@{}l@{}}Verteilung in\\ eine Zelle,\\ $N_c=4$\end{tabular}  & 1311,055 & 43,73                                                          \\
%8                                                                         & \begin{tabular}[c]{@{}l@{}}Verteilung in\\ eine Zelle,\\ $N_c=4$\end{tabular}  & 570,810 & 19,04                                                          \\ \midrule
%4                                                                         & \begin{tabular}[c]{@{}l@{}}Verteilung in\\ zwei Zellen,\\ $N_c=4$\end{tabular} & 2996,141 & 99,95                                                          \\
%8                                                                         & \begin{tabular}[c]{@{}l@{}}Verteilung in\\ zwei Zellen,\\ $N_c=4$\end{tabular} & 2727,642 & 90,99                                                         
%\end{tabular}
%\end{table}
%
%Aus den Messergebnissen aus Tabelle \ref{Erhaltungsmessungen} geht hervor, dass bei allen getesteten Transportmethoden Boids verloren gehen. Besonders schlimm ist dieses Problem bei der Verteilung in zwei Zellen. Hier war der Boid-Verlust auch mit bloßem Auge erkennbar.
%
%Ein Grund des Boid-Verlustes könnte die in Kapitel \ref{subsec_BilFiltUV} angesprochene bilineare Filterung sein, die sich in der Unreal Engine 4 nicht ohne weiteres deaktivieren lässt. 
%
%Wenn eine Zelle Boids von einer Nachbarzelle aufnimmt, geschieht dies indem sie ausliest wie viele Boids in der Nachbarzelle enthalten sind. Danach addiert sie den ausgelesenen Wert zu ihrem Boid-Bestand hinzu. Das Problem hierbei ist, dass der ausgelesene Wert ein bilinear gefilterter Wert ist. Es wird zwar versucht den exakten Pixelwert auszulesen, indem gezielt auf die Mitte eines Texels zugegriffen wird, dies ist allerdings, wie bereits erwähnt, nicht immer möglich. So könnte es vielleicht dazu kommen, dass Zellen 100\% von ihren Boids verlieren, aber andere Zellen nur 99\% davon aufnehmen. Kleine Verluste summieren sich dann über längere Zeit auf.
%
%Dies sind allerdings nur Vermutungen. Um das Problem genauer zu untersuchen blieb leider keine Zeit.
%
%
%\subsection{Simulationsqualitäten von unterschiedlichen Transportmethoden}
%
%Zwischen den beiden Transportmethoden \textit{Verteilung in eine Zelle} und  \textit{Verteilung in zwei Zellen} ist ein besonders deutlicher optischer Unterschied erkennbar, wie Abbildung \ref{SmoothAdvect} zeigt:
%
%\bild{SmoothAdvect}{10cm}{Links: Eine Gruppe Boids, die sich beim Transport in zwei Zellen verteilen. Dadurch entstehen sehr gleichmäßige Formen. Rechts: Eine Gruppe Boids, die sich in nur eine Zelle verteilen. Hier ist eine deutlich stärkere Körnung erkennbar.}
%
%Auf der linken Seit ist ein Schwarm zu sehen, der sich mittels \textit{Verteilung in zwei Zellen} fortbewegt. Auf der rechten Seite wird hingegen mit der \textit{Verteilung in eine Zelle} gearbeitet. Der rechte Schwarm ist deutlich körniger als der linke.
%
%Es hat sich außerdem gezeigt, dass sich der Einsatz von 8 möglichen Bewegungsrichtungen ($N_c=8$) lohnt. Der Rechenaufwind ist nicht sehr viel größer als bei $N_c=4$ und die Simulationsqualität steigt deutlich an. Wenn sich ein Individuum in nur 4 Richtungen bewegen kann, macht sich das deutlich in den Bewegungsmustern des Schwarms bemerkbar. Richtungswechsel wirken hart und abrupt.
%
%
%
%\clearpage
%\section{Instabilitäten}
%
%Erfahrungsgemäß ist die Simulation stabil, solange die folgenden Regeln befolgt werden:
%
%\begin{itemize}
%  \item Kapitel \ref{subsec_DefUVorschr}: Ein Boid darf beim Transport niemals eine Zelle überspringen.
%
%  \item Kapitel \ref{subsec_VerlasseneZellen}: Ein finales Steuersignal, $\vec{V}$, darf niemals größer als 1 sein.
%
%\end{itemize}
%
%Im Rahmen dieser Arbeit wurden keine Randbedingungen für das Simulationsmodell aufgestellt. Dies sollte nachgeholt werden, weil dadurch Probleme entstehen können. Im Test hat sich der simulierte Schwarm teilweise ins unendliche  \glqq aufgebläht\grqq{}, wenn er an den Rand der Textur geriet (siehe Abbildung \ref{Aufblaehen}). 
%
%\bild{Aufblaehen}{5cm}{Wenn eine Gruppe Boids dem Rand der Textur zu nahe kommt, kann sich der Schwarm aufblähen. Um dies zu vermeiden sollten Randbedingungen definiert werden.}
%
%
%
%
%
%
%\section{Performance-Messungen }
%\label{sec_performance}
%
%Die in dieser Arbeit genutzte Unreal Engine 4.18 beinhaltet einen umfangreichen Profiler \cite{GPUProfiling}. Mit Ihm konnte gemessen werden, wie lange die Abarbeitung der prototypischen Shader-Pipeline dauert. So entstanden einige Zeitmessungen, wobei stets mit unterschiedlichen Simulationseinstellungen experimentiert wurde. Der simulierte Schwarm wurde von unterschiedlich groß aufgelösten Texturen, $R$ mal $R$, beschrieben. Außerdem arbeiteten die Zellen mit unterschiedlich vielen Wahrnehmungsvektoren, $N_s$ und transportierten die Boids in unterschiedlich viele Richtungen, $N_c$.
%
%Die Messergebnisse sind in Tabelle \ref{gtx} aufgelistet. 
%
%
%\clearpage
%% Please add the following required packages to your document preamble:
%
%% \usepackage{booktabs}
%
%\begin{table}[]
%
%\centering
%
%\caption{Die prototypische Shader-Pipeline benötigt vor allem dann mehr Rechenzeit, wenn größere Texturen eingesetzt werden.}
%
%\label{gtx}
%
%\begin{tabular}{@{}lllllll@{}}
%
%\toprule
%
%                                                                                & \begin{tabular}[c]{@{}l@{}}$N_c=8$\\ $N_s=4$\\ $R=256$\end{tabular} & \begin{tabular}[c]{@{}l@{}}$N_c=8$\\ $N_s=4$\\ $R=1024$\end{tabular} & \begin{tabular}[c]{@{}l@{}}$N_c=8$\\ $N_s=4$\\ $R=4096$\end{tabular} & \begin{tabular}[c]{@{}l@{}}$N_c=8$\\ $N_s=8$\\ $R=256$\end{tabular} & \begin{tabular}[c]{@{}l@{}}$N_c=8$\\ $N_s=8$\\ $R=1024$\end{tabular} & \begin{tabular}[c]{@{}l@{}}$N_c=8$\\ $N_s=8$\\ $R=4096$\end{tabular} \\ \midrule
%
%\multicolumn{1}{l|}{Zeit-}                                                      & 0,39                                                                & 1,89                                                                 & 4,66                                                                 & 1,27                                                                & 1,55                                                                 & 4,6                                                                  \\
%
%\multicolumn{1}{l|}{messungen}                                                  & 0,44                                                                & 1,67                                                                 & 4,59                                                                 & 3,79                                                                & 5,87                                                                 & 5,45                                                                 \\
%
%\multicolumn{1}{l|}{in ms}                                                      & 0,42                                                                & 0,47                                                                 & 4,59                                                                 & 1,28                                                                & 1,57                                                                 & 4,78                                                                 \\
%
%\multicolumn{1}{l|}{}                                                           & 0,42                                                                & 2,11                                                                 & 5,16                                                                 & 1,37                                                                & 0,63                                                                 & 5,8                                                                  \\
%
%\multicolumn{1}{l|}{}                                                           & 0,39                                                                & 2,53                                                                 & 3,62                                                                 & 1,49                                                                & 2,12                                                                 & 4,63                                                                 \\
%
%\multicolumn{1}{l|}{}                                                           & 0,44                                                                & 2,01                                                                 & 5,07                                                                 & 1,80                                                                & 2,06                                                                 & 4,91                                                                 \\
%
%\multicolumn{1}{l|}{}                                                           & 0,44                                                                & 1,51                                                                 & 3,86                                                                 & 0,17                                                                & 2,13                                                                 & 4,57                                                                 \\
%
%\multicolumn{1}{l|}{}                                                           & 0,42                                                                & 1,51                                                                 & 5,14                                                                 & 1,22                                                                & 2,33                                                                 & 4,79                                                                 \\
%
%\multicolumn{1}{l|}{}                                                           & 0,42                                                                & 2,07                                                                 & 5,09                                                                 & 2,03                                                                & 2,28                                                                 & 4,85                                                                 \\
%
%\multicolumn{1}{l|}{}                                                           & 0,44                                                                & 2,19                                                                 & 3,64                                                                 & 0,39                                                                & 5,92                                                                 & 4,67                                                                 \\ \midrule
%
%\multicolumn{1}{l|}{\begin{tabular}[c]{@{}l@{}}Mittelwert\\ in ms\end{tabular}} & 0,422                                                               & 1,796                                                                & 4,542                                                                & 1,481                                                               & 2,646                                                                & 4,9050                                                               \\
%
%\multicolumn{1}{l|}{\begin{tabular}[c]{@{}l@{}}Streuung\\ in ms\end{tabular}}   & 0,0193                                                              & 0,5637                                                               & 0,6206                                                               & 0,9893                                                              & 1,7844                                                               & 0,4036                                                             
%
%\end{tabular}
%
%\end{table}
%
%
%Die in Tabelle \ref{gtx} zu sehenden Messergebnisse ergaben sich mit einem Computer, der mit einer NVIDIA GTX 970 Grafikkarte und einer intel Xeon E3-1230 v3 CPU ausgestattet war.
%
%Wie bereits erwähnt zeichnen sich Shader, die mit der Unreal Engine programmiert wurden, durch eine besonders gute Portabilität aus. Deswegen konnte der Prototyp ohne Umstände auf einem Laptop, der keine separate Grafikkarte besitzt, getestet werden. Es handelte sich dabei um ein \textit{ThinkPad X240}, das mit einer \textit{Intel Core i5-4210U}-CPU ausgestattet war.
%
%
%\clearpage
%Der Prototyp lief mit über 30 Bildern pro Sekunde auf diesem Testgerät. Die Messergebnisse werden in Tabelle \ref{ThinkPad} aufgelistet.
%
%
%\begin{table}[]
%\centering
%\caption{Die Verarbeitungsdauer der Shader-Pipeline auf einem etwas schwächerem Gerät.}
%\label{ThinkPad}
%\begin{tabular}{@{}ll@{}}
%\toprule
%                                & \begin{tabular}[c]{@{}l@{}}$N_c=8$\\ $N_s=4$\\ $R=256$\end{tabular} \\ \midrule
%                                & 1,89 ms                                                                \\
%                                & 1,43 ms                                                               \\
%                                & 1,30 ms                                                                \\
%                                & 1,41 ms                                                                                                                               \\
%                                & 1,85 ms                                                                                                                               \\
%                                & 1,86 ms                                                                                                                               \\
%                                & 1,76 ms                                                                                                                               \\
%                                & 1,32 ms                                                                                                                               \\
%                                & 1,91 ms                                                                                                                               \\
%                                & 1,32 ms                                                                                                                               \\ \midrule
%\multicolumn{1}{l|}{Mittelwert in ms} & 1,605                                                               \\
%\multicolumn{1}{l|}{Streuung in ms}   & 0,2682                                                             
%\end{tabular}
%\end{table}

% Die Arbeit besteht aus Kapiteln (chapter)
\chapter{Fazit und Ausblick}

Die für diese Arbeit gesetzten Ziele konnten größtenteils erreicht werden. 

...Zwar sind solche Simulatoren unbestreitbar weniger realistisch, dafür sinken die nötigen Kosten weil keine haptische Hardware benötigt wird. Zumal heutige VR-Geräte für immer mehr Privatpersonen erschwinglich sind. Das in dieser Arbeit verwendete VR-Gerät, die Oculus Quest 2, kostete beispielsweise 300 Euro...

% ----------------------------------- ALT --------------------------------------------

%%\section{Fazit}
%
%Im Rahmen dieser Arbeit konnten fast alle gesetzten Ziele erreicht werden. Anhand des konzipierten Simulationsmodells konnte ein Prototyp implementiert werden, der fast alle Anforderungen aus Kapitel \ref{sec_Anforderungen} erfüllt. 
%
%Der im Prototyp simulierte Schwarm kann durch Einsatz einfacher Zusatzverhaltensregeln von einem menschlichen Spieler gejagt oder verfolgt werden. Das bestehende System lässt sich zukünftig mit neuen Verhaltensregeln erweitern. Beispielsweise könnte ein Spieler ein Flussfeld zeichnen, um virtuelle \glqq Autobahnen\grqq{} für die Boids zu generieren. Gerät ein Individuum auf solch eine Bahn, folgt es dem Flussfeld, so wie Autos einer Autobahn folgen.
%
%Die Bewegungen des Schwarms können bei entsprechender Parametrisierung realistisch und lebendig gestaltet werden. Ohne mit zufällig generierten Steuervektoren arbeiten zu müssen, bewegen sich manche Schwärme scheinbar in zufällige Richtungen, als ob die virtuellen Lebewesen ein eigenes Ziel verfolgen würden.
%
%Die Performance des implementierten Prototyps erwies sich, entgegen aller Erwartungen, als überraschend gut. Die Messungen aus Kapitel \ref{sec_performance} lassen vermuten, dass sich das konzipierte Simulationsmodell durchaus in der Praxis anwenden lässt. Künftige Videospiele könnten von diesen interaktiven Schwarm-Wesen profitieren, zumal sich die entwickelte Simulation von herkömmlichen Shadern ausgeführt werden kann. Shader-Programme, die mithilfe der Unreal Engine 4 implementiert wurden, lassen sich ohne große Umstände auf einige Plattformen portieren. Von allen handelsüblichen Konsolen, bis hin zu Smartphones. Smartphone-Besitzer mithilfe des Touchscreens mit dem Schwarm interagieren zu lassen, ist durchaus umsetzbar. 
%
%Der erstellte Prototyp konnte aus zeitlichen Gründen nur als zweidimensionale Schwarmsimulation implementiert werden. Die Idee, den virtuellen Schwarm als dreidimensionale Leuchtpunkt-Wolke, mittels volumetrischer Render-Methodik, darzustellen, wird aber weiterhin mit großem Interesse verfolgt. Zu diesem Zeitpunkt ist jedoch noch nicht bekannt, wann die Unreal Engine 4 dreidimensionale Textur-Formate unterstützen wird.
%
%
%
%%\section{Ausblick}
%
%Trotz allem ist das bestehende System durchaus verbesserungswürdig. Besonders negativ aufgefallen ist der in Kapitel \ref{sec_LoseBoids} angesprochene Verlust von Boids. Diese Problematik konnte aus Zeitgründen nicht weiter untersucht werden, weswegen noch keine konkrete Erklärung für das Verschwinden von Individuen existiert.
%
%Außerdem könnte die Darstellung des Schwarms weiter verbessert werden. Die Schwarm-Textur wirkt zum Teil noch recht körnig und unregelmäßig. Es entstehen auch hin und wieder harte Kannten und Ecken.
%
%Eine weitere interessante Erweiterung wäre, die Wahrnehmungsvektoren der Boids wie folgt einzusetzen:  Wie in der Robotik üblich, könnten die Wahrnehmungssensoren nur an die Vorderseite eines Boids befestigt werden. Dann würde jedes Individuum nur den Bereich wahrnehmen können, der in dessen Bewegungsrichtung liegt. Das könnte den Realismus der Simulation steigern, weil dann die Bewegungsrichtung eines Boids auch seine \glqq Blickrichtung\grqq{} wäre.
%
%Denkt man über die Funktionsweise des konzipierten Modells genauer nach, lässt sich eine interessante Feststellung machen. In diesem Simulationsmodell entsteht keine Separation oder Kohäsion weil benachbarte Boids einen zu geringen oder zu großen Abstand zueinander haben. In diesem System handeln die Individuen dann wenn in einer benachbarten Zelle eine bestimmte Menge Boids enthalten ist. Hier ist also nicht die Entfernung, sondern vielmehr die Menge ausschlaggebend. Nichts desto trotz können auch in diesem System besonders realistisch wirkende Sachwarmbewegungen entstehen.
%
%Zum Schluss lässt sich sagen, dass sich das geniale Prinzip von Schwärmen erwiesen hat: Obwohl jede Zelle des Zellensystems nur wenige einfache Fähigkeiten beherrscht, entstand am Ende ein Gesamtsystem, das sich viel komplexer als seine einzelnen Bausteine verhalten kann. Jedes Boid kann nur wenige Bewegungsrichtungen einschlagen und besitzt sehr eingeschränkte Wahrnehmungsfähigkeiten. Aber dennoch kann sich eine große Gruppe aus Boids scheinbar in alle möglichen Richtungen bewegen und dabei einzigartige Formationen bilden.
%





%% Die Arbeit besteht aus Kapiteln (chapter)
\chapter{CheatSheet}

In diesem Kapitel wird das vorliegende System betrachtet und es werden Differentialgleichungen ermittelt, die das System mathematisch beschreiben. Besonderer Augenmerk wird hierbei auf die translatorische und rotatorische Bewegung des Systems gelegt. Gesucht sind zwei gekoppelte Differenzialgleichungen, die sowohl
die $X$-Position des Wagens, als auch die Drehung $\theta$ des Pendels beschreiben. In
mindestens einer der beiden Gleichungen muss die Stellkraft $U$ enthalten sein.
Zur Vereinfachung wird das Gesamtsystem in zwei Teilsysteme unterteilt. Durch
Freischneiden am Gelenk des Pendelstabes erhalten wir die Teilsysteme Wagen
und Pendel die nun einzeln Analysiert werden. 


\section{Modellierung des Wagens}
\subsection{Mathematische Modellierung des Wagens}
Zuerst wird das Teilsystem Wagen untersucht. Hierbei wird vor allem die translatorischen Bewegungen des Systems betrachten, um zu analysieren, wie sich der Wagen verhält. 

Auf ihn wirken die folgenden Kräfte:

\bild{TeilsystemWagen}{6cm}{Teilsystem Wagen}

Der Wagen kann sich nur in $X$-Richtung bewegen. Daher sind nur die Horizontalen Kräfte von Interesse. In positiver $X$-Richtung wirkt die Stellkraft $U$, in negativer Richtung wirkt die Kraft $F_x$. $F_x$ ist die Kraft, welche der Stab des Pendels auf den Wagen ausübt.

Zunächst wird Newton II auf das dynamische System angewandt:

\begin{equation}
M \cdot \ddot{x}_w =  U -F_x
\label{Fx0}
\end{equation}

$x_w$ ist die horizontale Position des Wagens. Aus Gleichung \ref{Fx0} ergibt sich für $\ddot{x}_w$:

\begin{equation}
\ddot{x}_w =  \frac{U - F_x}{M}
\label{x_wpp}
\end{equation}

Nun liegt eine Gleichung vor, in der $\ddot{x}_w$, beziehungsweise $x_w$ enthalten ist. Um die translatorische Bewegung des Systems vollständig beschreiben zu können, muss nun noch die Kraft $F_x$ hergeleitet werden. Da es sich hierbei um die Kraft handelt, mit welcher der Stab horizontal auf den Wagen einwirkt, wird nun das Teilsystem Pendel mathematisch untersucht, um einen Ausdruck für die Kraft $F_x$ zu finden:

\bild{TeilsystemPendel}{6cm}{Teilsystem Pendel}

Der Stab hat die Länge $l$ und sein Massenschwerpunkt liegt bei $l/2$. Auch in diesem Fall werden ausschließlich die horizontalen Kräfte berücksichtigt. Gemäß des dritten Newtonschen Axioms Actio gleich Reactio wirkt auf den Stab dieselbe Kraft $F_x$ wie auf den Wagen. Des Weiteren wirkt die horizontale Trägheitskraft $F_T$ auf das Pendel.

Laut Newton II gilt:

\begin{equation}
F_T = m \cdot \ddot{x} = F_x
\label{pendel0}
\end{equation}
%\clearpage
Wobei sich $x$ aus der Summe der Position des Wagens $x_w$ und der relativen Position
des Stabes $x_s$ zusammen setzt:

\bild{positionX}{6cm}{$x = x_w + w_s$}

Dieser Sachverhalt kommt daher, dass $F_T$  nicht nur von der Bewegung des Pendels ($x_s$), sondern auch von der Bewegung des Wagens ($x_w$) abhängt. $x$ beschreibt also die absolute Position des Schwerpunktes des Pendels.

Setzen wir nun Gleichung \ref{pendel0} in Gleichung \ref{x_wpp} ein, ergibt sich für $\ddot{x}_w$:

\begin{equation}
\ddot{x}_w =  \frac{U}{M} - \frac{m \cdot \ddot{x}}{M}=\frac{U}{M} - \frac{m(\ddot{x}_w+\ddot{x}_s)}{M}
\label{x_wpp2}
\end{equation}

Durch Ausklammern von $\ddot{x}_w$ und Umformen nach $\ddot{x}_w$ ergibt sich schlussendlich: 

\begin{equation}
\ddot{x}_w =  \frac{1}{M+m} \cdot (U-m\cdot \ddot{x}_s) 
\label{x_wpp3}
\end{equation}
Dieser Ausdruck lässt sich nun wie gewünscht für $F_x$ einsetzen, vorher muss jedoch noch $x_s$ mithilfe von $\theta$ ausgedrückt werden. Nach Betrachtung der gegebenen Geometrie ergibt sich:

\begin{equation}
x_s = f(\theta) = \frac{l}{2} \cdot \sin(\theta)
\label{x_s}
\end{equation}

Mit dem Wissen, dass $\theta$ von der Zeit abhängig ist, lässt sich auch $\ddot{x}_s$ bilden, indem $x_s=f(\theta)$ zwei mal nach der Zeit abgeleitet wird:

\begin{equation}
\dot{x}_s=\frac{l}{2} \cdot \dot{\theta} \cdot \cos(\theta)
\label{x_sp}
\end{equation}

\begin{equation}
\ddot{x}_s=\frac{l}{2} \cdot [\ddot{\theta}\cos(\theta) - \dot{\theta}^2\cdot \sin(\theta)]
\label{x_spp}
\end{equation}

Abschließend muss nur noch Gleichung \ref{x_spp} in Formel \ref{x_wpp3} eingesetzt werden:

\begin{equation}
\ddot{x}_w=\frac{1}{M+m} \cdot [U-\frac{1}{2}ml \cdot (\ddot{\theta}\cos(\theta)-\dot{\theta}^2\cdot\sin(\theta))]
\label{x_wpp_final}
\end{equation}

Damit ergibt sich nun eine Formel, welche die Bewegung des Wagens in $X$-Richtung beschreibt. Der Ausdruck ist von der Stellkraft $U$ und vom Winkel $\theta$ abhängig. Wie bereits erwähnt wird hierbei die Reibung des Wagens vernachlässigt.

Linearisiert man diese Gleichung um den Arbeitpunkt $\theta=0$ ergibt sich:

\begin{equation}
\ddot{x}_w \approx \frac{1}{M+m}[U-\frac{1}{2}ml\ddot{\theta}]
\label{x_wpp_final_lin}
\end{equation}



\clearpage
\subsection{Modellierung des Wagens mit Simulink}%..............Modellierung des Wagens mit Simulink.........................
Nachdem die Systemgleichung des Wagens hergeleitet wurde, wird mit Simulink ein entsprechendes Blockschaltbild für Gleichung \ref{x_wpp_final} entworfen:

\bild{Wagensimulink}{12cm}{Simulink Modell des Wagensystems}

Das Simulink-Modell für die linearisierte Gleichung (\ref{x_wpp_final_lin}) ist etwas kompakter:

\bild{WagensimulinkLin}{11cm}{Simulink Modell des linearisierten Wagensystems}


\begin{table}[h]
  \caption{Auswahl der Regler}
  \label{reglerTab}
  \renewcommand{\arraystretch}{1.2}
  \centering
  \sffamily
  \begin{footnotesize}
    \begin{tabular}{l l l}
    \toprule
    \textbf{Teilsystem} & \textbf{Übertragungsverhalten} & \textbf{Regler}\\
    \midrule
    Wagen	&	$I$	& 	$PI$\\
    Pendel	&	$P-T_2$		&	$PID$\\
    \bottomrule
    \end{tabular}
  \end{footnotesize}
  \rmfamily
\end{table} % das ist mein CheatSheet für Latex!!!!!!!!!!!!
% ------------------------------------------------------------------

\label{lastpage}

\clearpage

%% Backmatter mit normalem Zeilenabstand setzen
%\singlespacing
%
%% Römische Ziffern für die "Back-Matter", fortlaufend mit "Front-Matter"
%\pagenumbering{roman}
%\setcounter{page}{\value{frontmatterpage}}

% Literaturverzeichnis erzeugen (nach DIN)
\begin{flushleft}
%\bibliographystyle{dinat} % Literatur nach DIN
\bibliographystyle{apalike} % unsrt=Zitate mit [1], [2] etc.
%\bibliographystyle{alpha} % Zitate mit [Kor01], [Vix99] etc.
\bibliography{literatur}   % BibTeX-Datei mit Literaturquellen einbinden
\end{flushleft}


% Index ausgeben
%\cleardoublepage
%\phantomsection
%\addcontentsline{toc}{chapter}{Index}
%\printindex

\end{document}